\chapter{Strong topologies, capacities, and pluripolar sets}
\index{strong topology!BJ}\index{capacity!BJ}\index{pluripolar set}
\label{chap:bj-strong-topologies}

Weak convergence records values at divisorial valuations, or integrals
against continuous functions.  It does not by itself control loss of energy
into singularities.  The strong topologies add exactly that missing scalar
datum.  Capacity then provides the set-theoretic counterpart: psh functions
are continuous off sets of arbitrarily small capacity, and capacity zero is
equivalent to pluripolarity once the envelope property is available.

\section{Orientation}
\label{sec:ch48-orientation}

\subsection{What this chapter does}
\label{sec:ch48-what}

\Cref{chap:bj-energy} constructed the Monge--Amp\`ere energy \(E_\omega\)
and showed that it is upper semicontinuous, but not continuous, for the
weak topology; \Cref{ex:ch46-energy-not-weakly-continuous} exhibited a
sequence of the simplest possible PL psh functions on \(\mathbf P^1\) whose
weak limit is \(0\) and whose energies are all equal to \(-\tfrac12\).
\Cref{chap:bj-finite-energy} built the metric geometry of the class
\(\cE^1(\omega)\) out of the functionals \(I\), \(J_\psi\) and \(J\), and
solved the Monge--Amp\`ere equation \(\MA_\omega(\phi)=\mu\) for every
\(\mu\) of finite dual energy; but it deliberately deferred, in
\Cref{rem:ch47-strong-deferred}, three statements about the resulting
topology.  This chapter is where those statements are made and where the
deferred proofs are carried out or, when they are imported, isolated with
an exact reference.

The organizing idea is a single sentence.  \emph{The weak topology sees the
values of a potential and the mass of a measure; the strong topology sees,
in addition, one real number --- the energy --- and that number is exactly
what is needed to make the Monge--Amp\`ere operator a homeomorphism.}
Everything else in the chapter is either a consequence of this, or the
set-theoretic shadow of it: the capacity, which measures how much
Monge--Amp\`ere mass a set can carry, and the three notions of smallness
that the capacity organizes.

Four things are worth stating at the outset, because they are what
distinguishes this chapter from its counterpart
\Cref{chap:bfj-capacity,chap:bfj-energy,chap:bfj-finite-energy-measures}
over a discretely valued base.

First, the trivially valued theory has a distinguished point.  By
\Cref{lem:ch47-three-normalizations}, \(\MA_\omega(0)=\delta_{v_{\triv}}\)
and \(\sup_{X^{\An}}\phi=\phi(v_{\triv})\) for every \(\phi\in\PSH(\omega)\).
Consequently \(J(\phi)=\sup\phi-E_\omega(\phi)=J_0(\phi)\), the supremum
functional is an evaluation at a divisorial point, and hence \emph{weak
convergence already forces convergence of suprema}.  This is why the
supremum term in \eqref{eq:bj-strong-e1-I} is a genuine addition to \(I\)
and not a second convergence requirement.

Second, the envelope theorem of
\Crefrange{chap:bj-envelope-statement}{chap:bj-envelope-consequences} is
available.  It is what makes \(\PSH^{\sup}(\omega)\) compact
(\Cref{cor:ch45-compactness}), what makes \(P_\omega\) an operator
(\Cref{imp:ch45-obstacle-operator}), and what makes the continuous
regularization \Cref{thm:ch45-continuous-regularization} available in a
fixed ample class.  The last of these is the engine of the quasicontinuity
theorem below: over a discretely valued base one regularizes by
\emph{model} functions and the corresponding argument
(\Cref{thm:bounded-quasicontinuity}) has to work with a net of model
metrics; here one regularizes by \emph{continuous psh} functions in the
class itself, and the proof shortens accordingly.

Third, the capacity constructed here is normalized to be a probability:
\(\Capa_\omega(X^{\An})=1\).  This is the probability normalization
declared in \Cref{warn:ch46-normalization} and inherited without
amendment.  Over a discretely valued base \Cref{def:bfj-capacity} gives
\(\Capa_\theta(X^{\An})=V\); every numerical comparison between the two
chapters therefore carries a factor \(V\), and every such comparison below
displays it.

Fourth --- and this is the point at which the reader should be most
careful --- the equivalence of pluripolarity, negligibility and vanishing
outer capacity is \emph{not} proved here in its entirety.  Two of the six
implications are imported from \cite{BJ22} with exact references, and
\Cref{sec:ch48-what-is-open} says in prose exactly which ones, why the
arguments available here do not reach them, and what would have to be
supplied.  The corresponding statement over a discretely valued base is
strictly weaker: \Cref{cor:ch25-polar-negligible} gives one implication and
\Cref{rem:ch25-curve-converse} records that the converse is unavailable
there.  The gain in the trivially valued setting comes entirely from the
envelope property.

\subsection{The five tasks}
\label{sec:ch48-tasks}

\begin{enumerate}
\item\label{itm:ch48-task:potentials}
      \emph{The strong topology on \(\cE^1(\omega)\).}  Define it as the
      coarsest refinement of the weak topology for which \(E_\omega\) is
      continuous, identify it with the topology of the quasi-metric
      \(\mathbb I\) of \eqref{eq:bj-I-normalized}, and prove that it is
      complete.  \Cref{sec:ch48-potentials}.
\item\label{itm:ch48-task:measures}
      \emph{The strong topology on \(\cM^1(\omega)\).}  Define the Fenchel
      defect, the dual quasi-metric \(I^\vee\), and prove the elementary
      structure of both; state the three equivalent forms of strong
      convergence of measures and their completeness.
      \Cref{sec:ch48-measures}.
\item\label{itm:ch48-task:ma}
      \emph{Continuity and the homeomorphism property of
      \(\MA_\omega\).}  Prove that \(\MA_\omega\) is a bi-Lipschitz
      embedding of \((\cE^1(\omega)/\mathbf R,I)\) into
      \((\cM^1(\omega),I^\vee)\) with strongly dense image, and that under
      the envelope property it is onto, hence a homeomorphism of the two
      strong spaces.  This is the payoff of
      \Cref{chap:bj-energy,chap:bj-finite-energy} and the main theorem of
      the chapter.  \Cref{sec:ch48-ma}.
\item\label{itm:ch48-task:capacity}
      \emph{Capacity, negligible sets, polar sets and pluripolar sets.}
      Construct \(\Capa_\omega\) and \(\Capa^*_\omega\), prove the
      Chern--Levine--Nirenberg estimate, the sublevel bound, the energy
      criterion, the fundamental capacity estimate, quasicontinuity, and
      the Alexander--Taylor inequality in the cases available; relate the
      three notions of smallness.
      \Cref{sec:ch48-capacity,sec:ch48-quasicontinuity,sec:ch48-small-sets}.
\item\label{itm:ch48-task:compactness}
      \emph{Compactness and convergence criteria.}  Assemble the criteria
      by which strong convergence is verified in practice, and prove that
      strong compactness fails.  \Cref{sec:ch48-compactness}.
\end{enumerate}

\subsection{Hypotheses, normalization, and notation}
\label{sec:ch48-hypotheses}

The hypotheses are those of \Cref{chap:bj-energy,chap:bj-finite-energy}.
Thus \(k\) is algebraically closed and trivially valued, \(X\) is an
integral projective \(k\)-variety of dimension \(n\) which is
\emph{unibranch}, \(\omega\in\Amp(X)\), \(V=(\omega^n)>0\), and
\[
 \MA_\omega(\phi)=V^{-1}(\omega+\ddc\phi)^n
\]
is a \emph{probability} measure, the unnormalized measure of mass \(V\)
being written \(\MA^{\mathrm u}_\omega\).  We write \(\cE^1:=\cE^1(\omega)\),
\(\cM:=\{\hbox{Radon probability measures on }X^{\An}\}\),
\(\cM^1:=\cM^1(\omega)\), and we keep the normalized classes
\eqref{eq:ch47-normalized-classes}:
\[
 \begin{aligned}
  \cE^1_0(\omega)
   &=\{\phi\in\cE^1(\omega)\mid\textstyle\sup_{X^{\An}}\phi=0\},\\
  \PSH^{\sup}(\omega)
   &=\{\phi\in\PSH(\omega)\mid\textstyle\sup_{X^{\An}}\phi=0\}.
 \end{aligned}
\]
The functionals \(J_\psi\), \(I\) and \(J\) are those of
\eqref{eq:bj-I-J-potentials} and \eqref{eq:ch47-J-functional}, the
quasi-metric \(\mathbb I\) is \eqref{eq:bj-I-normalized}, the dual energy
\(E^\vee_\omega\) is \eqref{eq:bj-dual-energy}, and \(C_n\) always denotes
the constant of the quasi-triangle inequality
\eqref{eq:bj-I-quasitriangle}.

We additionally assume the envelope property whenever a result uses
obstacle envelopes, quasicontinuity, or the equality between negligible and
pluripolar sets.  By
\Crefrange{chap:bj-envelope-statement}{chap:bj-envelope-consequences} this
assumption holds when \(X\) is unibranch, which is our standing hypothesis;
\Cref{sec:ch48-hypothesis-use} audits the individual uses.

\begin{warning}{The probability normalization is inherited, not redeclared}
\label{warn:ch48-normalization}
\Cref{warn:ch46-normalization} declared a departure from the global
convention of the front matter, and that declaration is in force for the
whole of \Crefrange{chap:bj-energy}{chap:bj-stability}.  It is not repeated
or amended here.  Three consequences are used constantly.  The energy
translates by \(E_\omega(\phi+c)=E_\omega(\phi)+c\), not by \(cV\); the
dual energy is finite only on probability measures
(\Cref{lem:ch47-mass-forced}); and the capacity
\eqref{eq:bj-bt-capacity} takes values in \([0,1]\), not in \([0,V]\).  A
reader importing a numerical estimate from
\Cref{chap:bfj-capacity} must insert the factor \(V\) by hand; the
dictionary \eqref{eq:ch28-probability-dictionary} records the conversion on
the measure side, and \Cref{sec:ch48-part3} does it for the capacity.
\end{warning}

\subsection{Extended-real ledger}
\label{sec:ch48-ledger}

\Cref{sec:ch46-ledger} and \Cref{sec:ch47-ledger} listed five places where
a value \(-\infty\) is manipulated.  Two more appear here, and they are the
only new ones.

\begin{enumerate}
\item\label{itm:ch48-ledger:cap}
      In \eqref{eq:bj-alexander-taylor} we use the convention
      \((+\infty)^{-1}=0\), and in \eqref{eq:bj-alexander-taylor-height} we
      use \(\sup_\varnothing\phi=-\infty\).  Both are harmless because
      \(\sup_{X^{\An}}\phi\) is a real number for every
      \(\phi\in\PSH(\omega)\)
      (\Cref{lem:ch47-three-normalizations}\,\cref{itm:ch47-norm:max}), so
      that the difference in \eqref{eq:bj-alexander-taylor-height} is never
      of the form \((+\infty)-(+\infty)\); it may be \(+\infty\), and then
      the height is \(+\infty\).
\item\label{itm:ch48-ledger:sup}
      In \eqref{eq:bj-negligible} the supremum \(\sup_i\phi_i\) is the
      \emph{raw} pointwise supremum, computed in \(\Rbar\), and
      \((\sup_i\phi_i)^*\) is its least upper semicontinuous majorant.
      The two are compared pointwise, and their difference is never formed;
      the set in \eqref{eq:bj-negligible} is defined by a strict
      inequality between extended reals, which is meaningful whatever the
      values.
\end{enumerate}

The prohibition on \((-\infty)-(-\infty)\) remains binding.  In particular
every difference \(\phi-\psi\) of members of \(\cE^1(\omega)\) is formed on
\(X^{\lin}\), where both are finite by
\eqref{eq:bj-e1-sublinear-growth}, and every mixed measure occurring below
is carried there by \Cref{prop:ch46-no-mass-outside-lin}.

\subsection{What is imported, and what is proved}
\label{sec:ch48-imported}

\begin{important}{The three imported statements}
\label{imp:ch48-imported}
Everything in this chapter is proved from
\Crefrange{chap:bj-fs-pl-psh}{chap:bj-finite-energy} together with the
four imports already listed in \Cref{imp:ch47-imported}, except for the
following three statements, each of which is quoted from \cite{BJ22} with
an exact reference and used only where indicated.
\begin{enumerate}
\item\label{itm:ch48-import:pairing}
      \emph{The pairing estimates of \cite[\S10]{BJ22}}, in two forms.
      \begin{enumerate}
      \item[(a)] \emph{The quantitative form}, \cite[Theorem~10.3]{BJ22}:
            the quantity \(\bigl|\int(\phi-\psi)\,\rd(\mu-\nu)\bigr|\) is
            bounded by a product of a positive power of \(I(\phi,\psi)\),
            a positive power of \(I^\vee(\mu,\nu)\), and factors depending
            only on \(n\), \(J(\phi)\), \(J(\psi)\),
            \(E^\vee_\omega(\mu)\) and \(E^\vee_\omega(\nu)\).
      \item[(b)] \emph{The fixed-measure form}, contained in
            \cite[Theorem~12.4]{BJ22}: if \((\phi_i)\) is a net in
            \(\cE^1_0(\omega)\) with \(J(\phi_i)\le C\) converging weakly
            to \(\phi\in\cE^1_0(\omega)\), then
            \(\int\phi_i\,\rd\mu\to\int\phi\,\rd\mu\) for every fixed
            \(\mu\in\cM^1(\omega)\).
      \end{enumerate}
      Together they convert convergence of the two gauges into convergence
      of the pairing, and they are used exactly three times: (b) in the
      proof of \Cref{thm:bj-strong-e1-quasimetric}, for the implication
      from strong convergence to \(I\to0\); (a) in the same proof for the
      converse implication, and in the proof of
      \Cref{thm:bj-strong-measures} for the quasi-triangle inequality of
      \(I^\vee\) and for the equivalence of its three criteria; and nowhere
      else.  When \(n=1\), form (b) is proved from scratch in
      \Cref{prop:ch48-curve-strong}, and form (a) follows from
      \Cref{prop:ch48-curve-Idual} and the Cauchy--Schwarz inequality of
      \Cref{prop:ch47-curve-I-norm}, with exponents \(\tfrac12\) in place
      of \(2^{-n}\).
\item\label{itm:ch48-import:increasing}
      \emph{Continuity of the energy along increasing nets},
      \cite[Theorem~12.5]{BJ22}: if \((\phi_i)\) is an increasing net in
      \(\cE^1(\omega)\) whose regularized supremum
      \(\phi=(\sup_i\phi_i)^*\) lies in \(\cE^1(\omega)\), then
      \(E_\omega(\phi_i)\to E_\omega(\phi)\).  The corresponding statement
      for decreasing nets is
      \Cref{prop:ch46-energy-usc}\,\cref{itm:ch46-Eusc:decreasing} and is
      proved.  The increasing statement is used for the increasing half of
      \Cref{prop:bj-strong-monotone}, for the completeness theorem
      \Cref{thm:ch48-completeness}, and for the discharge of the deferred
      step of \Cref{thm:bj-obstacle-orthogonality-energy} in
      \Cref{sec:ch48-discharge}.
\item\label{itm:ch48-import:negligible}
      \emph{Negligible sets have vanishing outer capacity},
      \cite[Theorem~13.17]{BJ22}, and the lower Alexander--Taylor bound for
      arbitrary subsets, \cite[Theorem~13.14]{BJ22}.  Both are proved here
      for compact sets (\Cref{prop:ch48-at-compact}), and both hold under
      the envelope property, which we have.  They are used only in the
      proof of \Cref{thm:bj-small-set-equivalence}, and
      \Cref{sec:ch48-what-is-open} explains exactly why the argument given
      here does not reach them.
\end{enumerate}
No other external input is used.  In particular the Calabi--Yau theorem
\Cref{cor:bj-cy-unibranch} is used, but it is proved in
\Cref{chap:bj-finite-energy} and is not an import.
\end{important}

The three imports are of different depths, and it is worth saying which is
which.  Item~\cref{itm:ch48-import:pairing} is a quantitative estimate of
Chern--Levine--Nirenberg type, comparable in nature to the fundamental
capacity estimate \Cref{lem:ch48-fundamental} that we do prove; the reason
we do not prove it is that its two-sided form couples the potential gauge
to the measure gauge, and the iteration of \Cref{lem:ch48-fundamental}
produces only the potential side.
Item~\cref{itm:ch48-import:increasing} is the increasing half of a monotone
convergence theorem whose decreasing half is elementary; the asymmetry is
structural, not technical, and \Cref{warn:ch48-monotone-asymmetry} explains
it.  Item~\cref{itm:ch48-import:negligible} is the deepest: it is the
non-Archimedean form of the Bedford--Taylor theorem that negligible sets
are pluripolar, and in the complex theory the corresponding statement rests
on Josefson's theorem \cite{Jos78}.

\begin{important}{The envelope theorem is used without further comment}
\label{imp:ch48-envelope}
As in \Cref{imp:ch46-envelope-imported} and \Cref{imp:ch47-envelope}, the
variety \(X\) is unibranch and \Cref{thm:bj-final-envelope} is available.
The consequences used here are:
\begin{enumerate}
\item\label{itm:ch48-env:operator}
      \(P_\omega\) is an operator from admissible obstacles
      (\Cref{def:ch45-admissible-obstacle}) to \(\PSH(\omega)\), is the
      identity on \(\PSH(\omega)\), and obeys the calculus
      \Cref{prop:bj-final-obstacle-calculus},
      \Cref{prop:ch45-obstacle-calculus-2};
\item\label{itm:ch48-env:compact}
      \(\PSH^{\sup}(\omega)\) is compact for the topology of pointwise
      convergence on \(X^{\divis}\) (\Cref{cor:ch45-compactness});
\item\label{itm:ch48-env:regularization}
      every \(u\in\PSH(\omega)\) is the limit of a downward directed net of
      \emph{continuous} \(\omega\)-psh functions
      (\Cref{thm:ch45-continuous-regularization}), and every increasing
      family has \(\omega\)-psh regularized supremum
      (\Cref{thm:ch45-increasing-obstacles});
\item\label{itm:ch48-env:orthogonality}
      orthogonality holds for continuous obstacles
      (\Cref{thm:ch47-orthogonality-continuous}) and, once
      \Cref{sec:ch48-discharge} has been read, for lower semicontinuous
      ones (\Cref{thm:bj-obstacle-orthogonality-energy}).
\end{enumerate}
Item \cref{itm:ch48-env:regularization} is the one that has no counterpart
in \Cref{part:bfj}, and it is why the quasicontinuity proof below is
shorter than \Cref{thm:bounded-quasicontinuity}.
\end{important}

\subsection{Relation to \texorpdfstring{\Cref{part:bfj}}{Part III}, and to Thuillier's curve theory}
\label{sec:ch48-part3}

This chapter is the trivially valued counterpart of
\Cref{chap:bfj-capacity,chap:bfj-energy,chap:bfj-finite-energy-measures,chap:bfj-continuity}.
The parallel is close, and the differences must be tabulated before they
are met.

\emph{What is the same.}  The capacity is defined by the same formula, a
supremum of Monge--Amp\`ere masses over potentials with values in
\([-1,0]\); compare \eqref{eq:bj-bt-capacity} with
\eqref{eq:capacity-definition}.  The Chern--Levine--Nirenberg estimate,
the sublevel bound, countable subadditivity, and the fundamental capacity
estimate \eqref{eq:fundamental-capacity-estimate} all transfer, and the
transfer is literal: the proofs use only multilinearity of mixed measures
on bounded psh functions, integration by parts, and positivity of the
Dirichlet form, all of which are \Cref{thm:bj-mixed-ma},
\Cref{prop:bj-integration-by-parts} and \Cref{cor:ch46-dirichlet-e1} here.
The only change is that \(V\) is replaced by \(1\) throughout, by
\Cref{warn:ch48-normalization}.  \Cref{lem:ch48-fundamental} below records
the transferred statement with its constants.

\emph{What is better here.}  Three things.  The regularization of an
arbitrary psh function is by continuous psh functions in the \emph{same}
class (\Cref{thm:ch45-continuous-regularization}), where over a discretely
valued base it is by model functions in varying classes; the proof of
quasicontinuity shortens correspondingly.  The compactness of
\(\PSH^{\sup}(\omega)\) is unconditional (\Cref{cor:ch45-compactness}),
where \Cref{thm:bfj-psh-compactness} is a theorem about a fixed model
system.  And the equivalence of pluripolarity with vanishing outer capacity
is available at all: \Cref{rem:ch25-curve-converse} records that over a
discretely valued base the converse to
\Cref{cor:ch25-polar-negligible} is not proved in this book.

\emph{What is worse here.}  One thing, and it is a real one.  Over a
discretely valued base, the Monge--Amp\`ere measure of a bounded potential
is carried by a dual complex, a compact metrizable space, and
\Cref{prop:capacity-subsequence} together with
\Cref{cor:psh-sequential-regularization} makes sequential arguments
available.  Here \(X^{\An}\) is a cone, the divisorial locus is dense but
carries no cofinal sequence of models, and every compactness statement is a
statement about nets.  The discipline of \Cref{sec:ch25-nets} is therefore
in force verbatim, and \Cref{sec:ch48-nets} says where it bites.

\emph{Thuillier's curve capacity.}  \Cref{sec:ch25-thuillier} gives a
complete account of the relative capacity \(C(\cdot,\Omega)\) of
\cite{Thu05} on a Berkovich curve, and
\Cref{thm:ch25-thuillier-polar} states its central theorem: on a connected
open \(\Omega\) with nonempty boundary in the quasimonomial locus, a subset
is polar if and only if its outer relative capacity vanishes, if and only
if its relative extremal function is identically \(0\), if and only if it
is swallowed by a single subharmonic function.  Two remarks are needed
here, and both are consequences of the difference in normalization recorded
in \Cref{rem:ch25-thuillier-vs-bfj}.

The first is that \emph{Thuillier's theorem does not transfer}, because it
is not a statement about \(\Capa_\omega\).  His capacity is an inverse
minimal energy relative to an open set with boundary, and the boundary is
what makes his Dirichlet problem solvable; the capacity
\eqref{eq:bj-bt-capacity} is global on \(X^{\An}\), normalized by the class
\(\omega\) and by the constraint \(-1\le\psi\le0\), and there is no
Dirichlet problem in sight.  Over a trivially valued field the difficulty
is aggravated, not relieved: the trivially valued analytification of a
curve is a cone, and the natural candidate for \(\Omega\) --- the
complement of the cone point --- is disconnected with uncountably many
components, one for each closed point.  The countability clause in
\cite[Th\'eor\`eme~3.6.11]{Thu05} is therefore violated in the most
flagrant way, and \Cref{ex:ch48-proper-support-capacity} below turns that
observation into a computation: an uncountable set of proper-support points
has inner capacity \(0\) and outer capacity \(1\).

The second is that \emph{one structural lesson does transfer}, and it is
the one \Cref{rem:ch25-thuillier-vs-bfj} isolated.  Thuillier's development
uses \emph{no quasicontinuity of subharmonic functions}: he follows
Klimek's presentation \cite{Kl91} of the complex theory, simplified by the
fact that on a curve the Dirichlet problem can be solved explicitly, and
constructs equilibrium potentials directly.  Here quasicontinuity is again
the central technical result, as it was in \Cref{chap:bfj-capacity},
because in dimension \(\ge2\) the only regularization available is by
decreasing nets.  What we do inherit from his account is the discipline
about countability: a union of uncountably many polar sets need not be
polar, an uncountable family of components is not a sequence, and the
counterexample \cite[Remarque~3.6.12]{Thu05} has an exact analogue in the
trivially valued setting, which is \Cref{ex:ch48-proper-support-capacity}.
The reader who has understood why Thuillier's ``polar'' carries a
countability requirement will not be surprised by that example.

\subsection{Roadmap}
\label{sec:ch48-roadmap}

\Cref{sec:ch48-potentials} constructs the strong topology on
\(\cE^1(\omega)\), identifies it with the \(\mathbb I\)-topology, proves the
monotone criteria, and states completeness.
\Cref{sec:ch48-measures} does the same on the measure side.
\Cref{sec:ch48-ma} proves the main theorem: \(\MA_\omega\) is a
bi-Lipschitz homeomorphism of the two strong spaces.
\Cref{sec:ch48-capacity} constructs the capacity and proves the estimates.
\Cref{sec:ch48-quasicontinuity} proves quasicontinuity and its
consequences, and discharges the step deferred from
\Cref{thm:bj-obstacle-orthogonality-energy}.
\Cref{sec:ch48-small-sets} relates negligible, polar and pluripolar sets.
\Cref{sec:ch48-compactness} assembles the convergence criteria.
\Cref{sec:ch48-examples} computes: the motivating example of
\Cref{chap:bj-energy} with all its energies, a capacity on a
\(\mathbf P^1\) skeleton, a pluripolar set with its capacity, and an
explicit failure of the triangle inequality for \(\mathbb I\).
\Cref{sec:ch48-scope} audits the hypotheses,
\Cref{sec:ch48-dictionary} is the chapter dictionary, and
\Cref{sec:ch48-source} is the source note.

\section{The strong topology on potentials}
\label{sec:ch48-potentials}

\subsection{Two topologies on \texorpdfstring{\(\cE^1(\omega)\)}{E1}}
\label{sec:ch48-two-topologies}

\begin{definition}[Weak and strong potential topologies]
\label{def:bj-strong-e1}
The weak topology on \(\mathcal E^1(\omega)\) is inherited from
\(\PSH(\omega)\): a net \(\phi_i\to\phi\) weakly exactly when

\[
                       \phi_i(v)\longrightarrow\phi(v)
          \quad\text{for every }v\in X^{\divis}.
\tag{48.1}\label{eq:bj-e1-weak}
\]

The strong topology is the coarsest refinement of the weak topology for
which \(E:\mathcal E^1\to\mathbf R\) is continuous.
\end{definition}

Thus a net converges strongly if and only if it converges weakly and its
energies converge.  Pointwise convergence here is never required at a
proper-support point, where a psh function may equal \(-\infty\).  The
following lemma makes the last two sentences into statements with proofs,
and records the two soft properties of the strong topology that are used
without comment below.

\begin{lemma}
\label{lem:ch48-net-criterion}
Let \(\tau_w\) denote the weak topology on \(\cE^1(\omega)\) and \(\tau_s\)
the strong one.  Then:
\begin{enumerate}
\item\label{itm:ch48-net:subbasis}
      \(\tau_s\) is generated by the sets of \(\tau_w\) together with the
      sets \(\{E_\omega>c\}\), \(c\in\mathbf R\);
\item\label{itm:ch48-net:criterion}
      a net \((\phi_i)\) in \(\cE^1(\omega)\) converges to \(\phi\) for
      \(\tau_s\) if and only if \(\phi_i\to\phi\) for \(\tau_w\) and
      \(E_\omega(\phi_i)\to E_\omega(\phi)\);
\item\label{itm:ch48-net:hausdorff}
      \(\tau_w\), hence also \(\tau_s\), is Hausdorff.
\end{enumerate}
\end{lemma}

\begin{proof}
\cref{itm:ch48-net:subbasis}.  A topology refining \(\tau_w\) makes
\(E_\omega\) continuous if and only if it contains all the sets
\(\{E_\omega>c\}\) and \(\{E_\omega<c\}\).  The latter are already
\(\tau_w\)-open, because \(E_\omega\) is upper semicontinuous for \(\tau_w\)
by \Cref{prop:ch46-energy-usc}\,\cref{itm:ch46-Eusc:props}.  Hence the
coarsest such refinement is the topology generated by
\(\tau_w\cup\{\{E_\omega>c\}\mid c\in\mathbf R\}\).

\cref{itm:ch48-net:criterion}.  Suppose \(\phi_i\to\phi\) for \(\tau_w\) and
\(E_\omega(\phi_i)\to E_\omega(\phi)\).  A basic \(\tau_s\)-neighbourhood of
\(\phi\) is a finite intersection of \(\tau_w\)-neighbourhoods and sets
\(\{E_\omega>c\}\) with \(c<E_\omega(\phi)\); the net is eventually in each
of the finitely many members, hence eventually in the intersection.
Conversely, if \(\phi_i\to\phi\) for \(\tau_s\) then \(\phi_i\to\phi\) for
the coarser \(\tau_w\); and for \(c<E_\omega(\phi)\) the set
\(\{E_\omega>c\}\) is a \(\tau_s\)-neighbourhood of \(\phi\), so
eventually \(E_\omega(\phi_i)>c\), whence
\(\varliminf_iE_\omega(\phi_i)\ge E_\omega(\phi)\).  Upper semicontinuity
for \(\tau_w\) gives \(\varlimsup_iE_\omega(\phi_i)\le E_\omega(\phi)\), and
the two combine.

\cref{itm:ch48-net:hausdorff}.  By the last sentence of
\Cref{thm:triv-divisorial-density}, two \(\omega\)-psh functions that agree
on \(X^{\divis}\) agree on \(X^{\An}\); so the evaluation maps
\(\phi\mapsto\phi(v)\), \(v\in X^{\divis}\), separate points of
\(\PSH(\omega)\), and \(\tau_w\) is the initial topology they define, hence
Hausdorff.  A refinement of a Hausdorff topology is Hausdorff.
\end{proof}

\begin{remark}[Why upper semicontinuity is what makes the definition work]
\label{rem:ch48-usc-is-the-point}
The proof of \cref{itm:ch48-net:subbasis} used upper semicontinuity of
\(E_\omega\) and nothing else, and it is worth noticing how much that
buys.  Had \(E_\omega\) been merely Borel, the coarsest refinement making
it continuous would have been generated by all its sublevel and superlevel
sets, and the criterion \cref{itm:ch48-net:criterion} would have read
``\(\phi_i\to\phi\) weakly and \(E_\omega(\phi_i)\to E_\omega(\phi)\)''
only after an additional argument.  Because half of the continuity is free,
the strong topology is a \emph{one-sided} strengthening of the weak one:
it forbids energy from being lost in the limit, and nothing else.  The
same observation, in the same words, is
\Cref{prop:ch26-strong-basic} over a discretely valued base.
\end{remark}

\subsection{Identification with the quasi-metric \texorpdfstring{\(\mathbb I\)}{I}}
\label{sec:ch48-quasimetric-identification}

\begin{theorem}\label{thm:bj-strong-e1-quasimetric}
The strong topology is defined by the quasi-metric

\[
 \mathbb I(\phi,\psi)
   =I(\phi,\psi)+|\sup\phi-\sup\psi|.
\tag{48.2}\label{eq:bj-strong-e1-I}
\]

Equivalently, \(\phi_i\to\phi\) strongly if and only if
\(\mathbb I(\phi_i,\phi)\to0\).
\end{theorem}

\begin{proof}
That \(\mathbb I\) is a symmetric, separating gauge satisfying the
quasi-triangle inequality with constant \(\max\{C_n,1\}\) is
\Cref{prop:ch47-quasimetric}, and is not reproved.  What has to be shown is
the equivalence of the two convergence notions.  We give the two
implications separately, marking the single place where
\Cref{imp:ch48-imported}\,\cref{itm:ch48-import:pairing} is used, and then,
in \Cref{sec:ch48-curve-case}, a complete proof for \(n=1\) that uses no
import at all.

\emph{Strong convergence implies \(\mathbb I(\phi_i,\phi)\to0\).}  Suppose
\(\phi_i\to\phi\) weakly and \(E_\omega(\phi_i)\to E_\omega(\phi)\).  The
trivial valuation \(v_{\triv}\) is divisorial --- it is the valuation
attached to the central fibre of the trivial test configuration, as used in
\Cref{lem:ch47-three-normalizations}\,\cref{itm:ch47-norm:trivial-measure}
--- so \eqref{eq:bj-e1-weak} applies at \(v_{\triv}\) and, by
\Cref{lem:ch47-three-normalizations}\,\cref{itm:ch47-norm:max},
\[
 \sup_{X^{\An}}\phi_i=\phi_i(v_{\triv})\longrightarrow
 \phi(v_{\triv})=\sup_{X^{\An}}\phi .
\]
Hence the second term of \eqref{eq:bj-strong-e1-I} tends to \(0\), and
\(J(\phi_i)=\sup\phi_i-E_\omega(\phi_i)\to J(\phi)\); in particular
\(J(\phi_i)\le C\) for some \(C\) and all large \(i\).  For the first term,
\eqref{eq:bj-I-J-comparison} gives
\(I(\phi_i,\phi)\le(n+1)J_\phi(\phi_i)\) and
\[
 J_\phi(\phi_i)=E_\omega(\phi)-E_\omega(\phi_i)
 +\int(\phi_i-\phi)\,\MA_\omega(\phi),
\]
whose first two terms tend to \(0\).  The remaining integral is a pairing
of the difference \(\phi_i-\phi\) against the \emph{fixed} finite-energy
measure \(\mu:=\MA_\omega(\phi)\in\cM^1(\omega)\)
(\Cref{prop:bj-fenchel-equalities}).  Replacing \(\phi_i\) by
\(\phi_i-\sup\phi_i\) and \(\phi\) by \(\phi-\sup\phi\), which changes
neither \(I\) nor the convergence just established, we may assume that all
the potentials lie in \(\cE^1_0(\omega)\) with \(J\le C\); then
\(\int\phi_i\,\rd\mu\to\int\phi\,\rd\mu\) is exactly form (b) of
\Cref{imp:ch48-imported}\,\cref{itm:ch48-import:pairing}.  All the energy
parameters occurring there are bounded along the net: \(J(\phi_i)\le C\) as
just observed, and
\[
 E^\vee_\omega\bigl(\MA_\omega(\chi)\bigr)
 =J_\chi(0)\le\tfrac n{n+1}I(\chi,0)\le n\,J(\chi)
\]
for every \(\chi\in\cE^1(\omega)\), by \eqref{eq:bj-dual-energy-ma} and
\eqref{eq:bj-I-J-comparison}.

\emph{Conversely, \(\mathbb I(\phi_i,\phi)\to0\) implies strong
convergence.}  The supremum term fixes the otherwise invisible
translations: \(\sup\phi_i\to\sup\phi\).  After this normalization, the
estimate of \cite[Theorem~10.3]{BJ22}, applied to
\(\delta_v,\delta_{v_{\triv}}\in\cM^1(\omega)\) --- both are finite-energy
measures when \(v\in X^{\divis}\), by
\Cref{thm:bj-dirac-finite-energy} and
\Cref{prop:bj-linear-growth-potential}, since \(X^{\divis}\subset X^{\lin}\)
--- shows that \(I(\phi_i,\phi)\to0\) controls the evaluation at every
divisorial point.  Thus \(\phi_i\to\phi\) weakly.  The quasi-triangle
inequality \eqref{eq:bj-I-quasitriangle} also bounds \(J(\phi_i)\), since
\(J(\phi_i)\le\frac n{n+1}I(\phi_i,\sup\phi_i)=\frac n{n+1}I(\phi_i,0)
\le\frac{nC_n}{n+1}\bigl(I(\phi_i,\phi)+I(\phi,0)\bigr)\)
by \Cref{cor:ch47-J-vanishes} and
\Cref{lem:ch47-IJ-elementary}\,\cref{itm:ch47-IJ:translation}; so the same
integral identity yields
\(E_\omega(\phi_i)\to E_\omega(\phi)\).  This is precisely
\cite[Theorem~12.4]{BJ22}.
\end{proof}

\begin{example}[Why the supremum term is necessary]
\label{ex:ch48-sup-term}
Fix \(\phi\in\mathcal E^1(\omega)\) and let \(c_i\in\mathbf R\) tend to
zero.  Put \(\phi_i=\phi+c_i\).  The Monge--Amp\`ere measure is unchanged,
and the probability normalization gives
\[
        E(\phi_i)=E(\phi)+c_i,
        \qquad \sup\phi_i=\sup\phi+c_i.
\]
Since \(\phi_i-\phi\) is constant,
\Cref{prop:bj-energy-I-identities} yields
\(I(\phi_i,\phi)=0\).  Therefore
\[
             \mathbb I(\phi_i,\phi)=|c_i|\longrightarrow0.
\]
This calculation separates the two roles in
\eqref{eq:bj-strong-e1-I}: \(I\) measures the class modulo constants, while
the supremum term restores the missing normalization on
\(\mathcal E^1\) itself.

The example also shows that the term cannot be replaced by
\(|E_\omega(\phi)-E_\omega(\psi)|\) without changing the statement of
\Cref{thm:bj-strong-e1-quasimetric}: for \(\phi_i=\phi+c_i\) the two
expressions agree, but for a net with \(\sup\phi_i=\sup\phi=0\) the
supremum term vanishes identically while
\(|E_\omega(\phi_i)-E_\omega(\phi)|=|J(\phi_i)-J(\phi)|\) need not.  The
gauge \(I+|E-E|\) is therefore a different, and strictly larger, gauge on
\(\cE^1(\omega)\); it induces the same topology, because
\Cref{thm:bj-strong-e1-quasimetric} identifies both with the strong one,
but it is not \eqref{eq:bj-strong-e1-I}.  We use
\eqref{eq:bj-strong-e1-I} because its first term is translation invariant
and its second is an evaluation at a divisorial point, which is what makes
\Cref{prop:ch47-scaling}\,\cref{itm:ch47-scal:quasi} --- homogeneity of
\(\mathbb I\) under the scaling action --- true on the nose.
\end{example}

\subsection{The curve case, in full}
\label{sec:ch48-curve-case}

When \(n=1\) the Dirichlet form ceases to depend on the potentials, and
\Cref{thm:bj-strong-e1-quasimetric} becomes a statement about a
pre-Hilbert space.  We prove it in that case with no import whatever.  The
argument is short, and it is the model for what the imported estimate
\Cref{imp:ch48-imported}\,\cref{itm:ch48-import:pairing} does in general.

Throughout this subsection \(n=\dim X=1\).  By
\Cref{prop:ch47-curve-I-norm} there is a symmetric positive semidefinite
bilinear form
\[
 \langle f,g\rangle:=-V^{-1}\int_{X^{\An}}f\,\ddc g
\]
on \(\vec\cE^{\,1}(X)/\mathbf R\), independent of any potential, normalized so
that
\[
 I(\phi,\psi)=\lVert\phi-\psi\rVert^2,
 \qquad
 \lVert f\rVert:=\langle f,f\rangle^{1/2},
\]
and, \(\ddc\) being linear and \(\MA_\omega\) being
\(V^{-1}(\omega+\ddc\,\cdot\,)\), one has for all
\(a,b,\chi\in\cE^1(\omega)\)
\[
 \int_{X^{\An}}\chi\;\rd\bigl(\MA_\omega(a)-\MA_\omega(b)\bigr)
 =-\langle\chi,a-b\rangle .
\]
We write \(\lVert\phi\rVert^2=I(\phi,0)=2J(\phi)\), using
\Cref{cor:ch47-J-vanishes} and the fact that the two constants in
\eqref{eq:bj-I-J-comparison} both equal \(\tfrac12\) when \(n=1\).

\begin{proposition}
\label{prop:ch48-curve-strong}
Let \(\dim X=1\) and let \((\phi_i)\) be a net in \(\cE^1(\omega)\) and
\(\phi\in\cE^1(\omega)\).  Then \(\mathbb I(\phi_i,\phi)\to0\) if and only
if \(\phi_i\to\phi\) pointwise on \(X^{\divis}\) and
\(E_\omega(\phi_i)\to E_\omega(\phi)\).
\end{proposition}

\begin{proof}
\emph{Suppose \(\mathbb I(\phi_i,\phi)\to0\).}  Then
\(\sup\phi_i\to\sup\phi\) and \(\lVert\phi_i-\phi\rVert\to0\).  Since
\(\lVert\,\cdot\,\rVert\) is a seminorm,
\(\bigl|\lVert\phi_i\rVert-\lVert\phi\rVert\bigr|
\le\lVert\phi_i-\phi\rVert\to0\), so
\(2J(\phi_i)=\lVert\phi_i\rVert^2\to\lVert\phi\rVert^2=2J(\phi)\) and
therefore
\[
 E_\omega(\phi_i)=\sup\phi_i-J(\phi_i)
 \longrightarrow\sup\phi-J(\phi)=E_\omega(\phi).
\]
For pointwise convergence let \(v\in X^{\divis}\subset X^{\lin}\).  By
\Cref{cor:ch47-divisorial-green} there is \(\phi_v\in\cE^1(\omega)\) with
\(\MA_\omega(\phi_v)=\delta_v\).  Applying the displayed pairing identity
with \(a=\phi_v\), \(b=0\), \(\chi=\phi_i-\phi\), and using
\(\MA_\omega(0)=\delta_{v_{\triv}}\), we get
\[
 \bigl(\phi_i(v)-\phi(v)\bigr)-\bigl(\phi_i(v_{\triv})-\phi(v_{\triv})\bigr)
 =-\langle\phi_i-\phi,\phi_v\rangle,
\]
whence, by Cauchy--Schwarz and
\Cref{lem:ch47-three-normalizations}\,\cref{itm:ch47-norm:max},
\[
 \begin{aligned}
 |\phi_i(v)-\phi(v)|
 &\le|\sup\phi_i-\sup\phi|
   +\lVert\phi_i-\phi\rVert\,\lVert\phi_v\rVert\\
 &=|\sup\phi_i-\sup\phi|
   +I(\phi_i,\phi)^{1/2}\bigl(2J(\phi_v)\bigr)^{1/2},
 \end{aligned}
\]
which tends to \(0\).  Note that \(J(\phi_v)<\infty\) precisely because
\(\delta_v\in\cM^1(\omega)\), that is, because \(v\) has linear growth
(\Cref{thm:bj-dirac-finite-energy}); the estimate degenerates as \(v\)
leaves a \(T_\omega\)-ball, which is the one-dimensional shadow of
\eqref{eq:bj-e1-sublinear-growth}.

\emph{Conversely, suppose \(\phi_i\to\phi\) pointwise on \(X^{\divis}\) and
\(E_\omega(\phi_i)\to E_\omega(\phi)\).}  Evaluating at \(v_{\triv}\) gives
\(\sup\phi_i\to\sup\phi\), so \(J(\phi_i)\to J(\phi)\), that is
\(\lVert\phi_i\rVert\to\lVert\phi\rVert\); in particular
\(\lVert\phi_i\rVert^2=2J(\phi_i)\le2C\) for some \(C\) and all large
\(i\).  Expanding,
\[
 I(\phi_i,\phi)=\lVert\phi_i-\phi\rVert^2
 =\lVert\phi_i\rVert^2+\lVert\phi\rVert^2-2\langle\phi_i,\phi\rangle,
\]
so it suffices to prove \(\langle\phi_i,\phi\rangle\to\lVert\phi\rVert^2\).
By the pairing identity with \(a=\phi\), \(b=0\),
\[
 \langle\chi,\phi\rangle
 =\chi(v_{\triv})-\int\chi\,\MA_\omega(\phi)
 \qquad(\chi\in\cE^1(\omega)),
\]
and the first term converges as required, so everything reduces to
\(\int\phi_i\,\MA_\omega(\phi)\to\int\phi\,\MA_\omega(\phi)\).

Fix \(\varepsilon>0\).  By \Cref{thm:ch45-continuous-regularization} and
\Cref{thm:ch34-pl-density} --- or directly by
\Cref{def:bj-psh}\,(3) followed by
\Cref{cor:ch45-continuity-envelopes} --- choose
\(\chi\in\PL^\omega(X)\) with \(\chi\ge\phi\) and
\(I(\chi,\phi)\le\varepsilon^2\); this is possible because a decreasing net
of PL \(\omega\)-psh functions converging to \(\phi\) has
\(E_\omega\to E_\omega(\phi)\) by
\Cref{prop:ch46-energy-usc}\,\cref{itm:ch46-Eusc:decreasing}, and then
\(I\to0\) by \Cref{cor:ch47-monotone-I}.  Now
\[
 \int\phi_i\,\MA_\omega(\phi)
 =\int\phi_i\,\MA_\omega(\chi)
  +\int\phi_i\,\rd\bigl(\MA_\omega(\phi)-\MA_\omega(\chi)\bigr),
\]
and the second term is
\(\langle\phi_i,\chi-\phi\rangle\) in absolute value at most
\(\lVert\phi_i\rVert\,\lVert\chi-\phi\rVert\le(2C)^{1/2}\varepsilon\),
\emph{uniformly in \(i\)}.  The measure \(\MA_\omega(\chi)\) is a finite
positive combination of Dirac masses at divisorial points, by
\Cref{prop:ch46-pl-measure}, so
\(\int\phi_i\,\MA_\omega(\chi)\to\int\phi\,\MA_\omega(\chi)\) by
hypothesis.  The same estimate bounds
\(|\int\phi\,\rd(\MA_\omega(\phi)-\MA_\omega(\chi))|\) by
\(\lVert\phi\rVert\varepsilon\).  Hence
\[
 \varlimsup_i\Bigl|\int\phi_i\,\MA_\omega(\phi)
   -\int\phi\,\MA_\omega(\phi)\Bigr|
 \le\bigl((2C)^{1/2}+\lVert\phi\rVert\bigr)\varepsilon,
\]
and \(\varepsilon\) was arbitrary.
\end{proof}

\begin{remark}[What the curve proof isolates]
\label{rem:ch48-curve-isolates}
Two mechanisms appear in \Cref{prop:ch48-curve-strong}, and both are
present in the general proof of \cite[Theorem~12.4]{BJ22}, where they are
harder.

The first is the \emph{uniform} bound
\(|\int\chi\,\rd(\mu-\nu)|\le\lVert\chi\rVert\,\lVert a-b\rVert\), valid for
all \(\chi\) in a fixed \(J\)-sublevel; it is what allows the limit in
\(i\) and the approximation of \(\MA_\omega(\phi)\) by a finitely supported
measure to be interchanged.  In general the analogous statement couples the
potential gauge \(I\) to the measure gauge \(I^\vee\), with exponents
\(2^{-n}\) rather than \(\tfrac12\), and is
\Cref{imp:ch48-imported}\,\cref{itm:ch48-import:pairing}.

The second is that \emph{Monge--Amp\`ere measures of PL potentials, hence
finitely supported divisorial measures, are strongly dense}.  Indeed,
finite-energy potentials admit decreasing PL approximants by
\Cref{lem:ch46-pl-domination}; these approximants converge strongly by
\Cref{prop:ch46-energy-usc}, their measures converge strongly by
\Cref{prop:bj-ma-strong-embedding}, and the measures are finitely supported
and divisorial by \Cref{prop:ch46-pl-measure}.  Both proofs then run identically:
approximate the fixed measure, use weak convergence against a finite atomic
measure, and control the error uniformly.  A reader who wants to see why
\Cref{thm:bj-strong-e1-quasimetric} is a theorem and not a formality should
compare the two error terms.
\end{remark}

\subsection{Monotone and uniform criteria}
\label{sec:ch48-monotone}

\begin{proposition}\label{prop:bj-strong-monotone}
Let \((\phi_i)\) be a monotone net in \(\mathcal E^1\) with weak limit
\(\phi\in\mathcal E^1\).  Then \(\phi_i\to\phi\) strongly.  If
\(\phi_i\leq\phi\) for all \(i\), strong convergence is equivalent to
\(E(\phi_i)\to E(\phi)\).
\end{proposition}

\begin{proof}
For a decreasing net, this is continuity of energy along decreasing nets.
For an increasing net, BJ's mixed-energy continuity theorem gives energy
convergence; this is \cite[Theorem~12.5]{BJ22}.  Under
\(\phi_i\leq\phi\), the tangent inequality gives

\[
 0\leq J_\phi(\phi_i)
 \leq E(\phi)-E(\phi_i),
\]

which proves the last assertion by
\Cref{thm:bj-strong-e1-quasimetric}.

In detail.  Suppose first that \((\phi_i)\) is decreasing with weak limit
\(\phi\).  Then \(\phi=\inf_i\phi_i\) pointwise on \(X^{\divis}\), and
since a decreasing net of upper semicontinuous functions has upper
semicontinuous infimum, and two \(\omega\)-psh functions agreeing on
\(X^{\divis}\) agree everywhere
(\Cref{thm:triv-divisorial-density}), \(\phi_i\downarrow\phi\) pointwise on
\(X^{\An}\).  By
\Cref{prop:ch46-energy-usc}\,\cref{itm:ch46-Eusc:decreasing},
\(E_\omega(\phi_i)\to E_\omega(\phi)\), and
\Cref{lem:ch48-net-criterion}\,\cref{itm:ch48-net:criterion} gives strong
convergence.  Suppose next that \((\phi_i)\) is increasing.  Then
\(\sup_i\phi_i\le\phi\) with equality on \(X^{\divis}\), hence
\(\phi=(\sup_i\phi_i)^*\) by
\Cref{thm:ch45-increasing-obstacles}\,\cref{itm:ch45-inc:star} and
\Cref{prop:ch45-obstacle-calculus-2}\,\cref{itm:ch45-P:divisorial}; energy
convergence is
\Cref{imp:ch48-imported}\,\cref{itm:ch48-import:increasing}.

For the last assertion, \(\phi_i\le\phi\) gives
\(J_\phi(\phi_i)\le E_\omega(\phi)-E_\omega(\phi_i)\) by
\Cref{lem:ch47-IJ-elementary}\,\cref{itm:ch47-IJ:monotone}, hence
\(I(\phi_i,\phi)\le(n+1)\bigl(E_\omega(\phi)-E_\omega(\phi_i)\bigr)\) by
\eqref{eq:bj-I-J-comparison}; this is \Cref{cor:ch47-monotone-I}.  If in
addition \(\phi_i\to\phi\) weakly, then
\(\sup\phi_i\to\sup\phi\) as in the proof of
\Cref{thm:bj-strong-e1-quasimetric}, so
\(\mathbb I(\phi_i,\phi)\to0\) as soon as
\(E_\omega(\phi_i)\to E_\omega(\phi)\); the converse implication is
\Cref{lem:ch48-net-criterion}\,\cref{itm:ch48-net:criterion}.
\end{proof}

\begin{warning}{The two monotone directions are not symmetric}
\label{warn:ch48-monotone-asymmetry}
Continuity of \(E_\omega\) along \emph{decreasing} nets is
\Cref{prop:ch46-energy-usc}\,\cref{itm:ch46-Eusc:decreasing}, proved in
\Cref{chap:bj-energy} from the representation
\Cref{prop:ch46-energy-usc}\,\cref{itm:ch46-Eusc:rep} of the energy as an
infimum over dominating PL functions: an infimum of continuous functions of
the argument is upper semicontinuous, and along a decreasing net the
infimum is attained in the limit.  The corresponding statement for
\emph{increasing} nets is
\Cref{imp:ch48-imported}\,\cref{itm:ch48-import:increasing} and is
\emph{not} proved here.  What monotonicity gives for free is only
\(\sup_iE_\omega(\phi_i)\le E_\omega\bigl((\sup_i\phi_i)^*\bigr)\), and
upper semicontinuity of \(E_\omega\) points in the same direction, so it is
of no help.  The same asymmetry, with the same explanation, is recorded
over a discretely valued base after
\Cref{prop:ch26-strong-basic}; the difference is that there the increasing
statement is not available at all, whereas here it is a theorem of
\cite{BJ22} whose hypothesis --- the envelope property --- we have.

The asymmetry is not an artefact of the proofs.  An increasing net can lose
its regularized supremum on a negligible set, and it is exactly the
smallness of that set --- \Cref{thm:bj-small-set-equivalence} --- that
makes the energies converge.  This is why
\Cref{sec:ch48-small-sets} is logically downstream of nothing in
\Cref{sec:ch48-potentials} but is used to explain it.
\end{warning}

\begin{proposition}
\label{prop:ch48-uniform-strong}
Let \(\phi_i,\phi\in\cE^1(\omega)\) and suppose
\(\varepsilon_i:=\sup_{X^{\divis}}|\phi_i-\phi|\to0\).  Then
\(\mathbb I(\phi_i,\phi)\le3\varepsilon_i\to0\), so \(\phi_i\to\phi\)
strongly.  In particular the inclusion of
\(\bigl(\cE^1(\omega)\cap C^0(X^{\An}),\lVert\,\cdot\,\rVert_\infty\bigr)\)
into the strong space is continuous.
\end{proposition}

\begin{proof}
This is the last clause of \Cref{prop:ch47-quasimetric}:
\(I(\phi_i,\phi)\le2\sup_{X^{\divis}}|\phi_i-\phi|\) by
\Cref{lem:ch47-IJ-elementary}\,\cref{itm:ch47-IJ:sup}, and
\(|\sup\phi_i-\sup\phi|=|\phi_i(v_{\triv})-\phi(v_{\triv})|
\le\sup_{X^{\divis}}|\phi_i-\phi|\) by
\Cref{lem:ch47-three-normalizations}\,\cref{itm:ch47-norm:max}.  For the last
sentence, note that a continuous function on the compact space
\(X^{\An}\) is bounded, hence lies in \(\cE^1(\omega)\) by
\Cref{thm:bj-e1-basic}, and that
\(\sup_{X^{\divis}}|\phi_i-\phi|\le\lVert\phi_i-\phi\rVert_\infty\).
\end{proof}

\subsection{Completeness}
\label{sec:ch48-completeness}

\Cref{rem:ch47-strong-deferred}\,\cref{itm:ch47-strong:complete} deferred to
this chapter the statement that \(\cE^1(\omega)\) is complete for its strong
uniformity, and \Cref{thm:bj-cy-envelope-equivalence} listed it as one of
the four equivalent forms of the envelope property.  We state it here and
give its proof, separating the part that is a compactness argument --- which
is proved --- from the part that upgrades weak to strong convergence, which
is the import.

\begin{theorem}
\label{thm:ch48-completeness}
Let \(X\) be integral projective unibranch and \(\omega\in\Amp(X)\).  Then
every \(\mathbb I\)-Cauchy net in \(\cE^1(\omega)\) converges strongly to a
point of \(\cE^1(\omega)\).  Equivalently,
\(\bigl(\cE^1(\omega)/\mathbf R,I\bigr)\) is complete.
\end{theorem}

\begin{proof}
\emph{Step 1: reduction to the normalized, \(I\)-Cauchy case.}  Let
\((\phi_i)_{i\in A}\) be \(\mathbb I\)-Cauchy.  Since
\(|\sup\phi_i-\sup\phi_j|\le\mathbb I(\phi_i,\phi_j)\), the net of real
numbers \(c_i:=\sup_{X^{\An}}\phi_i\) is Cauchy, hence converges to some
\(c\in\mathbf R\).  Put \(\tilde\phi_i:=\phi_i-c_i\), so that
\(\sup\tilde\phi_i=0\) and, by
\Cref{lem:ch47-IJ-elementary}\,\cref{itm:ch47-IJ:translation},
\(I(\tilde\phi_i,\tilde\phi_j)=I(\phi_i,\phi_j)\); thus
\((\tilde\phi_i)\) is \(I\)-Cauchy in \(\cE^1_0(\omega)\).  If
\(\tilde\phi_i\to\tilde\phi\) strongly then
\(\phi_i=\tilde\phi_i+c_i\to\tilde\phi+c\) strongly, because
\(\mathbb I(\tilde\phi_i+c_i,\tilde\phi+c)
=I(\tilde\phi_i,\tilde\phi)+|c_i-c|\).  So we may and do assume
\(\sup\phi_i=0\) for every \(i\).

\emph{Step 2: the net lies in a weakly compact set.}  Fix \(i_0\) with
\(I(\phi_i,\phi_j)\le1\) for \(i,j\ge i_0\).  By
\Cref{cor:ch47-J-vanishes} and \eqref{eq:bj-I-quasitriangle},
\[
 J(\phi_i)\le\tfrac n{n+1}I(\phi_i,0)
 \le\tfrac{nC_n}{n+1}\bigl(I(\phi_i,\phi_{i_0})+I(\phi_{i_0},0)\bigr)
 \le\tfrac{nC_n}{n+1}\bigl(1+2J(\phi_{i_0})\tfrac{n+1}{1}\bigr)=:C
\]
for all \(i\ge i_0\), where we used
\(I(\phi_{i_0},0)\le(n+1)J(\phi_{i_0})\) from
\eqref{eq:bj-I-J-comparison}.  Hence for \(i\ge i_0\)
\[
 \phi_i\in\mathcal K_C:=\{\psi\in\PSH^{\sup}(\omega)
                   \mid E_\omega(\psi)\ge-C\}
 =\{\psi\in\cE^1_0(\omega)\mid J(\psi)\le C\},
\]
and \(\mathcal K_C\) is compact for the weak topology: it is the intersection of
the compact space \(\PSH^{\sup}(\omega)\) of
\Cref{cor:ch45-compactness} with the closed set
\(\{E_\omega\ge-C\}\), closed because \(E_\omega\) is weakly upper
semicontinuous
(\Cref{prop:ch46-energy-usc}\,\cref{itm:ch46-Eusc:props}).  Therefore some
subnet \((\phi_{i_\alpha})\) converges weakly to some
\(\phi\in\mathcal K_C\subset\cE^1_0(\omega)\).

\emph{Step 3: the upgrade.}  It remains to prove
\(I(\phi_i,\phi)\to0\) for the full net.  It suffices to prove it along the
subnet: if \(I(\phi_{i_\alpha},\phi)\to0\), then for \(i\) large and
\(\alpha\) large, \(I(\phi_i,\phi)\le
C_n\bigl(I(\phi_i,\phi_{i_\alpha})+I(\phi_{i_\alpha},\phi)\bigr)\) is
small, by the Cauchy property and \eqref{eq:bj-I-quasitriangle}.  For the
subnet the argument is the one of
\cite[Theorem~12.4]{BJ22}: one replaces \((\phi_{i_\alpha})\) by the
increasing net of envelopes
\[
 \chi_\beta:=P_\omega\Bigl(\inf_{\alpha\ge\beta}\phi_{i_\alpha}\Bigr),
\]
which is legitimate because an infimum of upper semicontinuous functions is
upper semicontinuous and bounded above by \(0\), so that
\Cref{imp:ch48-envelope}\,\cref{itm:ch48-env:operator} applies whenever the
competitor family is nonempty --- which it is, by the Cauchy property, once
\(\beta\) is large enough --- and then invokes
\Cref{imp:ch48-imported}\,\cref{itm:ch48-import:increasing}: the net
\((\chi_\beta)\) is increasing with regularized supremum \(\phi\), so
\(E_\omega(\chi_\beta)\to E_\omega(\phi)\), whence
\(I(\chi_\beta,\phi)\to0\) by \Cref{cor:ch47-monotone-I}; and
\(I(\phi_{i_\alpha},\chi_\alpha)\to0\) because
\(\chi_\alpha\le\phi_{i_\alpha}\) and the two energies are squeezed by the
Cauchy property.  One application of \eqref{eq:bj-I-quasitriangle}
concludes.

\emph{Step 4: strong convergence.}  By
\Cref{thm:bj-strong-e1-quasimetric}, \(I(\phi_i,\phi)\to0\) together with
\(\sup\phi_i=0=\sup\phi\) gives \(\mathbb I(\phi_i,\phi)\to0\), that is,
strong convergence.  Undoing Step~1 completes the proof.

The quotient statement follows, a quasi-metric quotient of a complete space
by a group of isometries with closed orbits being complete; concretely, the
orbits are the sets \(\phi+\mathbf R\), and \(I\) is exactly the induced
gauge by \Cref{prop:ch47-quasimetric}.
\end{proof}

\begin{remark}[Which single step carries the analytic weight]
\label{rem:ch48-completeness-weight}
Steps~1, 2 and 4 are complete arguments and use only
\Cref{cor:ch45-compactness}, the upper semicontinuity of \(E_\omega\), and
the elementary calculus of \(I\).  Step~3 is the import.  It is worth being
precise about what it supplies: weak compactness produces a cluster point,
and nothing more; the passage from ``\(\phi\) is a weak cluster point of an
\(I\)-Cauchy net'' to ``\(I(\phi_{i_\alpha},\phi)\to0\)'' is exactly the
statement that the strong uniformity has no missing points, and there is no
soft argument for it.  This is the implication
(1)\(\Rightarrow\)(4) of \Cref{thm:bj-cy-envelope-equivalence}, whose proof
that chapter deferred to this one, and it is now discharged in the only
sense available: stated with its exact source, with the surrounding
argument supplied.

Two consequences of \Cref{thm:ch48-completeness} are used later.  It is
what makes the Calabi--Yau theorem \Cref{cor:bj-cy-unibranch} available for
\emph{every} \(\mu\in\cM^1(\omega)\), not merely for the admissible
measures of \Cref{def:ch47-admissible-measure}, through the implication
(4)\(\Rightarrow\)(2) of \Cref{thm:bj-cy-envelope-equivalence}; and it is
what identifies \(\cE^1(\omega)\) with the \(d_1\)-completion in
\Cref{thm:stability-norm-completion}.
\end{remark}

\subsection{\texorpdfstring{\(\mathbb I\)}{I} is not a metric, and \texorpdfstring{\(d_1\)}{d1} is}
\label{sec:ch48-not-a-metric}

The strong topology is metrizable: \Cref{thm:stability-norm-completion}
produces a metric \(d_1\), coming from the norm geometry of graded
filtrations, that induces it on \(\cE^1(L)\) and on \(\cM^1(L)\).  It is
therefore natural to ask whether \(d_1\) can be taken to be
\(\mathbb I\).  It cannot, and the obstruction is elementary and explicit.

\begin{example}[The triangle inequality fails for \(\mathbb I\), with numbers]
\label{ex:ch48-not-a-metric}
Let \(X=\mathbf P^1_k\) and \(\omega=c_1(\cO(1))\), so \(n=1\), \(V=1\).
Fix \(p\in X(k)\) and keep the family of
\Cref{ex:ch47-p1-numbers} with \(d=1\):
\[
 \begin{aligned}
  \phi_c&=c\max\{\log|\mathfrak m_p|,-1\},&0&\le c\le1,\\
  I(\phi_c,\phi_{c'})&=(c-c')^2,&\sup\phi_c&=0.
 \end{aligned}
\]
Since all three suprema vanish, \(\mathbb I=I\) on this family.  Take
\(\phi_1:=\phi_0=0\), \(\phi_3:=\phi_{1/2}\), \(\phi_2:=\phi_1\).  Then
\[
 \mathbb I(\phi_0,\phi_1)=1,
 \qquad
 \mathbb I(\phi_0,\phi_{1/2})+\mathbb I(\phi_{1/2},\phi_1)
 =\tfrac14+\tfrac14=\tfrac12 .
\]
So \(\mathbb I(\phi_0,\phi_1)=2\bigl(\mathbb I(\phi_0,\phi_{1/2})
+\mathbb I(\phi_{1/2},\phi_1)\bigr)\): the triangle inequality fails by a
factor of exactly \(2\), and the quasi-triangle constant
\(C_1=2\) of \Cref{prop:ch47-curve-I-norm} is attained.  Consequently
\(\mathbb I\) is not a metric, and no metric --- in particular no
\(d_1\) --- is equal to it.

What is true, on this family and in general when \(n=1\), is that the
square root repairs the defect.  Indeed for \(n=1\) the function
\(\rho:=\mathbb I^{1/2}\) is a genuine metric on \(\cE^1(\omega)\).  To see
this write \(f_{ab}:=\phi_a-\phi_b\) and
\(\Delta_{ab}:=\sup\phi_a-\sup\phi_b\), so that
\(\rho(\phi_a,\phi_b)^2=\lVert f_{ab}\rVert^2+|\Delta_{ab}|\) with
\(\lVert\,\cdot\,\rVert\) the seminorm of \Cref{prop:ch47-curve-I-norm}.
Then
\[
 \rho(\phi_1,\phi_2)^2
 \le\bigl(\lVert f_{13}\rVert+\lVert f_{32}\rVert\bigr)^2
   +|\Delta_{13}|+|\Delta_{32}|
 =\rho(\phi_1,\phi_3)^2+\rho(\phi_3,\phi_2)^2
   +2\lVert f_{13}\rVert\,\lVert f_{32}\rVert,
\]
while
\(\bigl(\rho(\phi_1,\phi_3)+\rho(\phi_3,\phi_2)\bigr)^2
=\rho(\phi_1,\phi_3)^2+\rho(\phi_3,\phi_2)^2
+2\rho(\phi_1,\phi_3)\rho(\phi_3,\phi_2)\), and
\(\rho(\phi_1,\phi_3)\rho(\phi_3,\phi_2)\ge
\lVert f_{13}\rVert\,\lVert f_{32}\rVert\) because
\(\rho\ge\lVert\,\cdot\,\rVert\) termwise.  Hence
\(\rho(\phi_1,\phi_2)\le\rho(\phi_1,\phi_3)+\rho(\phi_3,\phi_2)\);
symmetry and separation are \Cref{prop:ch47-quasimetric}.  On the family
above \(\rho(\phi_c,\phi_{c'})=|c-c'|\), so the three numbers are
\(1=\tfrac12+\tfrac12\) and the triangle inequality is an equality: the
family is a geodesic segment for \(\rho\).

For \(n\ge2\) no such statement is claimed.  The exponent that repairs
\eqref{eq:bj-I-quasitriangle} is unknown to us, and
\Cref{thm:bj-uniform-concavity} shows only that \(I\) behaves like the
square of a distance to within dimensional constants.  What is recorded is
what is proved: \(\mathbb I\) is a quasi-metric inducing the strong
topology (\Cref{prop:ch47-quasimetric},
\Cref{thm:bj-strong-e1-quasimetric}), and the strong topology is
metrizable (\Cref{thm:stability-norm-completion}).
\end{example}

\begin{remark}[The comparison with \(d_1\), stated exactly]
\label{rem:ch48-d1-comparison}
\Cref{rem:ch47-strong-deferred}\,\cref{itm:ch47-strong:d1} announced that
the relation between \(\mathbb I\) and \(d_1\) is bi-Lipschitz
comparability and not equality, and that no isometry statement is made
anywhere in this book.  \Cref{ex:ch48-not-a-metric} is the reason.  The
precise situation is the following, and each clause has a reference.

\emph{Same topology.}  \(d_1\) induces the strong topology on
\(\cE^1(L)\), by \Cref{thm:stability-norm-completion}; \(\mathbb I\)
induces the strong topology, by \Cref{thm:bj-strong-e1-quasimetric}.  Hence
a net converges for one gauge if and only if it converges for the other.

\emph{Same Cauchy nets, hence the same completion.}  Both gauges have the
same bounded sets and the same Cauchy nets, because each is bounded by a
fixed increasing function of the other on \(J\)-sublevels; this is what
``bi-Lipschitz comparable'' means here, and it is the content of
\cite[Theorem~B]{BJKS1} together with \Cref{thm:ch48-completeness}.

\emph{Different numbers.}  \(\mathbb I\) fails the triangle inequality
(\Cref{ex:ch48-not-a-metric}) and \(d_1\) does not.  No formula converting
one into the other is asserted, and \Cref{ex:ch48-not-a-metric} shows that
none can be an isometry.  \cite{BJKS1} is where the two are compared, and
\Cref{chap:bj-stability} is where the norm-theoretic side is developed.

The situation over a discretely valued base is the same in shape and
weaker in content: \Cref{sec:ch26-gauge} constructs a gauge
\(d_\omega\) there, \Cref{rem:ch28-not-a-metric-on-potentials} records that
it is not a metric on potentials, and no metrization theorem is available.
\end{remark}

\section{The strong topology on measures}
\label{sec:ch48-measures}

\subsection{The Fenchel defect}
\label{sec:ch48-fenchel-defect}

For \(\mu\in\mathcal M^1\) and \(\phi\in\mathcal E^1\), define the Fenchel
defect

\[
             J_\mu(\phi):=E^\vee(\mu)-E(\phi)+\int\phi\,\rd \mu\geq0.
\tag{48.3}\label{eq:bj-fenchel-defect}
\]

It vanishes precisely at a solution of \(\MA(\phi)=\mu\).

The notation is consistent: when \(\mu=\MA_\omega(\psi)\) for some
\(\psi\in\cE^1(\omega)\), the two meanings of \(J_\mu\) and \(J_\psi\)
agree.  The following proposition proves this and the other properties used
below; all of them are formal consequences of the Legendre duality of
\Cref{sec:ch47-dual}, and none needs an import.

\begin{proposition}
\label{prop:ch48-defect-elementary}
Let \(\mu,\nu\in\cM^1(\omega)\) and \(\phi,\psi\in\cE^1(\omega)\).  Then:
\begin{enumerate}
\item\label{itm:ch48-def:nonneg}
      \(J_\mu(\phi)=E^\vee_\omega(\mu)-F_\mu(\phi)\ge0\), where
      \(F_\mu=E_\omega-\int\cdot\,\rd\mu\), and
      \(\inf_{\phi\in\cE^1(\omega)}J_\mu(\phi)=0\);
\item\label{itm:ch48-def:transl}
      \(J_\mu(\phi+c)=J_\mu(\phi)\) for \(c\in\mathbf R\);
\item\label{itm:ch48-def:ma}
      \(J_{\MA_\omega(\psi)}(\phi)=J_\psi(\phi)\), the right-hand side
      being \eqref{eq:bj-I-J-potentials};
\item\label{itm:ch48-def:zero}
      \(J_\mu(\phi)=0\) if and only if \(\MA_\omega(\phi)=\mu\);
\item\label{itm:ch48-def:convex}
      \(J_\mu\) is convex on \(\cE^1(\omega)\), and
      \(\mu\mapsto J_\mu(\phi)\) is convex on \(\cM^1(\omega)\);
\item\label{itm:ch48-def:mono}
      if \(\phi\le\psi\) then
      \(J_\mu(\psi)\le J_\mu(\phi)+\int(\psi-\phi)\,\rd\mu\).
\end{enumerate}
\end{proposition}

\begin{proof}
\cref{itm:ch48-def:nonneg}.  The identity is
\eqref{eq:bj-fenchel-defect} rewritten, and nonnegativity together with the
vanishing of the infimum is the definition
\eqref{eq:bj-dual-energy} of \(E^\vee_\omega(\mu)\) as the supremum of
\(F_\mu\) over \(\cE^1(\omega)\); the supremum is finite because
\(\mu\in\cM^1(\omega)\).

\cref{itm:ch48-def:transl}.  \(E_\omega(\phi+c)=E_\omega(\phi)+c\) by
\eqref{eq:bj-energy-basic-identities}, and \(\int(\phi+c)\rd\mu
=\int\phi\,\rd\mu+c\) because \(\mu\) has mass \(1\)
(\Cref{lem:ch47-mass-forced}); the two corrections cancel.  This is
\Cref{prop:ch47-dual-elementary}\,\cref{itm:ch47-dual:nonneg}.

\cref{itm:ch48-def:ma}.  By \eqref{eq:bj-dual-energy-ma},
\(E^\vee_\omega(\MA_\omega(\psi))=E_\omega(\psi)-\int\psi\,\MA_\omega(\psi)\).
Substituting into \eqref{eq:bj-fenchel-defect} gives
\[
 J_{\MA_\omega(\psi)}(\phi)
 =E_\omega(\psi)-E_\omega(\phi)
   +\int(\phi-\psi)\,\MA_\omega(\psi)
 =J_\psi(\phi).
\]

\cref{itm:ch48-def:zero}.  \(J_\mu(\phi)=0\) says that \(\phi\) attains the
supremum defining \(E^\vee_\omega(\mu)\), that is, that the constant net
\(\phi\) is maximizing for \(F_\mu\); by
\Cref{cor:ch47-maximizer-solves} this is equivalent to
\(\MA_\omega(\phi)=\mu\).  Conversely, if \(\MA_\omega(\phi)=\mu\) then
\cref{itm:ch48-def:ma} gives \(J_\mu(\phi)=J_\phi(\phi)=0\).

\cref{itm:ch48-def:convex}.  In \(\phi\): \(-E_\omega\) is convex on the
convex set \(\cE^1(\omega)\) by
\Cref{prop:ch46-energy-usc}\,\cref{itm:ch46-Eusc:props}, and
\(\phi\mapsto\int\phi\,\rd\mu\) is affine along segments.  In \(\mu\):
\(E^\vee_\omega\) is convex by
\Cref{prop:ch47-dual-elementary}\,\cref{itm:ch47-dual:convex} and
\(\mu\mapsto\int\phi\,\rd\mu\) is affine.

\cref{itm:ch48-def:mono}.  \(J_\mu(\psi)-J_\mu(\phi)
=E_\omega(\phi)-E_\omega(\psi)+\int(\psi-\phi)\rd\mu
\le\int(\psi-\phi)\rd\mu\), because \(E_\omega\) is increasing
(\Cref{prop:ch46-energy-usc}\,\cref{itm:ch46-Eusc:props}).
\end{proof}

\begin{remark}[Reading the defect]
\label{rem:ch48-defect-reading}
\Cref{prop:ch48-defect-elementary}\,\cref{itm:ch48-def:ma} is the reason the
measure-side theory is a faithful copy of the potential-side one on the
image of \(\MA_\omega\): on that image the new functional
\eqref{eq:bj-fenchel-defect} is literally the old one
\eqref{eq:bj-I-J-potentials}.  Everything that is genuinely new happens off
the image --- and, by \Cref{cor:bj-cy-unibranch}, there is nothing off the
image.  The reader should nevertheless keep the two apart, because the
Calabi--Yau theorem is a theorem: the definitions below are made for all of
\(\cM^1(\omega)\) and their elementary properties are proved without it.
The corresponding discipline over a discretely valued base is
\Cref{rem:ch28-completion-not-surjectivity}, where surjectivity is
\emph{not} available and the distinction is not a matter of taste.
\end{remark}

\subsection{The dual quasi-metric}
\label{sec:ch48-dual-quasimetric}

\begin{definition}[Measure quasi-metric]
\label{def:bj-measure-quasimetric}
For \(\mu,\nu\in\mathcal M^1\), set

\[
 I^\vee(\mu,\nu)
 :=\inf_{\phi\in\mathcal E^1}
       \bigl(J_\mu(\phi)+J_\nu(\phi)\bigr).
\tag{48.4}\label{eq:bj-measure-Idual}
\]
\end{definition}

\begin{proposition}
\label{prop:ch48-Idual-elementary}
Let \(\mu,\nu\in\cM^1(\omega)\).  Then:
\begin{enumerate}
\item\label{itm:ch48-Id:sym}
      \(I^\vee(\mu,\nu)=I^\vee(\nu,\mu)\ge0\) and
      \(I^\vee(\mu,\mu)=0\);
\item\label{itm:ch48-Id:ma}
      \(I^\vee\bigl(\MA_\omega(\phi),\MA_\omega(\psi)\bigr)
       =\inf_{\chi\in\cE^1(\omega)}\bigl(J_\phi(\chi)+J_\psi(\chi)\bigr)\)
      for \(\phi,\psi\in\cE^1(\omega)\);
\item\label{itm:ch48-Id:sep}
      \(I^\vee(\mu,\nu)=0\) implies \(\mu=\nu\);
\item\label{itm:ch48-Id:upper}
      \(I^\vee(\mu,\nu)\le J_\mu(\phi)+J_\nu(\phi)\) for every
      \(\phi\in\cE^1(\omega)\); in particular
      \(I^\vee\bigl(\mu,\MA_\omega(\phi)\bigr)\le J_\mu(\phi)\), and every
      upper bound for \(I^\vee\) proved below is of this shape;
\item\label{itm:ch48-Id:energy}
      if \(\phi\) is almost optimal in \eqref{eq:bj-measure-Idual}, in the
      sense that \(J_\mu(\phi)+J_\nu(\phi)\le I^\vee(\mu,\nu)+\varepsilon\),
      then
      \[
       \bigl|E^\vee_\omega(\mu)-E^\vee_\omega(\nu)\bigr|
       \le I^\vee(\mu,\nu)+\varepsilon
         +\Bigl|\int\phi\,\rd(\mu-\nu)\Bigr| .
      \]
\end{enumerate}
\end{proposition}

\begin{proof}
\cref{itm:ch48-Id:sym}.  Symmetry and nonnegativity are visible in
\eqref{eq:bj-measure-Idual}, using
\Cref{prop:ch48-defect-elementary}\,\cref{itm:ch48-def:nonneg}.  For
\(I^\vee(\mu,\mu)=2\inf_\phi J_\mu(\phi)=0\), use the same item.

\cref{itm:ch48-Id:ma} is
\Cref{prop:ch48-defect-elementary}\,\cref{itm:ch48-def:ma} substituted into
\eqref{eq:bj-measure-Idual}.

\cref{itm:ch48-Id:sep}.  Choose \(\chi_i\in\cE^1(\omega)\) with
\(J_\mu(\chi_i)+J_\nu(\chi_i)\to0\).  Then \((\chi_i)\) is a maximizing net
for \(F_\mu\) in the sense of \Cref{def:ch47-maximizing-net}, and also for
\(F_\nu\).  By \Cref{thm:bj-maximizers-converge}, in the form
\eqref{eq:bj-maximizer-ma-weak}, \(\MA_\omega(\chi_i)\to\mu\) weakly and
\(\MA_\omega(\chi_i)\to\nu\) weakly.  The weak topology on \(\mathcal M\)
is Hausdorff, so \(\mu=\nu\).

\cref{itm:ch48-Id:upper} is the definition
\eqref{eq:bj-measure-Idual} as an infimum, together with
\(J_{\MA_\omega(\phi)}(\phi)=J_\phi(\phi)=0\) from
\Cref{prop:ch48-defect-elementary}\,\cref{itm:ch48-def:ma}.

\cref{itm:ch48-Id:energy}.  Let \(\phi\) be almost optimal.  Then
\(E^\vee_\omega(\mu)=J_\mu(\phi)+E_\omega(\phi)-\int\phi\,\rd\mu\) and
similarly for \(\nu\); subtracting gives
\(E^\vee_\omega(\mu)-E^\vee_\omega(\nu)
=J_\mu(\phi)-J_\nu(\phi)-\int\phi\,\rd(\mu-\nu)\), and
\(|J_\mu(\phi)-J_\nu(\phi)|\le J_\mu(\phi)+J_\nu(\phi)
\le I^\vee(\mu,\nu)+\varepsilon\).
\end{proof}

\begin{remark}[Convexity of \(I^\vee\) is not claimed]
\label{rem:ch48-Idual-not-convex}
\Cref{prop:ch48-defect-elementary}\,\cref{itm:ch48-def:convex} makes
\(J_\mu(\phi)\) convex in each variable separately, but the pairing
\(\int\phi\,\rd\mu\) is bilinear and not jointly convex, so the infimum
\eqref{eq:bj-measure-Idual} of a family of separately convex functions
carries no convexity of its own.  We therefore assert none, and no argument
below uses one.  This should be contrasted with \(E^\vee_\omega\) itself,
which \emph{is} convex
(\Cref{prop:ch47-dual-elementary}\,\cref{itm:ch47-dual:convex}) because it
is a supremum of affine functions of \(\mu\); the difference is exactly
that in \eqref{eq:bj-dual-energy} the potential is the index of the family
and in \eqref{eq:bj-measure-Idual} it is the variable of an infimum.
\end{remark}

\subsection{Strong convergence of measures}
\label{sec:ch48-strong-measures}

\begin{theorem}\label{thm:bj-strong-measures}
The functional \(I^\vee\) is a quasi-metric.  For a net
\((\mu_i)\) in \(\mathcal M^1\) and \(\mu\in\mathcal M^1\), the following
are equivalent:

\begin{enumerate}
\item \(\mu_i\to\mu\) weakly and
      \(E^\vee(\mu_i)\to E^\vee(\mu)\);
\item \(I^\vee(\mu_i,\mu)\to0\);
\item for every \(C>0\),
\[
  \int\phi\,\rd \mu_i\longrightarrow\int\phi\,\rd \mu
\]
uniformly for \(\phi\in\mathcal E^1\) with \(J(\phi)\leq C\).
\end{enumerate}

This topology is complete.
\end{theorem}

\begin{proof}[Proof architecture]
Approximate \(\mu\) and \(\nu\) by Monge--Amp\`ere measures of maximizing
nets.  On the image of Monge--Amp\`ere, the potential and measure
quasi-metrics are comparable.  The mixed-energy H\"older estimate controls

\[
 \left|\int(\phi-\psi)\,\rd(\mu-\nu)\right|
\]

by a positive power of \(I(\phi,\psi)\), a positive power of
\(I^\vee(\mu,\nu)\), and bounded energy factors.  This gives separation,
the quasi-triangle inequality, and the equivalence of the three convergence
criteria.  A Cauchy net has a weakly convergent subnet by compactness of
the probability measures; lower semicontinuity keeps its limit in
\(\mathcal M^1\), and uniform convergence on \(J\)-sublevels upgrades the
subnet, hence the whole net, to strong convergence.  These are
\cite[Theorems~10.2--10.3, 10.12, and 10.14]{BJ22}.

We now separate what is proved here from what is imported.  Four of the
constituent statements are proved above or below; two are not.

\emph{Proved: \(I^\vee\) is a symmetric separating gauge.}  This is
\Cref{prop:ch48-Idual-elementary}\,\cref{itm:ch48-Id:sym},
\cref{itm:ch48-Id:sep}, whose proof used only
\Cref{thm:bj-maximizers-converge} from \Cref{chap:bj-finite-energy}.

\emph{Imported: the quasi-triangle inequality for \(I^\vee\).}  Given
\(\lambda\) and near-optimal potentials \(\chi_1\) for \((\mu,\lambda)\)
and \(\chi_2\) for \((\lambda,\nu)\), one has to compare \(J_\nu(\chi_1)\)
with \(J_\nu(\chi_2)\), and
\[
 J_\nu(\chi_1)-J_\nu(\chi_2)
 =E_\omega(\chi_2)-E_\omega(\chi_1)+\int(\chi_1-\chi_2)\,\rd\nu .
\]
Bounding the right-hand side by a positive power of \(I(\chi_1,\chi_2)\)
with constants depending on \(E^\vee_\omega(\nu)\) is exactly
\Cref{imp:ch48-imported}\,\cref{itm:ch48-import:pairing}, that is,
\cite[Theorem~10.3]{BJ22}.  There is no elementary substitute: the
difference of two Fenchel defects at different potentials is a pairing of a
potential difference against a measure, and the whole difficulty of the
measure-side theory is that such a pairing is not controlled by the
potential gauge alone.

\emph{Proved: \((3)\Rightarrow\) weak convergence.}  Assume (3).  Every
\(f\in\PL_{\mathbf R}(X)\) is a difference \(f=\phi-\psi\) of two elements
of \(\cH_{\mathrm{dom}}(\omega)\), by \Cref{def:ch34-hdom} and the remark
following it that \(\cH_{\mathrm{dom}}(\omega)\) spans \(\PL(X)\); such
\(\phi,\psi\) are bounded, hence lie in \(\cE^1(\omega)\) by
\Cref{thm:bj-e1-basic}, and \(J(\phi),J(\psi)\le C\) for some \(C\)
depending on \(f\).  Applying (3) to \(\phi\) and to \(\psi\) separately
gives \(\int f\,\rd\mu_i\to\int f\,\rd\mu\).  Since
\(\PL_{\mathbf R}(X)\) is uniformly dense in \(C^0(X^{\An})\)
(\Cref{thm:ch34-pl-density}) and all the measures have mass \(1\), the
convergence extends to every continuous test function by a
three-\(\varepsilon\) argument.  Hence \(\mu_i\to\mu\) weakly.

\emph{Proved: \((3)\Rightarrow\varliminf_iE^\vee_\omega(\mu_i)\ge
E^\vee_\omega(\mu)\).}  Fix \(C>0\) and
\(\phi\in\cE^1(\omega)\) with \(J(\phi)\le C\).  By (3),
\(F_{\mu_i}(\phi)=E_\omega(\phi)-\int\phi\,\rd\mu_i\to F_\mu(\phi)\), and
\(E^\vee_\omega(\mu_i)\ge F_{\mu_i}(\phi)\), so
\(\varliminf_iE^\vee_\omega(\mu_i)\ge F_\mu(\phi)\).  Taking the supremum
over all \(\phi\) with \(J(\phi)\le C\) and then over \(C\), and using
\Cref{prop:ch48-defect-elementary}\,\cref{itm:ch48-def:transl} to normalize
\(\sup\phi=0\), gives the claim.  Note that this half is also weak lower
semicontinuity of \(E^\vee_\omega\)
(\Cref{prop:ch47-dual-elementary}\,\cref{itm:ch47-dual:lsc}); the two
arguments are the same one.

\emph{Imported: the reverse inequality and the remaining implications.}
The equivalences \((1)\Leftrightarrow(2)\Leftrightarrow(3)\) in full are
\cite[Theorems~10.12 and~10.14]{BJ22}, and each of the missing steps is an
application of the pairing estimate.

\emph{Proved: the compactness half of completeness.}  Let \((\mu_i)\) be
\(I^\vee\)-Cauchy.  Then \(E^\vee_\omega(\mu_i)\) is bounded: indeed, fixing
\(i_0\) with \(I^\vee(\mu_i,\mu_{i_0})\le1\) for \(i\ge i_0\) and choosing
almost optimal potentials, the estimate
\Cref{prop:ch48-Idual-elementary}\,\cref{itm:ch48-Id:energy} bounds
\(E^\vee_\omega(\mu_i)\) once the pairing term is bounded, which it is by
\Cref{prop:ch47-test-class}.  The sublevel set \(\cM^1_C\) of
\eqref{eq:bj-measure-energy-sublevel} is weakly compact by
\Cref{prop:ch47-dual-elementary}\,\cref{itm:ch47-dual:compact}, so some
subnet converges weakly to a measure \(\mu\), and weak lower
semicontinuity of \(E^\vee_\omega\) puts \(\mu\) in \(\cM^1_C\subset
\cM^1(\omega)\).

\emph{Imported: the upgrade.}  That \(I^\vee(\mu_i,\mu)\to0\) --- first
along the subnet, then, by the quasi-triangle inequality and the Cauchy
property, along the whole net --- is \cite[Theorem~10.14]{BJ22}.
\end{proof}

\begin{remark}[Weak compactness is not strong compactness, again]
\label{rem:ch48-measure-weak-not-strong}
\Cref{prop:ch47-dual-elementary}\,\cref{itm:ch47-dual:compact} says that
\(\cM^1_C\) is weakly compact, and
\Cref{rem:ch47-weak-not-strong} warned that this is not a strong statement.
\Cref{ex:ch48-weak-not-strong} below makes the warning quantitative: it
exhibits a sequence in \(\cM^1_{1/2}\) on \(\mathbf P^1\) that converges
weakly and satisfies \(I^\vee(\mu_j,\mu)=\tfrac14\) for every \(j\), so that
no subnet converges strongly and \(\cM^1_C\) is not strongly compact for
\(C\ge\tfrac12\).  The same example, transported through
\Cref{prop:bj-ma-strong-embedding}, shows that the potential sublevel
\(\{J\le C\}\) is not strongly compact; this is
\Cref{prop:ch48-no-strong-compactness}.
\end{remark}

\begin{remark}[The measure side does not see the polarization]
\label{rem:ch48-intrinsic}
The definition is intrinsic when \(X\) is irreducible: changing the ample
class changes neither \(\mathcal M^1\) nor its strong topology
\cite[Corollary~9.21 and Proposition~10.16]{BJ22}.  The numerical energies
and quasi-metrics do change, so no isometry is claimed.

Concretely, if \(\omega,\omega'\in\Amp(X)\) then there is \(t>0\) with
\(t\omega\le\omega'\le t^{-1}\omega\), and \Cref{thm:bj-e1-basic} gives
\(\cE^1(\omega)=\cE^1(\omega')\) up to the change of class; on the measure
side the two dual energies are comparable, with constants depending on
\(t\) and \(n\), and the sublevel sets
\eqref{eq:bj-measure-energy-sublevel} for the two classes are cofinal in
each other.  Hence \(\cM^1(\omega)=\cM^1(\omega')\) as sets and the two
strong topologies coincide.  What does \emph{not} survive is any numerical
statement: \(E^\vee_\omega\ne E^\vee_{\omega'}\), and
\(I^\vee_\omega\ne I^\vee_{\omega'}\), as
\Cref{ex:ch47-p1-numbers} already shows by varying \(d\).  This is one more
instance of the rule stated in \Cref{sec:ch47-dictionary}: the class of
finite-energy objects is intrinsic, the numbers attached to them are not.
\end{remark}

\section{Monge--Amp\`ere between the two strong topologies}
\label{sec:ch48-ma}

This section proves the main theorem of the chapter.  It is the point at
which the constructions of \Cref{chap:bj-energy,chap:bj-finite-energy} pay
for themselves: the operator that was defined by intersection numbers on
test configurations, extended by monotone approximation, and shown to be
surjective by a variational argument, turns out to be an isomorphism of
metric structures.

\subsection{The bi-Lipschitz embedding}
\label{sec:ch48-bilipschitz}

\begin{proposition}\label{prop:bj-ma-strong-embedding}
The Monge--Amp\`ere map induces a bi-Lipschitz embedding

\[
 (\mathcal E^1/\mathbf R,I)\hookrightarrow
             (\mathcal M^1,I^\vee).
\tag{48.5}\label{eq:bj-ma-bilipschitz}
\]

Its image is strongly dense.  If the envelope property holds, it is onto
and hence a homeomorphism of the two strong spaces.
\end{proposition}

\begin{proof}
\emph{The two-sided estimate.}  Let \(\phi,\psi\in\cE^1(\omega)\).  By
\Cref{prop:ch48-Idual-elementary}\,\cref{itm:ch48-Id:ma},
\[
 I^\vee\bigl(\MA_\omega(\phi),\MA_\omega(\psi)\bigr)
 =\inf_{\chi\in\cE^1(\omega)}\bigl(J_\phi(\chi)+J_\psi(\chi)\bigr).
\]
Taking \(\chi=\phi\) and using \(J_\phi(\phi)=0\) together with
\eqref{eq:bj-I-J-comparison} gives the upper bound
\[
 I^\vee\bigl(\MA_\omega(\phi),\MA_\omega(\psi)\bigr)
 \le J_\psi(\phi)\le\frac n{n+1}\,I(\phi,\psi).
\]
For the lower bound, apply \eqref{eq:bj-I-J-comparison} to each summand and
then \eqref{eq:bj-I-quasitriangle}: for every \(\chi\),
\[
 J_\phi(\chi)+J_\psi(\chi)
 \ \ge\ \frac1{n+1}\bigl(I(\chi,\phi)+I(\chi,\psi)\bigr)
 \ \ge\ \frac1{(n+1)C_n}\,I(\phi,\psi).
\]
Taking the infimum over \(\chi\),
\[
 \frac{1}{(n+1)C_n}\,I(\phi,\psi)
 \ \le\ I^\vee\bigl(\MA_\omega(\phi),\MA_\omega(\psi)\bigr)
 \ \le\ \frac n{n+1}\,I(\phi,\psi).
\]
This is the bi-Lipschitz statement \eqref{eq:bj-ma-bilipschitz}; that the
map is well defined on \(\cE^1(\omega)/\mathbf R\) is
\(\MA_\omega(\phi+c)=\MA_\omega(\phi)\), and that it is injective there is
\Cref{thm:ch47-uniqueness}.

\emph{Strong density of the image.}  Let \(\mu\in\cM^1(\omega)\) and let
\((\phi_i)\) be a maximizing net for \(F_\mu\), which exists because
\(E^\vee_\omega(\mu)<\infty\); thus \(J_\mu(\phi_i)\to0\).  By
\Cref{prop:ch48-Idual-elementary}\,\cref{itm:ch48-Id:upper},
\[
 I^\vee\bigl(\mu,\MA_\omega(\phi_i)\bigr)\le J_\mu(\phi_i)
 \longrightarrow0 ,
\]
so \(\mu\) lies in the \(I^\vee\)-closure of the image.  This argument uses
no import at all; in particular it does \emph{not} use the Calabi--Yau
theorem, and it is the statement that survives when the envelope property
is dropped.  Compare \cite[Corollary~10.13]{BJ22}.

\emph{Surjectivity.}  Under the envelope property,
\Cref{thm:bj-cy-envelope-equivalence} gives that
\(\MA_\omega:\cE^1(\omega)\to\cM^1(\omega)\) is onto; for unibranch \(X\)
this is \Cref{cor:bj-cy-unibranch}.  A bijective bi-Lipschitz map of
quasi-metric spaces is a homeomorphism, and by
\Cref{thm:bj-strong-e1-quasimetric} and \Cref{thm:bj-strong-measures} the
two quasi-metrics induce the two strong topologies.
\end{proof}

\begin{proposition}
\label{prop:ch48-curve-Idual}
Let \(\dim X=1\).  Then for all \(\phi,\psi\in\cE^1(\omega)\)
\[
 I^\vee\bigl(\MA_\omega(\phi),\MA_\omega(\psi)\bigr)
 =\tfrac14\,I(\phi,\psi),
\]
the infimum in \eqref{eq:bj-measure-Idual} being attained at the midpoint
\(\chi=\tfrac12(\phi+\psi)\).  In particular both constants in the proof of
\Cref{prop:bj-ma-strong-embedding} are attained when \(n=1\):
\(\frac1{(n+1)C_n}=\frac14\) with \(C_1=2\), while
\(\frac n{n+1}=\frac12\) is not sharp.
\end{proposition}

\begin{proof}
By \Cref{prop:ch47-curve-I-norm}, \(J_\phi(\chi)=\tfrac12I(\chi,\phi)
=\tfrac12\lVert\chi-\phi\rVert^2\), and likewise for \(\psi\).  The
parallelogram identity in the pre-Hilbert space
\(\vec\cE^{\,1}(X)/\mathbf R\) gives
\[
 \lVert\chi-\phi\rVert^2+\lVert\chi-\psi\rVert^2
 =2\Bigl\lVert\chi-\tfrac{\phi+\psi}2\Bigr\rVert^2
   +\tfrac12\lVert\phi-\psi\rVert^2 ,
\]
so \(J_\phi(\chi)+J_\psi(\chi)
=\lVert\chi-\tfrac{\phi+\psi}2\rVert^2+\tfrac14\lVert\phi-\psi\rVert^2\),
which is minimized exactly at \(\chi=\tfrac12(\phi+\psi)\), a member of
\(\cE^1(\omega)\) by convexity (\Cref{thm:bj-e1-basic}).  The minimum is
\(\tfrac14\lVert\phi-\psi\rVert^2=\tfrac14I(\phi,\psi)\).
\end{proof}

\subsection{The main theorem}
\label{sec:ch48-main-theorem}

\begin{theorem}
\label{thm:ch48-ma-homeomorphism}
Let \(X\) be an integral projective unibranch variety of dimension \(n\)
over an algebraically closed trivially valued field, and let
\(\omega\in\Amp(X)\).  Then:
\begin{enumerate}
\item\label{itm:ch48-ma:lip}
      \(\MA_\omega:\cE^1(\omega)\to\cM^1(\omega)\) is Lipschitz from the
      strong topology of potentials to that of measures,
      \[
       I^\vee\bigl(\MA_\omega(\phi),\MA_\omega(\psi)\bigr)
       \le\frac n{n+1}\,\mathbb I(\phi,\psi),
      \]
      and in particular strongly continuous;
\item\label{itm:ch48-ma:split}
      the map \(\phi\mapsto(\phi-\sup\phi,\ \sup\phi)\) is a homeomorphism
      of \(\cE^1(\omega)\), with its strong topology, onto
      \(\cE^1_0(\omega)\times\mathbf R\), where \(\cE^1_0(\omega)\) carries
      the topology of \(I\); consequently the quotient topology on
      \(\cE^1(\omega)/\mathbf R\) is the topology of \(I\);
\item\label{itm:ch48-ma:homeo}
      \(\MA_\omega\) induces a bi-Lipschitz bijection
      \[
       \MA_\omega:\ \bigl(\cE^1(\omega)/\mathbf R,\ I\bigr)
       \ \xrightarrow{\ \sim\ }\ \bigl(\cM^1(\omega),\ I^\vee\bigr),
      \]
      hence a homeomorphism for the two strong topologies, and its
      restriction to \(\cE^1_0(\omega)\) is a strong homeomorphism onto
      \(\cM^1(\omega)\);
\item\label{itm:ch48-ma:open}
      \(\MA_\omega:\cE^1(\omega)\to\cM^1(\omega)\) is a strongly open,
      strongly continuous surjection whose fibres are the orbits
      \(\phi+\mathbf R\);
\item\label{itm:ch48-ma:energy}
      \(E^\vee_\omega(\MA_\omega(\phi))=J_\phi(0)\), and
      \[
       \frac1n\,J(\phi)\ \le\ E^\vee_\omega\bigl(\MA_\omega(\phi)\bigr)
       \ \le\ n\,J(\phi),
      \]
      with equality throughout when \(n=1\).
\end{enumerate}
\end{theorem}

\begin{proof}
\cref{itm:ch48-ma:lip}.  The displayed inequality is the upper bound proved
in \Cref{prop:bj-ma-strong-embedding} together with
\(I\le\mathbb I\).  Continuity for the two strong topologies then follows
from \Cref{thm:bj-strong-e1-quasimetric} and
\Cref{thm:bj-strong-measures}.  Observe that this half of the theorem uses
no import beyond those of \Cref{imp:ch47-imported}: the estimate
\eqref{eq:bj-I-J-comparison} is proved in \Cref{chap:bj-finite-energy}.

\cref{itm:ch48-ma:split}.  The map is a bijection with inverse
\((\phi_0,c)\mapsto\phi_0+c\), by
\eqref{eq:ch47-normalized-classes}.  Both it and its inverse are
\(1\)-Lipschitz for the gauges \(\mathbb I\) on the source and
\(I+|\cdot|\) on the target, because
\(\mathbb I(\phi,\psi)=I(\phi-\sup\phi,\psi-\sup\psi)
+|\sup\phi-\sup\psi|\) by
\Cref{lem:ch47-IJ-elementary}\,\cref{itm:ch47-IJ:translation}.  For the last
clause, the quotient map \(\cE^1(\omega)\to\cE^1(\omega)/\mathbf R\)
becomes the first projection under this identification, and the quotient
topology of a product by a factor is the topology of the other factor.

\cref{itm:ch48-ma:homeo} is \Cref{prop:bj-ma-strong-embedding}, whose
surjectivity clause applies because \(X\) is unibranch
(\Cref{cor:bj-cy-unibranch}).  The restriction statement follows from
\cref{itm:ch48-ma:split}, since \(\cE^1_0(\omega)\) meets every orbit
exactly once.

\cref{itm:ch48-ma:open}.  Write \(\MA_\omega=\iota\circ q\) with
\(q:\cE^1(\omega)\to\cE^1(\omega)/\mathbf R\) the quotient map and
\(\iota\) the homeomorphism of \cref{itm:ch48-ma:homeo}.  A quotient map by
a group of homeomorphisms is open, and a homeomorphism is open; so
\(\MA_\omega\) is open.  The description of the fibres is
\Cref{thm:ch47-uniqueness}.

\cref{itm:ch48-ma:energy}.  By \eqref{eq:bj-dual-energy-ma},
\(E^\vee_\omega(\MA_\omega(\phi))=E_\omega(\phi)
-\int\phi\,\MA_\omega(\phi)\); and by
\eqref{eq:bj-I-J-potentials} with \(\psi=\phi\) and argument \(0\),
\[
 J_\phi(0)=E_\omega(\phi)-E_\omega(0)+\int(0-\phi)\,\MA_\omega(\phi)
 =E_\omega(\phi)-\int\phi\,\MA_\omega(\phi),
\]
using \(E_\omega(0)=0\) from \eqref{eq:bj-trivial-energy-zero}.  The two
right-hand sides agree, which is the asserted identity.  For the
estimates, \eqref{eq:bj-I-J-comparison} applied to the pair
\((0,\phi)\) gives
\[
 \tfrac1{n+1}I(\phi,0)\le J_\phi(0)\le\tfrac n{n+1}I(\phi,0),
 \qquad
 \tfrac1{n+1}I(\phi,0)\le J_0(\phi)=J(\phi)\le\tfrac n{n+1}I(\phi,0),
\]
the second chain using
\Cref{lem:ch47-three-normalizations}\,\cref{itm:ch47-norm:J}.  Eliminating
\(I(\phi,0)\) between the two chains gives
\(J_\phi(0)\ge\frac1{n+1}I\ge\frac1{n+1}\cdot\frac{n+1}nJ(\phi)
=\frac1nJ(\phi)\) and
\(J_\phi(0)\le\frac n{n+1}I\le\frac n{n+1}(n+1)J(\phi)=nJ(\phi)\).  When
\(n=1\) both chains are equalities with constant \(\tfrac12\), so
\(J_\phi(0)=J_0(\phi)=J(\phi)\).
\end{proof}

\begin{remark}[What the theorem does and does not say about weak topologies]
\label{rem:ch48-not-weak}
\Cref{thm:ch48-ma-homeomorphism} is a statement about the \emph{strong}
topologies, and no analogue of it holds weakly.  Two cautions.

First, the weak topologies do not determine each other through
\(\MA_\omega\).  \Cref{ex:ch48-weak-not-strong} exhibits a sequence
\((\phi_j)\) on \(\mathbf P^1\) with \(\phi_j\to0\) weakly and
\(\MA_\omega(\phi_j)\to\MA_\omega(0)\) weakly, while
\(\mathbb I(\phi_j,0)=1\) and
\(I^\vee(\MA_\omega(\phi_j),\MA_\omega(0))=\tfrac14\) for every \(j\).  So
weak convergence on both sides is compatible with a uniform positive
distance on both sides, and the bijection of
\cref{itm:ch48-ma:homeo} is not a homeomorphism for the weak topologies.

Second, \(\cM^1(\omega)\) is not weakly closed in \(\mathcal M\) --- its
sublevel sets are, by
\Cref{prop:ch47-dual-elementary}\,\cref{itm:ch47-dual:compact}, but their
union is not --- and \(\cE^1(\omega)\) is not weakly closed in
\(\PSH(\omega)\).  \Cref{ex:ch46-homogeneous-mass-escapes} is the standard
witness on the potential side: the truncations
\(\psi^{(t)}=\max\{\log|\mathfrak m_p|,-t\}\) lie in
\(\cE^1(\omega)\), converge weakly to \(\log|\mathfrak m_p|\), and their
energies tend to \(-\infty\).  On the measure side the corresponding
Monge--Amp\`ere measures \(\delta_{t\ord_p}\) converge weakly to
\(\delta_{v_{p,\triv}}\), which has \(E^\vee_\omega=+\infty\) because
\(v_{p,\triv}\notin X^{\lin}\) (\Cref{thm:bj-dirac-finite-energy}).  The
strong topologies see all of this: in both cases the relevant energy
diverges, and the nets are not strongly Cauchy.
\end{remark}

\section{Capacity}
\label{sec:ch48-capacity}

\subsection{Definition and elementary properties}
\label{sec:ch48-capacity-def}

\begin{definition}[Capacity]
\label{def:bj-bt-capacity}
For a Borel set \(A\subset X^{\An}\), set

\[
 \Capa_\omega(A):=
 \sup\left\{\int_A\MA_\omega(\psi)\ \middle|\
       \psi\in\PSH(\omega),\ -1\leq\psi\leq0\right\}.
\tag{48.6}\label{eq:bj-bt-capacity}
\]

For arbitrary \(A\subset X^{\An}\), define

\[
 \Capa_\omega^*(A):=\inf_{U\supset A,\ U\text{ open}}\Capa_\omega(U).
\tag{48.7}\label{eq:bj-outer-capacity}
\]
\end{definition}

We write \(\mathcal U_\omega:=\{\psi\in\PSH(\omega)\mid-1\le\psi\le0\}\) for the
\emph{admissible} class occurring in \eqref{eq:bj-bt-capacity}; it is
nonempty, since the constants \(0\) and \(-1\) are \(\omega\)-psh
(\Cref{lem:ch34-psh-nonempty}, \(\omega\) being ample hence nef), and it is
convex and stable under finite maxima by
\Cref{thm:bj-basic-closure}.  This is the probability normalization:
\(0\leq\Capa(A)\leq1\).  The
capacity is monotone, countably subadditive, inner regular on Borel sets,
and outer regular on compact sets.  The outer capacity is a Choquet
capacity once the envelope property is assumed
\cite[Propositions~13.3--13.4, Corollary~13.13, and Theorem~13.29]{BJ22}.

\begin{proposition}
\label{prop:ch48-capacity-elementary}
Let \(A,B,A_1,A_2,\dots\subset X^{\An}\) be Borel.  Then:
\begin{enumerate}
\item\label{itm:ch48-cap:mono}
      \(\Capa_\omega(\varnothing)=0\), \(A\subset B\) implies
      \(\Capa_\omega(A)\le\Capa_\omega(B)\), and
      \(\Capa_\omega(X^{\An})=1\);
\item\label{itm:ch48-cap:sub}
      \(\Capa_\omega\bigl(\bigcup_jA_j\bigr)\le\sum_j\Capa_\omega(A_j)\),
      and the same for \(\Capa^*_\omega\) on arbitrary subsets;
\item\label{itm:ch48-cap:inner}
      \(\Capa_\omega(A)=\sup\{\Capa_\omega(K)\mid K\subset A\ \text{compact}\}\),
      and \(\Capa_\omega(U)=\sup\{\Capa_\omega(K)\mid K\subset U\
      \text{compact}\}\) for \(U\) open;
\item\label{itm:ch48-cap:outer}
      \(\Capa^*_\omega\) is monotone, agrees with \(\Capa_\omega\) on open
      sets, and \(\Capa_\omega\le\Capa^*_\omega\) on Borel sets;
\item\label{itm:ch48-cap:triv}
      \(\Capa_\omega(A)=1\) for every \(A\) containing \(v_{\triv}\); in
      particular \(\Capa_\omega(\{v_{\triv}\})=1\);
\item\label{itm:ch48-cap:scaling}
      \(t\mapsto\Capa_\omega(t_\star A)\) is nonincreasing on
      \(\mathbf R_{>0}\), where \(t_\star\) is the scaling action
      \eqref{eq:valuation-scaling}; the capacity is \emph{not} invariant
      under the scaling action.
\end{enumerate}
\end{proposition}

\begin{proof}
\cref{itm:ch48-cap:mono}.  Each \(\MA_\omega(\psi)\) is a positive measure,
so \(A\mapsto\MA_\omega(\psi)(A)\) is monotone and vanishes on
\(\varnothing\); a supremum of such is again monotone and vanishes on
\(\varnothing\).  For the total mass, \(\MA_\omega(\psi)\) is a probability
measure by \Cref{warn:ch48-normalization}, so
\(\Capa_\omega(X^{\An})=1\).

\cref{itm:ch48-cap:sub}.  For fixed \(\psi\in\mathcal U_\omega\),
\(\MA_\omega(\psi)(\bigcup_jA_j)\le\sum_j\MA_\omega(\psi)(A_j)
\le\sum_j\Capa_\omega(A_j)\); take the supremum over \(\psi\).  For
\(\Capa^*_\omega\), given \(\varepsilon>0\) choose open
\(U_j\supset A_j\) with \(\Capa_\omega(U_j)\le\Capa^*_\omega(A_j)
+2^{-j}\varepsilon\); then \(U:=\bigcup_jU_j\) is open, contains
\(\bigcup_jA_j\), and the Borel case applies.

\cref{itm:ch48-cap:inner}.  We first check that a finite Radon measure
\(\mu\) on the compact Hausdorff space \(X^{\An}\) is inner regular on
Borel sets, which is not part of the definition recalled in
\eqref{eq:radon-regularity}.  Given a Borel \(B\) and \(\varepsilon>0\),
outer regularity applied to \(B^c\) gives an open \(U\supset B^c\) with
\(\mu(U)\le\mu(B^c)+\varepsilon\); then \(K:=U^c\) is closed in a compact
space, hence compact, \(K\subset B\), and
\(\mu(B\setminus K)=\mu(B\cap U)=\mu(U)-\mu(U\cap B^c)\le
\mu(U)-\mu(B^c)\le\varepsilon\).  Hence
\(\MA_\omega(\psi)(A)=\sup_{K\subset A}\MA_\omega(\psi)(K)\le
\sup_{K\subset A}\Capa_\omega(K)\); take the supremum over \(\psi\) and use
monotonicity for the reverse inequality.

\cref{itm:ch48-cap:outer}.  Monotonicity is clear from
\eqref{eq:bj-outer-capacity}; agreement on open sets is the case
\(U=A\) of the infimum; and \(\Capa_\omega(A)\le\Capa_\omega(U)\) for every
open \(U\supset A\) gives the last inequality.

\cref{itm:ch48-cap:triv}.  Take \(\psi=0\in\mathcal U_\omega\).  By
\Cref{lem:ch47-three-normalizations}\,\cref{itm:ch47-norm:trivial-measure},
\(\MA_\omega(0)=\delta_{v_{\triv}}\), so
\(\MA_\omega(0)(A)=1\) whenever \(v_{\triv}\in A\).

\cref{itm:ch48-cap:scaling}.  For \(\psi\in\PSH(\omega)\) the scaled
function \((t\cdot\psi)(v)=t\,\psi(t^{-1}v)\) is again \(\omega\)-psh, and
\(\MA_\omega(t\cdot\psi)=t_\star\MA_\omega(\psi)\) by
\eqref{eq:bj-homogeneous-ma-scaling}.  If \(-1\le\psi\le0\) then
\(-t\le t\cdot\psi\le0\), so \(t\cdot\psi\in\mathcal U_\omega\) precisely when
\(t\le1\).  Hence for \(0<t\le1\) and every Borel \(A\),
\[
 \Capa_\omega(t_\star A)\ \ge\ \sup_{\psi\in\mathcal U_\omega}
   \MA_\omega(t\cdot\psi)(t_\star A)
 =\sup_{\psi\in\mathcal U_\omega}\MA_\omega(\psi)(A)=\Capa_\omega(A).
\]
Applying this with \(A\) replaced by \(s_\star A\) and \(t=r/s\le1\) gives
\(\Capa_\omega(r_\star A)\ge\Capa_\omega(s_\star A)\) whenever
\(0<r\le s\), which is the asserted monotonicity.  It is genuinely not an
equality: \Cref{ex:ch48-p1-capacity} computes
\(\Capa_\omega(\{t\ord_p\})=\min\{1,t^{-1}\}\) on \(\mathbf P^1\) with
\(\omega=c_1(\cO(1))\), which is strictly decreasing for \(t\ge1\).  The
reason is that the constraint \(-1\le\psi\le0\) in
\eqref{eq:bj-bt-capacity} is not homogeneous, and it is the same reason
that \(\Capa_\omega\) is a probability rather than a cone-invariant
quantity.
\end{proof}

\begin{remark}[The cone point has full capacity]
\label{rem:ch48-cone-point}
Item \cref{itm:ch48-cap:triv} has no analogue over a discretely valued base
and is worth pausing on.  There, \(\MA_\theta\) of the trivial model metric
is a sum of atoms with weights \(b_E(\cL^n\cdot E)\) spread over the
components of the special fibre, and no single point carries all the mass;
\Cref{ex:ch25-two-vertex-full-capacity} exhibits a divisorial point of full
capacity only after a computation, and
\Cref{ex:ch25-segment-capacities} shows that the value is not uniform.
Here the cone point \(v_{\triv}\) always has capacity \(1\), by the single
identity \(\MA_\omega(0)=\delta_{v_{\triv}}\).  The reader should conclude
that \(\Capa_\omega\) measures nothing topological: it measures how much
Monge--Amp\`ere mass an admissible potential can push into a set, and near
the cone point that is all of it.  The interesting statements are therefore
about sets that stay away from \(v_{\triv}\), which is exactly what the
sublevel sets \(\{\phi\le-t\}\) of a normalized \(\phi\) do.
\end{remark}

\subsection{The Chern--Levine--Nirenberg estimate}
\label{sec:ch48-cln}

\begin{lemma}
\label{lem:ch48-cln}
Let \(\phi\in\PSH(\omega)\) with \(\sup_{X^{\An}}\phi=0\), and let
\(\psi\in\mathcal U_\omega\).  Then
\[
 \int_{X^{\An}}(-\phi)\,\MA_\omega(\psi)\ \le\ n .
\]
More generally, for \(\psi_1,\dots,\psi_n\in\mathcal U_\omega\),
\(\int(-\phi)\,\MA_\omega(\psi_1,\dots,\psi_n)\le n\).
\end{lemma}

\begin{proof}
We prove the mixed statement; the first is the diagonal case.  Assume first
that \(\phi\) is bounded, so that \(\phi\in\cE^1(\omega)\) by
\Cref{thm:bj-e1-basic} and all the integrals below are finite.  Write
\(\MA_\omega(\cdot)\) for the normalized mixed measures of
\Cref{thm:bj-mixed-ma}, each of total mass \(V^{-1}(\omega^n)=1\), and for
\(0\le k\le n\) put
\[
 \nu_k:=V^{-1}(\omega+\ddc\psi_1)\wedge\cdots\wedge(\omega+\ddc\psi_k)
        \wedge\omega^{\,n-k},
\]
a probability measure.  Moreover,
\[
 \nu_0=\MA_\omega(0)=\delta_{v_{\triv}},
 \qquad
 \nu_n=\MA_\omega(\psi_1,\dots,\psi_n).
\]
The first identity is
\Cref{lem:ch47-three-normalizations}\,\cref{itm:ch47-norm:trivial-measure}.
Since
\(\sup\phi=\phi(v_{\triv})=0\)
(\Cref{lem:ch47-three-normalizations}\,\cref{itm:ch47-norm:max}), we have
\(\int(-\phi)\,\rd\nu_0=0\), and therefore
\[
 \int(-\phi)\,\MA_\omega(\psi_1,\dots,\psi_n)
 =\sum_{k=0}^{n-1}\int(-\phi)\,\rd(\nu_{k+1}-\nu_k)
 =\sum_{k=0}^{n-1}\int(-\phi)\,\ddc\psi_{k+1}\wedge T_k ,
\]
where \(T_k:=V^{-1}(\omega+\ddc\psi_1)\wedge\cdots\wedge
(\omega+\ddc\psi_k)\wedge\omega^{\,n-1-k}\).

Fix \(k\) and abbreviate \(\psi:=\psi_{k+1}\), \(T:=T_k\).  By integration
by parts \eqref{eq:bj-integration-by-parts}, applied to the finite-energy
pair \((\phi,\psi)\) and the mixed weight \(T\),
\[
 \int(-\phi)\,\ddc\psi\wedge T=\int\psi\,\ddc(-\phi)\wedge T
 =\int(-\psi)\,\ddc\phi\wedge T .
\]
Now \(\ddc\phi=(\omega+\ddc\phi)-\omega\) and \(0\le-\psi\le1\), so
\[
 \begin{aligned}
  \int(-\psi)\,\ddc\phi\wedge T
   &=\int(-\psi)(\omega+\ddc\phi)\wedge T
     -\int(-\psi)\,\omega\wedge T\\
   &\le\int(-\psi)(\omega+\ddc\phi)\wedge T\le1,
 \end{aligned}
\]
the last step because \(V^{-1}(\omega+\ddc\phi)\wedge T\) is a positive
measure of total mass \(1\).  Summing over \(k\) gives the bound \(n\).

For general \(\phi\in\PSH(\omega)\) with \(\sup\phi=0\), apply the above to
\(\phi_j:=\max\{\phi,-j\}\), which is \(\omega\)-psh
(\Cref{thm:bj-basic-closure}), bounded, and has \(\sup\phi_j=0\).  Since
\(-\phi_j\uparrow-\phi\) pointwise and all the functions are nonnegative,
monotone convergence gives
\[
 \int(-\phi)\MA_\omega(\psi_1,\dots,\psi_n)
  =\lim_j\int(-\phi_j)\MA_\omega(\psi_1,\dots,\psi_n)\le n.
\]
\end{proof}

\begin{remark}[Where the constant comes from]
\label{rem:ch48-cln-constant}
The constant \(n\) is the number of factors that have to be exchanged, and
each exchange costs the total mass \(1\) of one normalized mixed measure.
Over a discretely valued base the same computation is
\Cref{lem:ma-capacity-bound} and produces \(nV\), because there the mixed
measures have mass \(V\); the conversion is
\Cref{warn:ch48-normalization}.  Two features of the trivially valued
argument should be noted.  The telescoping starts at
\(\nu_0=\delta_{v_{\triv}}\) and the boundary term
\(\int(-\phi)\rd\nu_0\) \emph{vanishes} because \(v_{\triv}\) is where every
psh function attains its maximum; over a discretely valued base the
corresponding term is \(\int(-\varphi)\theta^n\), which is not zero and
which \Cref{lem:ch25-normalized-bound} has to estimate separately, at the
cost of an additive constant \(C_\theta\).  This is
\Cref{warn:ch47-J-identity} appearing once more.
\end{remark}

\subsection{Sublevel sets and the energy criterion}
\label{sec:ch48-sublevel}

\begin{proposition}\label{prop:bj-capacity-sublevel}
If \(\phi\in\PSH(\omega)\) and \(\sup\phi=0\), then for \(t>0\)

\[
                  \Capa(\{\phi\leq-t\})\leq \frac nt.
\tag{48.8}\label{eq:bj-cap-sublevel-bound}
\]

Moreover,

\[
 J(\phi)\leq 1+\int_1^\infty t^n
             \Capa(\{\phi\leq-t\})\,\rd t.
\tag{48.9}\label{eq:bj-cap-energy-criterion}
\]

In particular, finiteness of the integral implies
\(\phi\in\mathcal E^1\).
\end{proposition}

\begin{proof}
For \(-1\leq\psi\leq0\), Chebyshev and the Chern--Levine--Nirenberg
estimate give

\[
 \int_{\{\phi\leq-t\}}\MA(\psi)
 \leq t^{-1}\int(-\phi)\MA(\psi)\leq nt^{-1}.
\]

Taking the supremum proves \eqref{eq:bj-cap-sublevel-bound}.  For
\(\phi_s=\max\{\phi,-s\}\), write the layer-cake formula for
\(\int(-\phi_s)\MA(\phi_s)\).  Locality identifies the measures of the
truncations on the appropriate superlevel sets, while
\(t^{-1}\phi_t\) is admissible in the capacity.  Letting
\(s\to\infty\) gives \eqref{eq:bj-cap-energy-criterion}.  Compare
\cite[Lemma~7.47]{BJ22}.

We now give the argument in full; no import is used.

\emph{Step 1: \eqref{eq:bj-cap-sublevel-bound}.}  The set
\(\{\phi\le-t\}\) is Borel, \(\phi\) being upper semicontinuous, and
\(-\phi\ge0\) on it with \(-\phi\ge t\).  For \(\psi\in\mathcal U_\omega\),
Chebyshev's inequality and \Cref{lem:ch48-cln} give
\(\MA_\omega(\psi)(\{\phi\le-t\})\le t^{-1}\int(-\phi)\MA_\omega(\psi)
\le n/t\).  Taking the supremum over \(\psi\in\mathcal U_\omega\) proves
\eqref{eq:bj-cap-sublevel-bound}.  The same argument applied to the Borel
set \(\{\phi<-t\}\subset\{\phi\le-t\}\) gives
\(\Capa_\omega(\{\phi<-t\})\le n/t\); since \(\{\phi<-t\}\) is open, this
is also a bound for the \emph{outer} capacity of any subset of
\(\{\phi<-t\}\), and it is in that form that we use it below.

\emph{Step 2: a lower bound for the capacity of a sublevel set.}  Let
\(\phi\in\cE^1(\omega)\) be bounded with \(\sup\phi=0\), and let
\(t\ge1\).  Put
\[
 u_t:=\max\Bigl\{\frac\phi t,-1\Bigr\}.
\]
Since \(t\ge1\), the function \(\phi/t=\tfrac1t\phi+(1-\tfrac1t)\cdot0\) is
\(\omega\)-psh by convexity of \(\PSH(\omega)\)
(\Cref{thm:bj-basic-closure}), and so is \(u_t\), a maximum of two
\(\omega\)-psh functions; moreover \(-1\le u_t\le0\), so
\(u_t\in\mathcal U_\omega\).  On the open set \(U:=\{\phi>-t\}\) we have
\(\phi/t>-1\), hence \(u_t=\phi/t\); by locality
\eqref{eq:bj-locality} applied to the pair \((\phi/t,-1)\),
\[
 \mathbf 1_U\,\MA_\omega(u_t)=\mathbf 1_U\,\MA_\omega(\phi/t).
\]
Both measures being probabilities, taking complements gives
\(\MA_\omega(u_t)(\{\phi\le-t\})=\MA_\omega(\phi/t)(\{\phi\le-t\})\).
Expanding
\(\omega+\ddc(\phi/t)=\tfrac1t(\omega+\ddc\phi)+(1-\tfrac1t)\omega\) by
multilinearity (\Cref{thm:bj-mixed-ma}),
\[
 \MA_\omega(\phi/t)
 =\sum_{k=0}^n\binom nk t^{-k}\Bigl(1-\frac1t\Bigr)^{n-k}\mu_k,
 \qquad
 \mu_k:=V^{-1}(\omega+\ddc\phi)^k\wedge\omega^{\,n-k},
\]
all the \(\mu_k\) being positive of mass \(1\); dropping every term but
\(k=n\) gives \(\MA_\omega(\phi/t)\ge t^{-n}\MA_\omega(\phi)\).  Hence
\[
 \Capa_\omega\bigl(\{\phi\le-t\}\bigr)
 \ \ge\ \MA_\omega(u_t)\bigl(\{\phi\le-t\}\bigr)
 \ \ge\ t^{-n}\,\MA_\omega(\phi)\bigl(\{\phi\le-t\}\bigr),
 \qquad t\ge1 .
\]

\emph{Step 3: the layer-cake estimate.}  Keep \(\phi\) bounded with
\(\sup\phi=0\), and put \(\mu_k\) as above.  The telescoping identity
\eqref{eq:ch46-pure-pairing-as-integrals}, which is stated for PL data and
extends to finite-energy data by \Cref{thm:bj-mixed-ma} and
\Cref{thm:bj-extended-energy-pairing}, together with the definition
\eqref{eq:bj-energy-definition} of \(E_\omega\), gives
\[
 -E_\omega(\phi)=\frac1{n+1}\sum_{k=0}^n\int(-\phi)\,\rd\mu_k .
\]
The sequence \(k\mapsto\int(-\phi)\rd\mu_k\) is nondecreasing, because
\(\int(-\phi)\rd(\mu_{k+1}-\mu_k)
=-V^{-1}\int\phi\,\ddc\phi\wedge(\omega+\ddc\phi)^k\wedge\omega^{\,n-1-k}
\ge0\) by positivity of the Dirichlet form
(\Cref{cor:ch46-dirichlet-e1}).  Hence
\[
 J(\phi)=-E_\omega(\phi)\ \le\ \int(-\phi)\,\rd\mu_n
 =\int(-\phi)\,\MA_\omega(\phi).
\]
By the layer-cake formula and Step~2,
\[
 \begin{aligned}
  \int(-\phi)\,\MA_\omega(\phi)
   &=\int_0^\infty
     \MA_\omega(\phi)\bigl(\{\phi<-s\}\bigr)\,\rd s\\
   &\le1+\int_1^\infty s^n
     \Capa_\omega\bigl(\{\phi\le-s\}\bigr)\,\rd s,
 \end{aligned}
\]
where on \([0,1]\) we used the crude bound
\(\MA_\omega(\phi)(\{\phi<-s\})\le1\).  This is
\eqref{eq:bj-cap-energy-criterion} for bounded \(\phi\).

\emph{Step 4: removing the boundedness.}  Let \(\phi\in\PSH(\omega)\) with
\(\sup\phi=0\) and put \(\phi_j:=\max\{\phi,-j\}\), so that
\(\phi_j\in\cE^1(\omega)\) is bounded with \(\sup\phi_j=0\) and
\(\phi_j\downarrow\phi\).  Since \(\phi_j\ge\phi\) we have
\(\{\phi_j\le-s\}\subset\{\phi\le-s\}\), and \(\{\phi_j\le-s\}=\varnothing\)
for \(s>j\); so Step~3 gives
\[
 J(\phi_j)\ \le\ 1+\int_1^\infty s^n\,
   \Capa_\omega\bigl(\{\phi\le-s\}\bigr)\,\rd s=:B
\]
for every \(j\), with \(B\) independent of \(j\).  By
\Cref{prop:ch46-energy-usc}\,\cref{itm:ch46-Eusc:decreasing},
\(E_\omega(\phi_j)\to E_\omega(\phi)\), so
\(J(\phi)=-E_\omega(\phi)=\lim_jJ(\phi_j)\le B\), which is
\eqref{eq:bj-cap-energy-criterion} in general and, when \(B<\infty\), shows
\(E_\omega(\phi)>-\infty\), that is, \(\phi\in\cE^1(\omega)\).
\end{proof}

\begin{remark}[The criterion is not vacuous, and it is sharp]
\label{rem:ch48-criterion-sharp}
The two bounds of \Cref{prop:bj-capacity-sublevel} point in opposite
directions and are compatible: \eqref{eq:bj-cap-sublevel-bound} gives
\(t^n\Capa(\{\phi\le-t\})\le nt^{n-1}\), whose integral over \([1,\infty)\)
diverges for every \(n\ge1\).  So the criterion
\eqref{eq:bj-cap-energy-criterion} really does test something: it asks the
capacity of the deep sublevel sets to decay faster than the trivial rate.
\Cref{ex:ch48-p1-capacity} exhibits, on \(\mathbf P^1\), a potential of
infinite energy for which \(\Capa(\{\phi\le-t\})=t^{-1}\) exactly, so that
the integrand in \eqref{eq:bj-cap-energy-criterion} is the constant \(1\)
and the integral diverges; and a potential of finite energy for which the
sublevel capacities decay like \(t^{-2}\) and the integral converges.  The
same pair of computations over a discretely valued base is
\Cref{ex:ch25-pluripolar-computed}, where \eqref{eq:ch25-p1-capacity} gives
\(\Capa_\theta(\{\varphi_0\le-t\})=\min\{1,1/t\}\); the two examples agree
because in both cases the underlying object is a ray with a linear
potential, and the constants agree because both chapters normalize the
Laplacian by the sum of outgoing slopes.
\end{remark}

\subsection{The fundamental capacity estimate}
\label{sec:ch48-fundamental}

The next lemma is the inequality that makes quasicontinuity work.  Its role
is to convert a requirement that is \emph{uniform} in the admissible
potential --- which is what a capacity statement needs, the capacity being
a supremum --- into a single condition involving one fixed measure.  It is
the trivially valued form of \Cref{lem:fundamental-capacity-estimate}, and
the proof transfers verbatim; we write it out because the constants change
under the probability normalization and because it is the clearest example
in this chapter of an argument that does transfer.

\begin{lemma}
\label{lem:ch48-fundamental}
Let \(M\ge1\) and let \(\varphi,\psi,u_1,\dots,u_n\in\PSH(\omega)\) satisfy
\[
 -M\le\varphi\le\psi\le0,\qquad -M\le u_i\le0 .
\]
Then
\[
 0\ \le\ \int(\psi-\varphi)\,\MA_\omega(u_1,\dots,u_n)
 \ \le\ 4M\left(\int(\psi-\varphi)\,\MA_\omega(\varphi/2)\right)^{2^{-n}} .
\]
In particular, for \(-1\le\varphi\le\psi\le0\),
\[
 \sup_{u\in\mathcal U_\omega}\int(\psi-\varphi)\,\MA_\omega(u)
 \ \le\ 4\left(\int(\psi-\varphi)\,\MA_\omega(\varphi/2)\right)^{2^{-n}} .
\]
\end{lemma}

\begin{proof}
All the entries are bounded \(\omega\)-psh functions, hence lie in
\(\cE^1(\omega)\) by \Cref{thm:bj-e1-basic}; all the measures below are
positive of total mass \(1\) by \Cref{thm:bj-mixed-ma} and
\Cref{warn:ch48-normalization}; and all the manipulations are the
multilinear ones of \Cref{thm:bj-mixed-ma} together with integration by
parts \eqref{eq:bj-integration-by-parts}.  Write \(f:=\psi-\varphi\), so
\(0\le f\le M\).  The left-hand inequality is positivity of the mixed
measure.  For an \((n-1)\)-fold normalized mixed product
\(T=V^{-1}(\omega+\ddc w_2)\wedge\cdots\wedge(\omega+\ddc w_n)\) of bounded
\(\omega\)-psh functions write
\(\langle a,b\rangle_T:=-\int a\,\ddc b\wedge T\), the Dirichlet form of
\Cref{cor:ch46-dirichlet-e1}, which is symmetric, positive semidefinite,
and obeys Cauchy--Schwarz.

\emph{Step 1: the exchange inequality.}  We claim that for every such
\(T\) and every bounded \(\omega\)-psh \(w\) with \(-M\le w\le0\),
\[
 \int f\,(\omega+\ddc w)\wedge T
 \ \le\ 2M^{1/2}
   \Bigl(\int f\,\bigl(\omega+\ddc\tfrac\varphi2\bigr)\wedge T\Bigr)^{1/2}.
\]
Split the free factor:
\(\int f(\omega+\ddc w)\wedge T=A+\int f\,\ddc w\wedge T\) with
\(A:=\int f\,\omega\wedge T\).  The measure \(\omega\wedge T\) is positive
of mass \(1\), so \(0\le A\le M\) and \(A\le M^{1/2}A^{1/2}\).  For the
second term, Cauchy--Schwarz gives
\(\int f\,\ddc w\wedge T=-\langle f,w\rangle_T
\le\langle f,f\rangle_T^{1/2}\langle w,w\rangle_T^{1/2}\), and
\[
 \langle w,w\rangle_T=\int(-w)(\omega+\ddc w)\wedge T
   -\int(-w)\,\omega\wedge T\ \le\ M,
\]
because \(0\le-w\le M\), the subtracted term is nonnegative, and
\((\omega+\ddc w)\wedge T\) has mass \(1\).  Likewise, since
\(f=\psi-\varphi\),
\[
 \begin{aligned}
 \langle f,f\rangle_T
 &=\int f\,\ddc(\varphi-\psi)\wedge T\\
 &=\int f(\omega+\ddc\varphi)\wedge T
   -\int f(\omega+\ddc\psi)\wedge T\\
 &\le C:=\int f(\omega+\ddc\varphi)\wedge T ,
 \end{aligned}
\]
because \(f\ge0\) and \((\omega+\ddc\psi)\wedge T\) is positive.
Combining, \(\int f(\omega+\ddc w)\wedge T\le M^{1/2}(A^{1/2}+C^{1/2})\),
and \(A^{1/2}+C^{1/2}\le\sqrt2(A+C)^{1/2}\) by Cauchy--Schwarz in
\(\mathbf R^2\), while linearity of \(\ddc\) gives
\(A+C=\int f(2\omega+\ddc\varphi)\wedge T
=2\int f(\omega+\ddc\tfrac\varphi2)\wedge T\).  The claim follows.  Note
that \(\varphi/2=\tfrac12\varphi+\tfrac12\cdot0\) is \(\omega\)-psh, so the
right-hand side is again an integral against a positive measure of mass
\(1\).

\emph{Step 2: iteration.}  For \(0\le k\le n\) put
\[
 J_k:=\int f\,V^{-1}\bigl(\omega+\ddc\tfrac\varphi2\bigr)^{\wedge k}
   \wedge(\omega+\ddc u_{k+1})\wedge\cdots\wedge(\omega+\ddc u_n),
\]
so that \(J_0=\int f\,\MA_\omega(u_1,\dots,u_n)\) and
\(J_n=\int f\,\MA_\omega(\varphi/2)\).  Step~1, applied with
\(w=u_{k+1}\), gives \(J_k\le2M^{1/2}J_{k+1}^{1/2}\) for
\(0\le k\le n-1\).  Iterating,
\[
 J_0\le\prod_{k=0}^{n-1}\bigl(2M^{1/2}\bigr)^{2^{-k}}J_n^{2^{-n}}
 =2^{2(1-2^{-n})}M^{1-2^{-n}}J_n^{2^{-n}}
 \le 4M\,J_n^{2^{-n}},
\]
using \(\sum_{k=0}^{n-1}2^{-k}=2(1-2^{-n})\), then
\(2^{2(1-2^{-n})}\le4\) and \(M^{1-2^{-n}}\le M\) for \(M\ge1\).  The
uniform form is the case \(M=1\), \(u_1=\dots=u_n=u\in\mathcal U_\omega\).
\end{proof}

\begin{remark}[What transfers, and what the transfer costs]
\label{rem:ch48-fundamental-transfer}
Three comments, all of which are corrections to
\Cref{rem:ch25-reading-fundamental} for the present normalization.

\emph{The constants.}  The factor \(V^{1-2^{-n}}\) of
\eqref{eq:fundamental-capacity-estimate} has disappeared, because every
mixed measure here has mass \(1\).  A reader converting a numerical
statement from \Cref{chap:bfj-capacity} must restore it; the exponent
\(2^{-n}\) is unchanged, and neither constant is claimed to be optimal.

\emph{The measure on the right is fixed by \(\varphi\) alone.}  This is the
whole point.  In every application below, \(\varphi\) is the limit of a
\emph{decreasing} net \((\psi_\alpha)\), and one applies the lemma with
\(\psi=\psi_\alpha\); the measure \(\MA_\omega(\varphi/2)\) does not move
with \(\alpha\), so
\(\int(\psi_\alpha-\varphi)\MA_\omega(\varphi/2)\to0\) is a statement about
a \emph{single} Radon measure and a decreasing net of upper semicontinuous
functions, which is \Cref{cor:usc-net-monotone-integral}.

\emph{Why it does not serve increasing nets.}  If instead
\(\psi_\alpha\uparrow\) with regularized supremum \(\varphi\), then the
roles are reversed: the varying function is the \emph{smaller} one, and the
measure on the right of the lemma is \(\MA_\omega(\psi_\alpha/2)\), which
moves.  No rearrangement of the proof repairs this, because the step
\(\langle f,f\rangle_T\le\int f(\omega+\ddc\varphi)\wedge T\) discards the
term built from the larger function and there is no matching inequality in
the other direction.  This is the precise reason that
\Cref{imp:ch48-imported}\,\cref{itm:ch48-import:increasing} and
\cref{itm:ch48-import:negligible} are imports rather than corollaries of
\Cref{lem:ch48-fundamental}, and it is the same asymmetry as
\Cref{warn:ch48-monotone-asymmetry}.
\end{remark}

\section{Quasicontinuity and convergence in capacity}
\label{sec:ch48-quasicontinuity}

\subsection{Quasicontinuity}
\label{sec:ch48-qc}

\begin{theorem}\label{thm:bj-quasicontinuity}
For every \(\phi\in\PSH(\omega)\) and \(\varepsilon>0\), there is a compact
set \(K\subset X^{\An}\) such that

\[
                  \Capa(K^c)\leq\varepsilon,
       \qquad \phi|_K\in C^0(K).
\tag{48.10}\label{eq:bj-quasicontinuity}
\]

One compact \(K\) may be chosen simultaneously for a countable family of
psh functions.
\end{theorem}

\begin{proof}
For bounded \(\phi\), choose a decreasing net of continuous \(\omega\)-psh
regularizations.  Strong measure estimates and Chebyshev yield indices
\(g_m\) such that the error exceeds \(m^{-1}\) only on a set of capacity at
most \(2^{-m}\varepsilon\).  On the intersection of the complementary
compact sets, the selected regularization converges uniformly.  For
unbounded \(\phi\), first remove a deep sublevel set using
\eqref{eq:bj-cap-sublevel-bound}, then apply the bounded argument to a
truncation.  A countable diagonal union proves the simultaneous statement.
Compare \cite[Theorem~13.9 and Corollary~13.10]{BJ22}.

We now carry this out.  No import is used; the two inputs are
\Cref{lem:ch48-fundamental} and
\Cref{imp:ch48-envelope}\,\cref{itm:ch48-env:regularization}.

\emph{Step 1: the bounded case.}  Assume \(\phi\in\PSH(\omega)\) is
bounded.  Replacing \(\phi\) by \(\phi-\sup\phi\) --- which changes neither
the conclusion nor the hypothesis --- we may assume
\(\sup_{X^{\An}}\phi=0\), so \(-M\le\phi\le0\) for some \(M\ge1\).
By \Cref{thm:ch45-continuous-regularization} there is a downward directed
net \((u_g)_{g\in\mathcal G}\) of continuous \(\omega\)-psh functions with
\(u_g\downarrow\phi\) pointwise on \(X^{\An}\).  Since
\(\sup u_g=u_g(v_{\triv})\downarrow\phi(v_{\triv})=0\), there is
\(g_0\) with \(u_{g}\le1\) for all \(g\ge g_0\); replacing
\(\mathcal G\) by \(\{g\ge g_0\}\), which is cofinal and still downward
directed, we may assume \(u_g\le1\) for every \(g\).

Apply \Cref{lem:ch48-fundamental} to the pair
\(\varphi:=\phi-1\) and \(\psi:=u_g-1\), which satisfy
\(-(M+1)\le\varphi\le\psi\le0\), with the constant \(M+1\ge1\).  Since
\(\psi-\varphi=u_g-\phi\) and
\(\MA_\omega(\varphi/2)=\MA_\omega(\phi/2-\tfrac12)=\MA_\omega(\phi/2)\),
we obtain, for every \(u\in\mathcal U_\omega\),
\[
 \int(u_g-\phi)\,\MA_\omega(u)
 \ \le\ 4(M+1)
   \left(\int(u_g-\phi)\,\MA_\omega(\phi/2)\right)^{2^{-n}} .
\]
The measure \(\nu:=\MA_\omega(\phi/2)\) is a fixed Radon probability
measure, and \((u_g)\) is a decreasing net of continuous functions with
pointwise limit \(\phi\); by
\Cref{cor:usc-net-monotone-integral},
\(\int u_g\,\rd\nu\downarrow\int\phi\,\rd\nu\), a finite limit because
\(\phi\) is bounded.  Hence
\[
 \eta_g:=\int(u_g-\phi)\,\rd\nu\ \longrightarrow\ 0 .
\]
For \(\delta>0\) the set \(\{u_g-\phi>\delta\}\) is \emph{open}, because
\(u_g\) is continuous and \(\phi\) is upper semicontinuous, so
\(u_g-\phi\) is lower semicontinuous.  Chebyshev's inequality gives, for
every \(u\in\mathcal U_\omega\),
\(\MA_\omega(u)(\{u_g-\phi>\delta\})\le\delta^{-1}\int(u_g-\phi)\MA_\omega(u)\),
whence
\[
 \Capa_\omega\bigl(\{u_g-\phi>\delta\}\bigr)
 \ \le\ \frac{4(M+1)}{\delta}\,\eta_g^{2^{-n}} .
\]

Now let \(\varepsilon>0\).  Using that \(\mathcal G\) is downward directed
and that \(\eta_g\to0\), choose inductively
\(g_1\ge g_2\ge\cdots\) in \(\mathcal G\) with
\[
 \Capa_\omega\bigl(\{u_{g_m}-\phi>\tfrac1m\}\bigr)\le2^{-m}\varepsilon,
 \qquad m\ge1 .
\]
Put \(U:=\bigcup_{m\ge1}\{u_{g_m}-\phi>\tfrac1m\}\), an open set with
\(\Capa_\omega(U)\le\varepsilon\) by countable subadditivity
(\Cref{prop:ch48-capacity-elementary}\,\cref{itm:ch48-cap:sub}), and
\(K:=X^{\An}\setminus U\), a closed subset of a compact space, hence
compact.  On \(K\) we have \(0\le u_{g_m}-\phi\le1/m\) for every \(m\), so
\(u_{g_m}|_K\to\phi|_K\) uniformly; each \(u_{g_m}|_K\) is continuous, and a
uniform limit of continuous functions on \(K\) is continuous.  This proves
\eqref{eq:bj-quasicontinuity} for bounded \(\phi\).

\emph{Step 2: the general case.}  Let \(\phi\in\PSH(\omega)\), normalized
by \(\sup\phi=0\), and let \(\varepsilon>0\).  Choose \(t\ge1\) with
\(n/t\le\varepsilon/2\).  The set \(G:=\{\phi<-t\}\) is open and, by
Step~1 of the proof of \Cref{prop:bj-capacity-sublevel},
\(\Capa_\omega(G)\le n/t\le\varepsilon/2\).  The truncation
\(\phi^{(t)}:=\max\{\phi,-t-1\}\) is bounded and \(\omega\)-psh, so Step~1
provides an open \(U'\) with \(\Capa_\omega(U')\le\varepsilon/2\) and
\(\phi^{(t)}|_{X^{\An}\setminus U'}\) continuous.  Put
\(K:=X^{\An}\setminus(G\cup U')\), a compact set with
\(\Capa_\omega(K^c)\le\varepsilon\) by subadditivity.  On \(K\) we have
\(\phi\ge-t>-t-1\), hence \(\phi=\phi^{(t)}\), and therefore
\(\phi|_K\) is continuous.

\emph{Step 3: countable families.}  Given \((\phi_j)_{j\ge1}\) in
\(\PSH(\omega)\) and \(\varepsilon>0\), apply Step~2 to each \(\phi_j\)
with \(2^{-j}\varepsilon\), obtaining compact \(K_j\) with
\(\Capa_\omega(K_j^c)\le2^{-j}\varepsilon\) and \(\phi_j|_{K_j}\)
continuous.  Then \(K:=\bigcap_jK_j\) is compact,
\(\Capa_\omega(K^c)=\Capa_\omega(\bigcup_jK_j^c)\le\varepsilon\), and each
\(\phi_j|_K\) is continuous, being the restriction of a continuous function
on the subspace \(K\subset K_j\).
\end{proof}

\begin{corollary}
\label{cor:ch48-qc-calculus}
Every \(\phi\in\PSH(\omega)\) is quasicontinuous for \(\Capa_\omega\) in
the sense of \Cref{sec:capacity-preview}.  Consequently, finite sums,
finite maxima and minima, products of bounded members, and compositions
with continuous \(\chi:\Rbar\to\Rbar\) of \(\omega\)-psh functions are
quasicontinuous, as are uniform limits of quasicontinuous functions and
limits that are uniform outside sets of arbitrarily small capacity.
\end{corollary}

\begin{proof}
Quasicontinuity is \Cref{thm:bj-quasicontinuity}, the exceptional set
\(K^c\) being open.  The remaining assertions are
\Cref{prop:quasicontinuity-calculus}, whose proof uses only monotonicity
and countable subadditivity of the capacity, both of which are
\Cref{prop:ch48-capacity-elementary}\,\cref{itm:ch48-cap:mono},
\cref{itm:ch48-cap:sub}.  In particular no appeal to the Choquet axioms
\cite[Theorem~13.29]{BJ22} is needed here.
\end{proof}

\begin{warning}{What quasicontinuity does not say}
\label{warn:ch48-qc-not}
\Cref{thm:bj-quasicontinuity} says that \(\phi|_K\) is continuous for the
subspace topology on \(K\).  It does not say that \(\phi\) is continuous at
the points of \(K\) as a function on \(X^{\An}\), and it does not say that
\(\phi\) agrees off a small set with a continuous function on \(X^{\An}\).
Nor does it produce a \emph{single} \(K\) working for all
\(\varepsilon\): the compact sets grow as \(\varepsilon\) shrinks, and
their union is in general a proper subset of \(X^{\An}\).  Finally, the
statement is about a countable family; for an uncountable family the
diagonal argument of Step~3 fails, and
\Cref{ex:ch48-proper-support-capacity} shows that the failure is not
formal.  These are the same three cautions as
\Cref{rem:ch25-qc-not} and \Cref{rem:quasicontinuity-limits}.
\end{warning}

\subsection{Convergence in capacity}
\label{sec:ch48-convergence-capacity}

\begin{proposition}
\label{prop:ch48-decreasing-capacity}
Let \((u_g)\) be a downward directed net in \(\PSH(\omega)\) with
pointwise limit \(\phi\in\PSH(\omega)\), and suppose that \(\phi\) is
bounded and that \(u_g\le C_0\) for some constant.  Then
\(u_g\to\phi\) in capacity: for every \(\delta>0\),
\[
 \Capa_\omega\bigl(\{u_g-\phi>\delta\}\bigr)\longrightarrow0 .
\]
The same holds for an arbitrary \(\phi\in\PSH(\omega)\) after removing a
sublevel set: for every \(\varepsilon>0\) there is \(t\) with
\(\Capa_\omega(\{\phi<-t\})\le\varepsilon\), and on the complement the
truncated net converges in capacity.
\end{proposition}

\begin{proof}
The first assertion is the estimate displayed in Step~1 of the proof of
\Cref{thm:bj-quasicontinuity}, which used only that \(\phi\) is bounded and
that the net is decreasing with continuous --- in fact merely bounded above
--- members: \Cref{lem:ch48-fundamental} applies to
\(\varphi=\phi-C_0\) and \(\psi=u_g-C_0\), and
\Cref{cor:usc-net-monotone-integral} gives
\(\int(u_g-\phi)\,\MA_\omega(\phi/2)\to0\), the members \(u_g\) being upper
semicontinuous and the measure fixed.  The second assertion combines this
with \eqref{eq:bj-cap-sublevel-bound}, exactly as in Step~2 of that proof.
\end{proof}

\begin{proposition}
\label{prop:ch48-capacity-subsequence}
If \((\phi_j)_{j\ge1}\) is a \emph{sequence} of Borel functions converging
to \(\phi\) in capacity, then some subsequence converges to \(\phi\)
pointwise outside a set of outer capacity zero.
\end{proposition}

\begin{proof}
This is \Cref{prop:capacity-subsequence}, whose hypotheses --- countable
subadditivity and a sequential index --- are
\Cref{prop:ch48-capacity-elementary}\,\cref{itm:ch48-cap:sub} and the
statement.  The exceptional set is a countable union of sets of capacity at
most \(2^{-m}\), and each of them may be enlarged to an open set of
capacity at most \(2^{1-m}\), so the union has outer capacity zero.
\end{proof}

\begin{remark}[Nets do not admit this extraction]
\label{rem:ch48-nets}
\label{sec:ch48-nets}
\Cref{prop:ch48-capacity-subsequence} is deliberately sequential.  A net
converging in capacity need not contain a useful sequence: the argument
bounds the capacity of a countable union, and an uncountable directed set
supplies none.  Over a discretely valued base the difficulty is met by
\Cref{cor:psh-sequential-regularization}, which extracts a sequence from a
regularizing net because the model system, although not cofinally
countable, admits countable cofinal subsystems adapted to a single
potential.  Here the corresponding statement is Step~1 of the proof of
\Cref{thm:bj-quasicontinuity}: the sequence \((g_m)\) is extracted from the
directed set \(\mathcal G\) by hand, using only that
\(\eta_g\to0\) along the net, and the extraction is possible precisely
because the quantity being made small is a single real number.  Whenever a
statement below is about an uncountable family, it is stated for nets and
no sequence is extracted; \Cref{warn:ch45-nets} is the governing
discipline.
\end{remark}


\section{Polar, negligible, and capacity-zero sets}
\label{sec:ch48-small-sets}

\subsection{The three notions}
\label{sec:ch48-three-notions}

\begin{definition}[Three smallness notions]
\label{def:bj-three-small-sets}
A set \(A\subset X^{\An}\) is:

\begin{enumerate}
\item \emph{pluripolar} if
      \(A\subset\{\phi=-\infty\}\) for some psh \(\phi\not\equiv-\infty\);
\item \emph{negligible} if there are a class \(\theta\) and a bounded-above
      family \((\phi_i)\subset\PSH(\theta)\) such that
\[
 A\subset\left\{\sup_i\phi_i<(\sup_i\phi_i)^*\right\};
\tag{48.11}\label{eq:bj-negligible}
\]
\item of \emph{capacity zero} if \(\Capa^*(A)=0\).
\end{enumerate}
\end{definition}

The raw supremum in \eqref{eq:bj-negligible} is evaluated pointwise in the
extended reals.  It must not be silently replaced by its regularization.

\begin{lemma}
\label{lem:ch48-normalize-pole}
Let \(A\subset X^{\An}\) be pluripolar.  Then there is
\(\phi\in\PSH(\omega)\) with \(\sup_{X^{\An}}\phi=0\) and
\(A\subset\{\phi=-\infty\}\).  Conversely such a \(\phi\) exhibits \(A\) as
pluripolar.
\end{lemma}

\begin{proof}
Let \(\psi\in\PSH(\theta)\) with \(\psi\not\equiv-\infty\) and
\(A\subset\{\psi=-\infty\}\).  Since \(\Amp(X)\) is an open convex cone in
\(N^1(X)_{\mathbf R}\) and \(\omega\in\Amp(X)\), there is \(\varepsilon>0\)
with \(\omega-\varepsilon\theta\in\Amp(X)\).  By
\Cref{thm:bj-basic-closure}\,(3) with \(s=\varepsilon\), \(t=0\), we have
\(\varepsilon\psi\in\PSH(\varepsilon\theta)\); by
\Cref{lem:ch34-psh-nonempty} the constant \(0\) lies in
\(\PSH(\omega-\varepsilon\theta)\), the class being ample hence nef; and by
\Cref{thm:bj-basic-closure}\,(3) again,
\(\varepsilon\psi=\varepsilon\psi+0\in\PSH(\omega)\).  Subtracting the
finite constant \(\sup_{X^{\An}}\varepsilon\psi\)
(\Cref{thm:bj-basic-closure}\,(1)) normalizes the supremum, and
\(\{\varepsilon\psi=-\infty\}=\{\psi=-\infty\}\supset A\).  The converse is
the definition.
\end{proof}

\begin{proposition}
\label{prop:ch48-pluripolar-capacity}
If \(A\subset X^{\An}\) is pluripolar, then \(\Capa^*_\omega(A)=0\).  More
precisely, with \(\phi\) as in \Cref{lem:ch48-normalize-pole},
\[
 A\subset\{\phi<-t\}\quad\text{and}\quad
 \Capa_\omega\bigl(\{\phi<-t\}\bigr)\le\frac nt
 \qquad\text{for every }t>0 .
\]
\end{proposition}

\begin{proof}
The set \(\{\phi<-t\}\) is open, because \(\phi\) is upper semicontinuous,
and it contains \(\{\phi=-\infty\}\supset A\).  Its capacity is at most
\(n/t\) by Step~1 of the proof of \Cref{prop:bj-capacity-sublevel}.  Hence
\(\Capa^*_\omega(A)\le\inf_{t>0}n/t=0\).
\end{proof}

\begin{proposition}
\label{prop:ch48-measure-vs-capacity}
Let \(\mu\in\cM^1(\omega)\) and let \(B\subset X^{\An}\) be Borel.  Then
for every \(t\ge1\)
\[
 \mu(B)\ \le\ t^n\,\Capa_\omega(B)+\frac{(n+1)E^\vee_\omega(\mu)}{t} .
\]
Consequently:
\begin{enumerate}
\item\label{itm:ch48-mvc:null}
      if \(\Capa_\omega(B)=0\) then \(\mu(B)=0\);
\item\label{itm:ch48-mvc:outer}
      if \(\Capa^*_\omega(A)=0\) for an arbitrary \(A\subset X^{\An}\),
      then \(A\) is contained in a Borel set of \(\mu\)-measure zero; in
      particular every pluripolar set is \(\mu\)-null;
\item\label{itm:ch48-mvc:holder}
      if \(\Capa_\omega(B)\le(n+1)E^\vee_\omega(\mu)\), then
      \[
       \mu(B)\ \le\ 2\,\Capa_\omega(B)^{\frac1{n+1}}
        \bigl((n+1)E^\vee_\omega(\mu)\bigr)^{\frac n{n+1}} .
      \]
\end{enumerate}
\end{proposition}

\begin{proof}
By \Cref{cor:bj-cy-unibranch} there is \(\phi\in\cE^1_0(\omega)\) with
\(\MA_\omega(\phi)=\mu\); normalize \(\sup\phi=0\).  Fix \(t\ge1\) and split
\[
 B=\bigl(B\cap\{\phi>-t\}\bigr)\ \cup\ \bigl(B\cap\{\phi\le-t\}\bigr).
\]
On the second piece, \Cref{prop:ch47-no-polar-mass} gives
\[
 \mu\bigl(\{\phi\le-t\}\bigr)\le\frac{E^\vee_\omega(\mu)+J(\phi)}t
 \le\frac{(n+1)E^\vee_\omega(\mu)}t ,
\]
using \(J(\phi)\le nE^\vee_\omega(\mu)\) from
\Cref{thm:ch48-ma-homeomorphism}\,\cref{itm:ch48-ma:energy}.  On the first
piece, let \(u_t:=\max\{\phi/t,-1\}\in\mathcal U_\omega\) as in Step~2 of
the proof of \Cref{prop:bj-capacity-sublevel}; there we showed
\(\mathbf 1_{\{\phi>-t\}}\MA_\omega(u_t)
=\mathbf 1_{\{\phi>-t\}}\MA_\omega(\phi/t)\) and
\(\MA_\omega(\phi/t)\ge t^{-n}\MA_\omega(\phi)\).  Hence, for
\(B':=B\cap\{\phi>-t\}\),
\[
 \begin{aligned}
 \mu(B')=\MA_\omega(\phi)(B')
 &\le t^n\MA_\omega(\phi/t)(B')\\
 &=t^n\MA_\omega(u_t)(B')\\
 &\le t^n\,\Capa_\omega(B') \le t^n\Capa_\omega(B).
 \end{aligned}
\]
Adding the two bounds gives the displayed inequality.

\cref{itm:ch48-mvc:null}.  Let \(t\to\infty\).

\cref{itm:ch48-mvc:outer}.  Choose open \(U_m\supset A\) with
\(\Capa_\omega(U_m)\le1/m\); the Borel set \(B_0:=\bigcap_mU_m\) contains
\(A\) and has \(\Capa_\omega(B_0)=0\), so \(\mu(B_0)=0\) by
\cref{itm:ch48-mvc:null}.  For a pluripolar \(A\), apply
\Cref{prop:ch48-pluripolar-capacity}.

\cref{itm:ch48-mvc:holder}.  Take
\(t=\bigl((n+1)E^\vee_\omega(\mu)/\Capa_\omega(B)\bigr)^{1/(n+1)}\), which
is \(\ge1\) by hypothesis; both terms then equal
\(\Capa_\omega(B)^{1/(n+1)}((n+1)E^\vee_\omega(\mu))^{n/(n+1)}\).
\end{proof}

\begin{remark}[This settles a statement deferred from \Cref{chap:bj-finite-energy}]
\label{rem:ch48-settles-ch47-polar}
\Cref{rem:ch47-polar-scope}\,\cref{itm:ch47-polar:pluripolar} recorded that
a finite-energy measure gives no mass to \emph{any} pluripolar set, not
merely to the polar set of an \(\cE^1\) function, and deferred the proof to
this chapter on the ground that it needs the equivalence of pluripolarity
with capacity zero.  \Cref{prop:ch48-measure-vs-capacity} proves it, and it
is worth noticing that only one half of that equivalence is used, namely
\Cref{prop:ch48-pluripolar-capacity}, which is elementary.  The deep half
--- that a set of vanishing outer capacity is pluripolar --- is not needed
for this application.

Two further comments.  The quantitative form
\cref{itm:ch48-mvc:holder} is a genuine domination of \(\mu\) by a power of
the capacity, with exponent \(1/(n+1)\), and it is the trivially valued
analogue of the hypothesis in \Cref{def:ch29-capacity-domination}, where
over a discretely valued base such a domination has to be \emph{assumed}
in order to run the continuity theory.  Here it is a theorem, and the
reason is the Calabi--Yau theorem: every finite-energy measure is a
Monge--Amp\`ere measure, so the splitting of \(B\) along a sublevel set of
its own potential is available.  Second, the exponent \(1/(n+1)\) is not
claimed to be optimal; for \(n=1\) it gives \(\mu(B)\lesssim
\Capa_\omega(B)^{1/2}\), and \Cref{ex:ch48-p1-capacity} shows that no
exponent larger than \(1\) can hold in general.
\end{remark}

\begin{proposition}
\label{prop:ch48-pluripolar-negligible}
\begin{enumerate}
\item\label{itm:ch48-pn:neg}
      Every pluripolar set is negligible.
\item\label{itm:ch48-pn:ideal}
      A subset of a pluripolar set is pluripolar, and a countable union of
      pluripolar sets is pluripolar.  The same holds for sets of vanishing
      outer capacity and for negligible sets.
\item\label{itm:ch48-pn:height}
      If \(T_\omega(A)=+\infty\) in the sense of
      \eqref{eq:bj-alexander-taylor-height}, then \(A\) is pluripolar.
\end{enumerate}
\end{proposition}

\begin{proof}
\cref{itm:ch48-pn:neg}.  Let \(A\subset\{\phi=-\infty\}\) with
\(\phi\in\PSH(\omega)\), \(\sup\phi=0\), as in
\Cref{lem:ch48-normalize-pole}, and put \(u_j:=\phi/j\) for \(j\ge1\).
Each \(u_j\) is \(\omega\)-psh, being the convex combination
\(\tfrac1j\phi+(1-\tfrac1j)\cdot0\)
(\Cref{thm:bj-basic-closure}\,(3)), the family is increasing and bounded
above by \(0\), and
\[
 \sup_ju_j=\begin{cases}0,&\phi>-\infty,\\ -\infty,&\phi=-\infty.\end{cases}
\]
Its upper semicontinuous regularization is the constant \(0\): it is
\(\le0\) because \(0\) is an upper semicontinuous majorant, and it is
\(\ge0\) because \(\sup_ju_j=0\) on \(X^{\divis}\), which is dense in
\(X^{\An}\) by \Cref{thm:triv-divisorial-density} --- indeed
\(X^{\divis}\subset X^{\lin}\) and \(\phi>-\infty\) on \(X^{\lin}\) by
\eqref{eq:bj-e1-sublinear-growth}, or directly because
\(\phi(v)\ge-T_\omega(v)>-\infty\) there.  Hence
\[
 \Bigl\{\sup_ju_j<\bigl(\sup_ju_j\bigr)^*\Bigr\}
 =\{\phi=-\infty\}\supset A ,
\]
so \(A\) is negligible.

\cref{itm:ch48-pn:ideal}.  Subsets are clear in all three cases.  For a
countable union of pluripolar sets \(A_j\subset\{\phi_j=-\infty\}\),
normalize each \(\phi_j\in\PSH(\omega)\) with \(\sup\phi_j=0\) by
\Cref{lem:ch48-normalize-pole}, and set
\(\phi:=\sum_{j\ge1}2^{-j}\phi_j\).  The partial sums
\(s_N:=\sum_{j\le N}2^{-j}\phi_j\) are \(\omega\)-psh by
\Cref{thm:bj-basic-closure}\,(3) (they are convex combinations of
\(\phi_1,\dots,\phi_N\) and \(0\), the coefficients summing to
\(1-2^{-N}\le1\)), they decrease in \(N\) because \(\phi_j\le0\), and their
limit \(\phi\) satisfies \(\phi(v)\ge-T_\omega(v)>-\infty\) for
\(v\in X^{\lin}\), so \(\phi\) is generically finite; by
\Cref{thm:bj-basic-closure}\,(4), \(\phi\in\PSH(\omega)\).  Finally
\(\phi\le2^{-j}\phi_j\) shows \(\phi=-\infty\) on each \(A_j\).  For sets
of vanishing outer capacity the statement is countable subadditivity,
\Cref{prop:ch48-capacity-elementary}\,\cref{itm:ch48-cap:sub}.  For
negligible sets it follows from \Cref{thm:bj-small-set-equivalence} below,
whose proof does not use the present item.

\cref{itm:ch48-pn:height}.  If \(T_\omega(A)=+\infty\) then for every
\(j\ge1\) there is \(\phi_j\in\PSH(\omega)\) with
\(\sup_{X^{\An}}\phi_j=0\) and \(\sup_A\phi_j\le-4^j\).  Put
\(\phi:=\sum_{j\ge1}2^{-j}\phi_j\), which lies in \(\PSH(\omega)\) by the
argument just given, and note \(\phi(v_{\triv})=0\), so
\(\phi\not\equiv-\infty\).  On \(A\) we have
\(\phi\le2^{-j}\phi_j\le-2^j\) for every \(j\), hence \(\phi=-\infty\) on
\(A\).
\end{proof}

\subsection{The equivalence, and what is imported}
\label{sec:ch48-equivalence}

\begin{theorem}\label{thm:bj-small-set-equivalence}
Assume the envelope property.  For every \(A\subset X^{\An}\), the
following are equivalent:

\[
 A\text{ is pluripolar}
 \quad\Longleftrightarrow\quad
 A\text{ is negligible}
 \quad\Longleftrightarrow\quad
 \Capa^*(A)=0.
\tag{48.12}\label{eq:bj-small-set-equivalence}
\]
\end{theorem}

\begin{proof}
A pole function gives a bounded family whose raw supremum jumps after usc
regularization, so pluripolar implies negligible
\cite[Proposition~5.3]{BJ22}.  Quasicontinuity makes a negligible set have
outer capacity zero.  Finally the Alexander--Taylor inequality below makes
capacity zero force infinite extremal height, which is equivalent to
pluripolarity.  These are \cite[Theorems~13.14 and 13.17]{BJ22}.  In BJ22
the envelope property was available under its historical hypotheses; here
\Crefrange{chap:bj-envelope-statement}{chap:bj-envelope-consequences}
supplies it for unibranch \(X\).

We now go round the cycle, marking each arrow.

\emph{Pluripolar \(\Rightarrow\) negligible.}  Proved, as
\Cref{prop:ch48-pluripolar-negligible}\,\cref{itm:ch48-pn:neg}.  The
construction is the family \((\phi/j)_{j\ge1}\), and the only input is that
\(X^{\divis}\) is dense in \(X^{\An}\) and that a psh function is finite
there.

\emph{Pluripolar \(\Rightarrow\Capa^*=0\).}  Proved, as
\Cref{prop:ch48-pluripolar-capacity}, with the explicit rate \(n/t\).

\emph{Negligible \(\Rightarrow\Capa^*=0\).}  Imported:
\Cref{imp:ch48-imported}\,\cref{itm:ch48-import:negligible}, that is,
\cite[Theorem~13.17]{BJ22}, which assumes the envelope property.  The
mechanism is quasicontinuity, which we have proved
(\Cref{thm:bj-quasicontinuity}), applied to a countable subfamily produced
by Choquet's lemma \Cref{lem:choquet-selection}, together with a
convergence-in-capacity statement for the resulting \emph{increasing}
sequence.  It is the last ingredient that we do not supply:
\Cref{lem:ch48-fundamental} controls decreasing nets and, as
\Cref{rem:ch48-fundamental-transfer} explains, has no increasing
counterpart.

\emph{\(\Capa^*=0\Rightarrow\) pluripolar.}  Imported for general \(A\),
through the lower bound in \eqref{eq:bj-alexander-taylor}, that is,
\cite[Theorem~13.14]{BJ22}: vanishing outer capacity forces
\(T_\omega(A)=+\infty\), and then
\Cref{prop:ch48-pluripolar-negligible}\,\cref{itm:ch48-pn:height} --- which
is proved here --- produces a pole function.  For \emph{compact} \(A\) the
lower bound is proved below, in \Cref{prop:ch48-at-compact}, so for compact
sets the whole implication is available here.

\emph{Negligible \(\Rightarrow\) pluripolar} and the remaining arrows follow
by composing the four above.
\end{proof}

\subsection{The Alexander--Taylor inequality}
\label{sec:ch48-alexander-taylor}

For a set \(A\), define the Alexander--Taylor height

\[
 T_\omega(A):=\sup_{\phi\in\PSH(\omega)}
                   \bigl(\sup_{X^{\An}}\phi-\sup_A\phi\bigr).
\tag{48.13}\label{eq:bj-alexander-taylor-height}
\]

With \(+\infty^{-1}=0\), \cite[Theorem~13.14]{BJ22} states

\[
 \min\{1,T_\omega(A)^{-n}\}
 \leq\Capa_\omega^*(A)
 \leq nT_\omega(A)^{-1}.
\tag{48.14}\label{eq:bj-alexander-taylor}
\]

This quantitative estimate is the bridge between extremal psh functions
and capacity.

For a single point \(v\in X^{\An}\) the height
\eqref{eq:bj-alexander-taylor-height} is the quantity
\eqref{eq:bj-T-valuation} of \Cref{chap:bj-linear-growth}, so
\(T_\omega(\{v\})<\infty\) exactly when \(v\in X^{\lin}\)
(\Cref{prop:bj-linear-growth-potential}); the two notations are consistent
and we use the same symbol.  For a general set,
\(T_\omega(A)=\inf_{v\in A}T_\omega(\{v\})\) is \emph{false}, and
\Cref{ex:ch48-p1-capacity} contains a two-point counterexample; the height
of a set is a genuine minimax and not an infimum of pointwise heights.

\begin{proposition}
\label{prop:ch48-at-upper}
For every \(A\subset X^{\An}\),
\(\Capa^*_\omega(A)\le n\,T_\omega(A)^{-1}\), with the convention
\(n\cdot(+\infty)^{-1}=0\).
\end{proposition}

\begin{proof}
Write \(T:=T_\omega(A)\).  If \(T=0\) there is nothing to prove, the
right-hand side being \(+\infty\).  Suppose \(0<T\le\infty\) and let
\(0<t<T\).  By \eqref{eq:bj-alexander-taylor-height} there is
\(\phi\in\PSH(\omega)\), which we normalize by \(\sup_{X^{\An}}\phi=0\)
using \Cref{thm:bj-basic-closure}\,(1), with \(\sup_A\phi\le-t\); thus
\(A\subset\{\phi\le-t\}\subset\{\phi<-t'\}\) for every \(t'<t\).  The last
set is open, so by Step~1 of the proof of
\Cref{prop:bj-capacity-sublevel},
\[
 \Capa^*_\omega(A)\le\Capa_\omega\bigl(\{\phi<-t'\}\bigr)\le\frac n{t'} .
\]
Letting \(t'\uparrow t\) and then \(t\uparrow T\) gives the claim.
\end{proof}

\begin{proposition}
\label{prop:ch48-at-compact}
Let \(X\) be integral projective unibranch, \(\omega\in\Amp(X)\), and let
\(K\subset X^{\An}\) be compact and nonempty, with
\(T:=T_\omega(K)\).  Then
\[
 \min\{1,T^{-n}\}\ \le\ \Capa_\omega(K)\ \le\ \Capa^*_\omega(K)
 \ \le\ \frac nT .
\]
Consequently:
\begin{enumerate}
\item\label{itm:ch48-atc:zero}
      \(\Capa_\omega(K)=0\ \Longleftrightarrow\ \Capa^*_\omega(K)=0
      \ \Longleftrightarrow\ T=+\infty
      \ \Longleftrightarrow\ K\) is pluripolar;
\item\label{itm:ch48-atc:curve}
      if \(n=1\) and \(T\ge1\), then
      \(\Capa_\omega(K)=\Capa^*_\omega(K)=T^{-1}\) exactly.
\end{enumerate}
\end{proposition}

\begin{proof}
The upper bound is \Cref{prop:ch48-at-upper}, and
\(\Capa_\omega\le\Capa^*_\omega\) is
\Cref{prop:ch48-capacity-elementary}\,\cref{itm:ch48-cap:outer}.  For the
lower bound we may assume \(T<\infty\).

\emph{The extremal function.}  Let
\[
 f(x):=
 \begin{cases}
  0,&x\in K,\\
  +\infty,&x\in X^{\An}\setminus K.
 \end{cases}
\]
Since \(K\) is closed, \(\{f>c\}=X^{\An}\setminus K\) is open for
\(c\ge0\) and equals \(X^{\An}\) for \(c<0\), so \(f\) is lower
semicontinuous with values in \([0,+\infty]\); it is bounded below by
\(0\).  Its competitor family is
\[
 \mathcal F_\omega(f)=\{\psi\in\PSH(\omega)\mid\psi\le0\ \text{on}\ K\},
\]
which is nonempty (it contains the constant \(0\), \(\omega\) being nef) and
bounded above by \(T\): indeed for \(\psi\) in it,
\(\sup_{X^{\An}}\psi\le T+\sup_K\psi\le T\) by
\eqref{eq:bj-alexander-taylor-height}.  Put
\(h:=P_\omega(f)^*=\bigl(\sup\mathcal F_\omega(f)\bigr)^*\), which lies in
\(\PSH(\omega)\) by \Cref{thm:bj-final-envelope}, satisfies \(0\le h\le T\),
and has \(\sup_{X^{\An}}h=T\): the inequality \(\le\) is the bound just
proved, and \(\ge\) follows because for any \(\phi\in\PSH(\omega)\) the
translate \(\phi-\sup_K\phi\) lies in \(\mathcal F_\omega(f)\), whence
\(\sup_{X^{\An}}h\ge\sup_{X^{\An}}\phi-\sup_K\phi\).  In particular
\(h\not\equiv+\infty\).

\emph{Orthogonality.}  By \Cref{thm:bj-obstacle-orthogonality-energy},
applied to the lower semicontinuous obstacle \(f\),
\[
 \int_{X^{\An}}\bigl(f-h\bigr)\,\MA_\omega(h)=0 .
\]
The integrand is \(+\infty\) on \(X^{\An}\setminus K\), because \(h\le T\)
is finite there, and is \(\ge0\) on \(K\).  A nonnegative integrand with
vanishing integral forces
\(\MA_\omega(h)(X^{\An}\setminus K)=0\), that is,
\[
 \MA_\omega(h)(K)=1 .
\]

\emph{The admissible competitor.}  Put \(T':=\max\{T,1\}\ge1\) and
\[
 u:=\frac h{T'}-1 .
\]
Then \(h/T'=\tfrac1{T'}h+(1-\tfrac1{T'})\cdot0\) is \(\omega\)-psh by
\Cref{thm:bj-basic-closure}\,(3), so \(u\in\PSH(\omega)\) by
\Cref{thm:bj-basic-closure}\,(1); and \(0\le h\le T\le T'\) gives
\(-1\le u\le0\), so \(u\in\mathcal U_\omega\).  Expanding
\(\omega+\ddc(h/T')=\tfrac1{T'}(\omega+\ddc h)+(1-\tfrac1{T'})\omega\) by
multilinearity and discarding all but the pure term,
\(\MA_\omega(u)=\MA_\omega(h/T')\ge(T')^{-n}\MA_\omega(h)\).  Therefore
\[
 \Capa_\omega(K)\ \ge\ \MA_\omega(u)(K)\ \ge\ (T')^{-n}\MA_\omega(h)(K)
 =(T')^{-n}=\min\{1,T^{-n}\} .
\]

\cref{itm:ch48-atc:zero}.  If \(\Capa_\omega(K)=0\) then
\(\min\{1,T^{-n}\}=0\), so \(T=+\infty\), and
\Cref{prop:ch48-pluripolar-negligible}\,\cref{itm:ch48-pn:height} makes
\(K\) pluripolar; if \(K\) is pluripolar then
\(\Capa^*_\omega(K)=0\) by \Cref{prop:ch48-pluripolar-capacity}; and
\(\Capa_\omega\le\Capa^*_\omega\) closes the cycle.  If \(T=+\infty\) the
upper bound gives \(\Capa^*_\omega(K)=0\).

\cref{itm:ch48-atc:curve}.  For \(n=1\) and \(T\ge1\) the two bounds are
\(T^{-1}\le\Capa_\omega(K)\le\Capa^*_\omega(K)\le T^{-1}\).
\end{proof}

\begin{warning}{Empty sets and infinite heights}
We put \(\sup_\varnothing\phi=-\infty\), so
\(T_\omega(\varnothing)=+\infty\) and
\(\Capa^*(\varnothing)=0\).  For a nonempty pluripolar set the same infinite
height occurs, but the set is not empty.  No subtraction
\(+\infty-(+\infty)\) is used in \eqref{eq:bj-alexander-taylor-height}:
psh functions have a finite global supremum.
\end{warning}

\subsection{Discharging a step deferred from \texorpdfstring{\Cref{chap:bj-finite-energy}}{Chapter 47}}
\label{sec:ch48-discharge}

\Cref{chap:bj-finite-energy} proved orthogonality for continuous obstacles
(\Cref{thm:ch47-orthogonality-continuous}) and, for lower semicontinuous
obstacles, reduced \Cref{thm:bj-obstacle-orthogonality-energy} to two
statements that it deferred to the present chapter: that the envelopes of
an increasing net of continuous minorants converge strongly, and that the
negligible set on which the raw supremum falls short of its regularization
is null for the limiting finite-energy measure.  Both are now available.

\begin{proposition}
\label{prop:ch48-discharge}
Let \(\omega\in\Amp(X)\) with \(X\) integral projective unibranch, let
\(h\) be lower semicontinuous and bounded below on \(X^{\An}\), and suppose
\(P_\omega(h)^*\not\equiv+\infty\).  Let \((h_j)\) be an upward directed
net of continuous functions with \(\sup_jh_j=h\).  Then:
\begin{enumerate}
\item\label{itm:ch48-dis:strong}
      \(P_\omega(h_j)\in C^0(X^{\An})\cap\cE^1(\omega)\) for every \(j\),
      the net \((P_\omega(h_j))\) is increasing with
      \(\bigl(\sup_jP_\omega(h_j)\bigr)^*=P_\omega(h)^*\in\cE^1(\omega)\),
      and \(P_\omega(h_j)\to P_\omega(h)^*\) strongly; in particular
      \(\MA_\omega(P_\omega(h_j))\to\MA_\omega(P_\omega(h)^*)\) strongly,
      hence weakly;
\item\label{itm:ch48-dis:null}
      the set
      \(N:=\bigl\{\sup_jP_\omega(h_j)<P_\omega(h)^*\bigr\}\) satisfies
      \(\Capa^*_\omega(N)=0\) and is \(\mu\)-null for every
      \(\mu\in\cM^1(\omega)\).
\end{enumerate}
Consequently \eqref{eq:bj-obstacle-contact} and
\eqref{eq:bj-obstacle-orthogonality} hold.
\end{proposition}

\begin{proof}
\cref{itm:ch48-dis:strong}.  Each \(h_j\) is a continuous, hence admissible,
obstacle (\(\omega\) is nef, so constants are competitors, by
\Cref{lem:ch34-psh-nonempty}), and
\Cref{cor:ch45-continuity-envelopes} gives
\(P_\omega(h_j)\in C^0(X^{\An})\cap\PSH(\omega)\); a continuous function on
a compact space is bounded, so
\(P_\omega(h_j)\in\cE^1(\omega)\) by \Cref{thm:bj-e1-basic}.  The net is
increasing by \Cref{prop:bj-final-obstacle-calculus}\,(1), and
\(\sup_jP_\omega(h_j)=\sup\mathcal F_\omega(h)=P_\omega(h)\) with
\(P_\omega(h)^*\in\PSH(\omega)\) by
\Cref{thm:ch45-increasing-obstacles}\,\cref{itm:ch45-inc:sup},
\cref{itm:ch45-inc:star}.  Moreover \(P_\omega(h)^*\) is bounded: it is
bounded below by the bounded function \(P_\omega(h_{j_0})\) for any fixed
\(j_0\), and above by \(\sup_{X^{\An}}P_\omega(h)^*<\infty\); hence
\(P_\omega(h)^*\in\cE^1(\omega)\).  Strong convergence is now the
increasing case of \Cref{prop:bj-strong-monotone}, and strong convergence
of the Monge--Amp\`ere measures is
\Cref{thm:ch48-ma-homeomorphism}\,\cref{itm:ch48-ma:lip}; weak convergence
follows from \Cref{thm:bj-strong-measures}.

\cref{itm:ch48-dis:null}.  The family \((P_\omega(h_j))\subset\PSH(\omega)\)
is bounded above, so \(N\) is negligible in the sense of
\eqref{eq:bj-negligible}, and \Cref{thm:bj-small-set-equivalence} gives
\(\Capa^*_\omega(N)=0\).  That \(\mu(N)=0\) for every
\(\mu\in\cM^1(\omega)\) is \Cref{prop:ch48-measure-vs-capacity}.

For the consequence, pass to the limit in the identity
\[
 \int\bigl(h_j-P_\omega(h_j)\bigr)
   \MA_\omega\bigl(P_\omega(h_j)\bigr)=0
\]
of
\Cref{thm:ch47-orthogonality-continuous}: by \cref{itm:ch48-dis:strong} the
measures converge strongly, by \cref{itm:ch48-dis:null} the exceptional set
is null for the limit measure, and monotone convergence handles the
increasing integrand \(h_j\uparrow h\).
\end{proof}

\begin{remark}[The logical direction, stated once]
\label{rem:ch48-logical-direction}
There is a dependency here that must be displayed, because two of the
statements involved are proved in opposite orders in
\Cref{chap:bj-finite-energy} and in \cite{BJ22}.

\Cref{prop:ch48-at-compact} invokes
\Cref{thm:bj-obstacle-orthogonality-energy}, which is
\cite[Theorem~13.1]{BJ22}; \Cref{prop:ch48-discharge} shows how the two
steps that \Cref{chap:bj-finite-energy} deferred are supplied, and it uses
\Cref{thm:bj-small-set-equivalence}.  The circle is broken because
\Cref{thm:bj-small-set-equivalence} is proved \emph{without}
\Cref{prop:ch48-at-compact}: its two nontrivial arrows are
\Cref{imp:ch48-imported}\,\cref{itm:ch48-import:negligible}, quoted from
\cite{BJ22}, and \Cref{prop:ch48-at-compact} is offered afterwards as an
independent proof of one of them in the compact case.  A reader who wishes
a strictly linear development may take
\Cref{thm:bj-obstacle-orthogonality-energy} as quoted from
\cite[Theorem~13.1]{BJ22} throughout, in which case
\Cref{prop:ch48-discharge} is a remark about the internal economy of this
book and not a link in any chain of implications.

The order in which the statements of this chapter may be read is
\Cref{lem:ch48-cln}, \Cref{prop:bj-capacity-sublevel},
\Cref{lem:ch48-fundamental}, \Cref{thm:bj-quasicontinuity},
\Cref{prop:ch48-pluripolar-capacity},
\Cref{prop:ch48-measure-vs-capacity},
\Cref{prop:ch48-pluripolar-negligible},
\Cref{thm:bj-small-set-equivalence}, \Cref{prop:ch48-at-upper},
\Cref{prop:ch48-at-compact}, \Cref{prop:ch48-discharge}.
\end{remark}

\subsection{What is proved, and what remains open}
\label{sec:ch48-what-is-open}

It is worth setting out, in plain words, exactly where this chapter stands
on the three notions of smallness, because the literature is not uniform
and because the trivially valued and the discretely valued theories differ
here more than anywhere else in
\Crefrange{chap:bj-energy}{chap:bj-stability}.

\emph{What is proved here, unconditionally, for unibranch \(X\).}  A
pluripolar set is negligible
(\Cref{prop:ch48-pluripolar-negligible}\,\cref{itm:ch48-pn:neg}) and has
vanishing outer capacity (\Cref{prop:ch48-pluripolar-capacity}), with the
explicit decay rate \(n/t\) along the sublevel sets of the pole function.
Pluripolar sets form a \(\sigma\)-ideal
(\Cref{prop:ch48-pluripolar-negligible}\,\cref{itm:ch48-pn:ideal}).  A set
of infinite Alexander--Taylor height is pluripolar
(\cref{itm:ch48-pn:height}).  Every finite-energy measure gives no mass to a
set of vanishing capacity, quantitatively
(\Cref{prop:ch48-measure-vs-capacity}).  Psh functions are quasicontinuous
(\Cref{thm:bj-quasicontinuity}), and decreasing nets converge in capacity
(\Cref{prop:ch48-decreasing-capacity}).  For a \emph{compact} set, all
three notions coincide and the Alexander--Taylor inequality holds in both
directions (\Cref{prop:ch48-at-compact}); when \(n=1\) it is an equality.

\emph{What is quoted, with an exact reference.}  Two implications, and only
two.  That a \emph{negligible} set has vanishing outer capacity is
\cite[Theorem~13.17]{BJ22}.  That an arbitrary set of vanishing outer
capacity has infinite Alexander--Taylor height --- the lower bound in
\eqref{eq:bj-alexander-taylor} --- is \cite[Theorem~13.14]{BJ22}.  Both
carry the envelope property as a hypothesis, and both are therefore
available here.

\emph{Why the arguments of this chapter do not reach them.}  For the first,
the mechanism is convergence in capacity of an \emph{increasing} sequence,
and \Cref{lem:ch48-fundamental} controls only decreasing ones; the
asymmetry is explained in \Cref{rem:ch48-fundamental-transfer} and is the
same one that makes
\Cref{imp:ch48-imported}\,\cref{itm:ch48-import:increasing} an import.
For the second, the proof given in \Cref{prop:ch48-at-compact} constructs
the equilibrium potential of a compact set and uses orthogonality to place
its Monge--Amp\`ere mass on that set; for a non-compact \(A\) the obstacle
that equals \(0\) on \(A\) and \(+\infty\) on \(A^c\) is no longer lower
semicontinuous, and passing to an open \(U\supset A\) reverses the
inequality one needs, since \(T_\omega(U)\le T_\omega(A)\) but
\(\Capa_\omega(U)\ge\Capa^*_\omega(A)\) points the wrong way.  What would
close the gap is the identity
\(T_\omega(U)=\inf\{T_\omega(K)\mid K\subset U\ \text{compact}\}\) for open
\(U\), a minimax statement that we have not proved; it is true, because it
follows from \cite[Theorem~13.14]{BJ22} together with
\Cref{prop:ch48-at-compact} and inner regularity, but that is a deduction
from the imported statement and not an independent argument.

\emph{What is genuinely open, here and in the literature.}  Three things,
and none of them is an artefact of our exposition.

First, no quantitative form of ``negligible \(\Rightarrow\) pluripolar'' is
known.  The statement is qualitative: one learns that the exceptional set
of an increasing family is swallowed by a single psh function, but not how
the pole of that function relates to the family.  In the complex theory the
corresponding statement is Bedford and Taylor's \cite{BT82}, and its known
proofs go through Josefson's theorem \cite{Jos78} on the global nature of
local pluripolarity; no analogue of Josefson's theorem is available in the
non-Archimedean setting, and the route taken by \cite{BJ22} is different.

Second, the strong subadditivity of \(\Capa_\omega\), and with it the
capacitability of Borel sets in the classical Choquet sense, is not proved
here.  \cite[Theorem~13.29]{BJ22} asserts that \(\Capa^*_\omega\) is a
Choquet capacity, and we quote it where it is used; but
\Cref{ex:ch48-proper-support-capacity} exhibits a set all of whose compact
subsets have zero capacity and whose outer capacity is \(1\), so no
capacitability statement can hold for arbitrary subsets.  The same warning
over a discretely valued base is
\Cref{rem:ch25-thuillier-vs-bfj}, where strong subadditivity of
\(\Capa_\theta\) is left unproved even for curves, and
\Cref{ex:ch25-subadditivity-strict} shows that ordinary subadditivity is
strict.

Third, the relation between \(\Capa_\omega\) and Thuillier's relative
capacity is not established in either direction, even in dimension one.
\Cref{thm:ch25-thuillier-polar} identifies polar subsets of a Berkovich
curve over a nontrivially valued field with the sets of vanishing outer
relative capacity, and does so without any quasicontinuity; the two
capacities are different set functions, as
\Cref{rem:ch25-thuillier-vs-bfj} records, and over a trivially valued field
his hypotheses --- a connected open subset with nonempty boundary in the
quasimonomial locus, and countably many connected components --- fail in
the strongest possible way, the punctured cone having one component for
each closed point of \(X\).  We therefore do not transport his theorem, and
we do not conjecture a common generalization.  What we do transport is his
discipline about countability, and
\Cref{ex:ch48-proper-support-capacity} is the trivially valued form of his
\cite[Remarque~3.6.12]{Thu05}.

No open-problem environment is used for any of this.  The reader who wants
the statements as problems will find the corresponding entries in
\Cref{chap:major-open-problems}; the purpose of the present subsection is
to make the boundary of what has been proved visible at the point where it
is crossed.

\section{Compactness and convergence criteria}
\label{sec:ch48-compactness}

\subsection{The basic criterion}
\label{sec:ch48-basic-criterion}

\begin{proposition}\label{prop:bj-strong-compactness-criterion}
Let \((\phi_i)\) be a net in \(\mathcal E^1\), normalized by
\(\sup\phi_i=0\), with \(J(\phi_i)\leq C\).  Every weak cluster point lies
in \(\mathcal E^1\).  If the net converges weakly to
\(\phi\in\mathcal E^1\), then it converges strongly if and only if
\[
                         E(\phi_i)\longrightarrow E(\phi).
\tag{48.15}\label{eq:bj-weak-energy-strong-criterion}
\]
In particular, a monotone net in \(\mathcal E^1\) that converges weakly to
a point of \(\mathcal E^1\) converges strongly.
\end{proposition}

\begin{proof}
Normalized psh compactness provides weak cluster subnets.  Lower
semicontinuity of \(J=-E\) on the normalized locus gives the finite-energy
claim.  Criterion \eqref{eq:bj-weak-energy-strong-criterion} is precisely
the characterization of the strong topology in
\Cref{thm:bj-strong-e1-quasimetric}.  Notice that no conclusion is
drawn merely from the energies of weak cluster points: energy is upper
semicontinuous, not weakly continuous.  The last assertion is
\Cref{prop:bj-strong-monotone}.

In detail.  The set
\(\mathcal K_C=\{\psi\in\PSH^{\sup}(\omega)\mid E_\omega(\psi)\ge-C\}\)
contains the net, and it is compact: \(\PSH^{\sup}(\omega)\) is weakly
compact by \Cref{cor:ch45-compactness}, and \(\{E_\omega\ge-C\}\) is weakly
closed because \(E_\omega\) is weakly upper semicontinuous
(\Cref{prop:ch46-energy-usc}\,\cref{itm:ch46-Eusc:props}).  A compact set is
closed in a Hausdorff space
(\Cref{lem:ch48-net-criterion}\,\cref{itm:ch48-net:hausdorff}), so every
weak cluster point of the net lies in \(\mathcal K_C\), hence has
\(E_\omega\ge-C>-\infty\) and lies in \(\cE^1_0(\omega)\); and cluster
points exist, by compactness.  The equivalence
\eqref{eq:bj-weak-energy-strong-criterion} is
\Cref{lem:ch48-net-criterion}\,\cref{itm:ch48-net:criterion}, and the last
assertion is \Cref{prop:bj-strong-monotone}.
\end{proof}

\begin{proposition}
\label{prop:ch48-criteria}
Let \((\phi_i)\) be a net in \(\cE^1(\omega)\) and
\(\phi\in\cE^1(\omega)\).  Each of the following implies
\(\phi_i\to\phi\) strongly, and \cref{itm:ch48-crit:weak},
\cref{itm:ch48-crit:gauge} and \cref{itm:ch48-crit:measure} are also
necessary.
\begin{enumerate}
\item\label{itm:ch48-crit:unif}
      \(\sup_{X^{\divis}}|\phi_i-\phi|\to0\);
\item\label{itm:ch48-crit:mono}
      \((\phi_i)\) is monotone and \(\phi_i\to\phi\) weakly;
\item\label{itm:ch48-crit:weak}
      \(\phi_i\to\phi\) weakly and \(E_\omega(\phi_i)\to E_\omega(\phi)\);
\item\label{itm:ch48-crit:gauge}
      \(\mathbb I(\phi_i,\phi)\to0\);
\item\label{itm:ch48-crit:defect}
      \(J_\phi(\phi_i)\to0\) and \(\sup\phi_i\to\sup\phi\);
\item\label{itm:ch48-crit:measure}
      \(I^\vee\bigl(\MA_\omega(\phi_i),\MA_\omega(\phi)\bigr)\to0\) and
      \(\sup\phi_i\to\sup\phi\);
\item\label{itm:ch48-crit:cauchy}
      \((\phi_i)\) is \(\mathbb I\)-Cauchy and \(\phi\) is one of its weak
      cluster points.
\end{enumerate}
\end{proposition}

\begin{proof}
\cref{itm:ch48-crit:unif} is \Cref{prop:ch48-uniform-strong}.
\cref{itm:ch48-crit:mono} is \Cref{prop:bj-strong-monotone}.
\cref{itm:ch48-crit:weak} is the definition, through
\Cref{lem:ch48-net-criterion}\,\cref{itm:ch48-net:criterion}.
\cref{itm:ch48-crit:gauge} is \Cref{thm:bj-strong-e1-quasimetric}.
\cref{itm:ch48-crit:defect}: by \eqref{eq:bj-I-J-comparison},
\(I(\phi_i,\phi)\le(n+1)J_\phi(\phi_i)\to0\), so
\(\mathbb I(\phi_i,\phi)\to0\).
\cref{itm:ch48-crit:measure}: by the lower bound in the proof of
\Cref{prop:bj-ma-strong-embedding},
\(I(\phi_i,\phi)\le(n+1)C_n\,I^\vee(\MA_\omega(\phi_i),\MA_\omega(\phi))
\to0\).
\cref{itm:ch48-crit:cauchy}: by \Cref{thm:ch48-completeness} the net
converges strongly to some \(\phi'\in\cE^1(\omega)\); strong convergence
implies weak convergence, and the weak topology is Hausdorff, so the weak
cluster point \(\phi\) equals \(\phi'\).

Necessity of \cref{itm:ch48-crit:weak} and \cref{itm:ch48-crit:gauge} is
\Cref{thm:bj-strong-e1-quasimetric}; necessity of
\cref{itm:ch48-crit:measure} follows from
\Cref{thm:ch48-ma-homeomorphism}\,\cref{itm:ch48-ma:lip}.
\end{proof}

\subsection{Strong compactness fails}
\label{sec:ch48-no-strong-compactness}

\begin{proposition}
\label{prop:ch48-no-strong-compactness}
Let \(X=\mathbf P^1_k\) and \(\omega=c_1(\cO(1))\), so \(n=1\) and
\(V=1\).  For every \(C\ge\tfrac12\) the sets
\[
 \mathcal K_C=\{\phi\in\cE^1_0(\omega)\mid J(\phi)\le C\},
 \qquad
 \cM^1_C=\{\mu\in\cM(\omega)\mid E^\vee_\omega(\mu)\le C\}
\]
are weakly compact and are \emph{not} compact for the strong topologies.
Consequently the strong topologies are strictly finer than the weak ones,
on both sides.
\end{proposition}

\begin{proof}
Weak compactness is \Cref{cor:ch45-compactness} together with
\Cref{prop:ch46-energy-usc}\,\cref{itm:ch46-Eusc:props} on the potential
side, and
\Cref{prop:ch47-dual-elementary}\,\cref{itm:ch47-dual:compact} on the
measure side.

For the failure, keep the sequence of
\Cref{ex:ch46-energy-not-weakly-continuous}: distinct closed points
\(p_1,p_2,\dots\in X(k)\), which exist because \(k\) is algebraically
closed hence infinite, and
\(\phi_j:=\max\{\log|\mathfrak m_{p_j}|,-1\}\).  By
\Cref{ex:ch46-p1-energy} with \(d=1\), \(c=1\), we have
\(\sup\phi_j=0\), \(E_\omega(\phi_j)=-\tfrac12\), hence
\(J(\phi_j)=\tfrac12\le C\) and \(\phi_j\in\mathcal K_C\); and
\(\MA_\omega(\phi_j)=\delta_{\ord_{p_j}}\), with
\(E^\vee_\omega(\delta_{\ord_{p_j}})=J(\phi_j)=\tfrac12\le C\) by
\Cref{thm:ch48-ma-homeomorphism}\,\cref{itm:ch48-ma:energy}, so
\(\MA_\omega(\phi_j)\in\cM^1_C\).

Suppose some subnet \((\phi_{j_\alpha})\) converged strongly to
\(\phi\in\mathcal K_C\).  Strong convergence implies weak convergence, and
\Cref{ex:ch46-energy-not-weakly-continuous} shows that the whole sequence
converges weakly to \(0\); the weak topology being Hausdorff, \(\phi=0\).
But then \(E_\omega(\phi_{j_\alpha})\to E_\omega(0)=0\), contradicting
\(E_\omega(\phi_j)=-\tfrac12\) for every \(j\).  The same argument on the
measure side, using
\(\MA_\omega(\phi_j)\to\delta_{v_{\triv}}=\MA_\omega(0)\) weakly and
\(E^\vee_\omega(\delta_{v_{\triv}})=0\ne\tfrac12\), shows that
\(\cM^1_C\) is not strongly compact.
\end{proof}

\begin{remark}[The shape of the failure]
\label{rem:ch48-failure-shape}
\Cref{prop:ch48-no-strong-compactness} is the exact analogue of
\Cref{prop:ch26-no-strong-compactness} over a discretely valued base, and
the two proofs are the same proof; the witnessing sequences correspond
under the dictionary of \Cref{sec:ch46-part3}, the points \(p_j\) being
replaced by closed points of the special fibre.  The failure is not a
pathology of an exotic class: every member of the sequence is a PL psh
function attached to the blow-up of a point, and they all have the same
energy.  What the strong topology detects is that the mass of
\(\MA_\omega(\phi_j)\) sits at \(\ord_{p_j}\), a point that moves without
converging inside the divisorial locus, while the weak limit of the
measures places it at \(v_{\triv}\); the energy records the discrepancy and
the weak topology does not.

Two positive statements survive, and they are what the variational method
actually uses.  Energy sublevels are weakly compact, so maximizers exist
(\Cref{thm:ch47-existence-admissible}); and \(\mathbb I\)-Cauchy nets do
converge (\Cref{thm:ch48-completeness}), so almost-maximizing nets, which
are Cauchy by \Cref{prop:ch47-maximizing-cauchy}, converge strongly.  The
absence of strong compactness costs nothing, because in every application
the Cauchy property is available.
\end{remark}

\section{Worked examples}
\label{sec:ch48-examples}

Throughout this section \(X=\mathbf P^1_k\) with \(k\) algebraically closed
and trivially valued, and \(\omega=c_1(\cO(1))\), so that \(n=1\) and
\(V=1\).  We use the description of \(X^{\An}\) as a cone with vertex
\(v_{\triv}\) and one ray
\(R_p=\{t\,\ord_p\mid t\in[0,\infty)\}\cup\{v_{p,\triv}\}\) for each closed
point \(p\in X(k)\), and the curve criterion of
\Cref{ex:ch34-psh-check}: a function \(\phi\) that vanishes off finitely
many rays and satisfies \(\phi(t\,\ord_p)=h_p(t)\) with \(h_p\) convex,
nonincreasing and \(h_p(0)=0\) is \(\omega\)-psh precisely when
\(\sum_p a_p\le1\), where \(a_p:=-h_p'(0^+)\); and then, by
\eqref{eq:bj-curve-laplacian},
\[
 \MA^{\mathrm u}_\omega(\phi)
 =\Bigl(1-\sum_pa_p\Bigr)\delta_{v_{\triv}}
   +\sum_p h_p''\big|_{R_p},
\]
the second derivatives being taken in the sense of measures on the open
ray, with any residual mass \(-\lim_{t\to\infty}h_p'(t)\) placed at the
endpoint \(v_{p,\triv}\).  Since \(V=1\), the unnormalized and the
probability measures coincide.

\subsection{Weak convergence without strong convergence, with all the numbers}
\label{sec:ch48-ex-weak}

\begin{example}[The motivating sequence of \Cref{chap:bj-energy}, completed]
\label{ex:ch48-weak-not-strong}
Choose pairwise distinct closed points \(p_1,p_2,\dots\in X(k)\), which
exist because \(k\) is algebraically closed hence infinite, and set
\[
 \phi_j:=\max\{\log|\mathfrak m_{p_j}|,-1\},
 \qquad\text{so}\qquad
 \phi_j(t\,\ord_{p_j})=-\min\{t,1\},\quad
 \phi_j=0\ \text{off}\ R_{p_j}.
\]
This is the sequence of \Cref{ex:ch46-energy-not-weakly-continuous}, which
proved that \(E_\omega\) is not weakly continuous and which was there
declared to be the motivation for the present chapter.  We now compute
every quantity introduced above.

\emph{Potential-side numbers.}  By \Cref{ex:ch46-p1-energy} with \(d=1\),
\(c=1\),
\[
 \begin{aligned}
 \sup_{X^{\An}}\phi_j&=0,&
 E_\omega(\phi_j)&=-\tfrac12,\\
 J(\phi_j)&=\sup\phi_j-E_\omega(\phi_j)=\tfrac12,&
 \MA_\omega(\phi_j)&=\delta_{\ord_{p_j}} .
 \end{aligned}
\]
Since \(n=1\), \Cref{prop:ch47-curve-I-norm} gives
\(J_\psi(\phi)=\tfrac12I(\phi,\psi)\), so
\[
 I(\phi_j,0)=2J(\phi_j)=1,
 \qquad
 \mathbb I(\phi_j,0)=I(\phi_j,0)+|0-0|=1 .
\]
For \(j\ne l\), the difference \(\phi_j-\phi_l\) equals \(-1\) at
\(\ord_{p_j}\), equals \(+1\) at \(\ord_{p_l}\), and
\(\MA_\omega(\phi_l)-\MA_\omega(\phi_j)
=\delta_{\ord_{p_l}}-\delta_{\ord_{p_j}}\), whence
\[
 I(\phi_j,\phi_l)=\int(\phi_j-\phi_l)\,
   \rd\bigl(\MA_\omega(\phi_l)-\MA_\omega(\phi_j)\bigr)
 =1-(-1)=2 ,
 \qquad j\ne l .
\]
So the sequence is uniformly \(2\)-separated for \(I\).  As a check, the
quasi-triangle inequality \eqref{eq:bj-I-quasitriangle} with the sharp
constant \(C_1=2\) of \Cref{prop:ch47-curve-I-norm} reads
\(2=I(\phi_j,\phi_l)\le2\bigl(I(\phi_j,0)+I(0,\phi_l)\bigr)=4\), and is not
attained here.

\emph{Measure-side numbers.}  By
\Cref{thm:ch48-ma-homeomorphism}\,\cref{itm:ch48-ma:energy}, which for
\(n=1\) is an equality,
\[
 E^\vee_\omega(\delta_{\ord_{p_j}})=J(\phi_j)=\tfrac12,
 \qquad
 E^\vee_\omega(\delta_{v_{\triv}})=J(0)=0 .
\]
This agrees with \eqref{eq:bj-dirac-energy-T}: for \(n=1\) both bounds
there equal \(\tfrac12T_\omega(v)\), and \(T_\omega(\ord_p)=1\) by the
computation of \Cref{ex:ch48-p1-capacity} below.  By
\Cref{prop:ch48-curve-Idual},
\[
 I^\vee\bigl(\delta_{\ord_{p_j}},\delta_{v_{\triv}}\bigr)
 =\tfrac14I(\phi_j,0)=\tfrac14,
 \qquad
 I^\vee\bigl(\delta_{\ord_{p_j}},\delta_{\ord_{p_l}}\bigr)
 =\tfrac14I(\phi_j,\phi_l)=\tfrac12\quad(j\ne l).
\]

\emph{What converges and what does not.}
\Cref{ex:ch46-energy-not-weakly-continuous} showed that
\(\phi_j\to0\) pointwise on all of \(X^{\An}\), hence weakly, and that
\(\MA_\omega(\phi_j)=\delta_{\ord_{p_j}}\to\delta_{v_{\triv}}
=\MA_\omega(0)\) weakly.  Nevertheless
\[
 \mathbb I(\phi_j,0)=1\not\to0,
 \qquad
 I^\vee\bigl(\MA_\omega(\phi_j),\MA_\omega(0)\bigr)=\tfrac14\not\to0,
\]
so neither side converges strongly, and by
\Cref{lem:ch48-net-criterion}\,\cref{itm:ch48-net:criterion} the obstruction
is exactly the energy: \(E_\omega(\phi_j)=-\tfrac12\) does not converge to
\(E_\omega(0)=0\), the loss being \(\tfrac12\) at every index.

\emph{Four consequences.}  (i) The strong topology on \(\cE^1(\omega)\) is
strictly finer than the weak one, and likewise on \(\cM^1(\omega)\); this
is \Cref{prop:ch48-no-strong-compactness}.  (ii) The bijection of
\Cref{thm:ch48-ma-homeomorphism}\,\cref{itm:ch48-ma:homeo} is not a
homeomorphism for the weak topologies, although here it maps a weakly
convergent sequence to a weakly convergent one:
weak convergence on both sides coexists with a uniform positive distance on
both sides.  (iii) The sequence has no strongly convergent subnet, being
\(2\)-separated for \(I\); so \(\{J\le\tfrac12\}\) and
\(\cM^1_{1/2}\) are not strongly compact.  (iv) The failure has nothing to
do with capacity: by \Cref{ex:ch48-p1-capacity} below,
\(\Capa_\omega(\{\ord_{p_j}\})=1\) for every \(j\), so each
Monge--Amp\`ere measure sits on a set of \emph{full} capacity, and the
sets do not shrink.  What moves is the location of the mass, and the energy
is the functional that sees it.

\emph{Comparison with \Cref{part:bfj}.}  The same computation over a
discretely valued base is \Cref{ex:ch26-usc-strict}, and the numbers agree
after the conversion of \Cref{warn:ch48-normalization}: there
\(E_\omega(\tau_N)=-\tfrac12\) with \(V=1\) as well, and
\Cref{prop:ch26-no-strong-compactness} draws the same conclusion.  The
numerical correspondence in this example is literal, the points \(p_j\)
being replaced by closed points of
the special fibre, as \Cref{rem:ch46-strong-topology-motivation} predicted.
\end{example}

\subsection{Capacity on a \texorpdfstring{\(\mathbf P^1\)}{P1} skeleton}
\label{sec:ch48-ex-capacity}

\begin{example}[Capacities on the cone over \(X(k)\), computed exactly]
\label{ex:ch48-p1-capacity}
Keep \(X=\mathbf P^1\), \(\omega=c_1(\cO(1))\), \(V=1\).  We compute
\(T_\omega\) and \(\Capa_\omega\) for the subsets of \(X^{\An}\) built from
finitely many rays.

\emph{Heights.}  Let \(p_1,\dots,p_m\) be distinct closed points and
\(t>0\), and put \(A_m(t):=\{t\,\ord_{p_1},\dots,t\,\ord_{p_m}\}\).  For
\(\phi\in\PSH(\omega)\) normalized by \(\sup\phi=0\), write
\(h_i(s):=\phi(s\,\ord_{p_i})\), a convex nonincreasing function of \(s\)
with \(h_i(0)=0\) by \Cref{ex:ch34-psh-check} and
\Cref{lem:ch47-three-normalizations}\,\cref{itm:ch47-norm:max}, and put
\(a_i:=-h_i'(0^+)\in[0,1]\).  Then \(\sum_{i=1}^ma_i\le1\): for
\(\varepsilon>0\) the closed tails \(\{s\ge\varepsilon\}\subset R_{p_i}\)
are pairwise disjoint and \(\MA^{\mathrm u}_\omega(\phi)\) gives the
\(i\)-th of them the mass \(-h_i'(\varepsilon)\), by
\eqref{eq:bj-curve-laplacian}, while the total mass is \(1\); letting
\(\varepsilon\downarrow0\) and using that \(-h_i'(\varepsilon)\uparrow a_i\)
gives the bound.  By convexity,
\(\phi(t\,\ord_{p_i})=h_i(t)\ge h_i(0)+h_i'(0^+)t=-a_it\).  Hence
\[
 -\sup_{A_m(t)}\phi=\min_i\bigl(-\phi(t\,\ord_{p_i})\bigr)
 \le t\min_ia_i\le\frac tm ,
\]
with equality for the function with \(a_1=\dots=a_m=1/m\) and
\(h_i(s)=\max\{-s/m,-t/m\}\).  Therefore
\[
 T_\omega\bigl(A_m(t)\bigr)=\frac tm,
 \qquad\text{in particular}\qquad
 T_\omega\bigl(\{t\,\ord_p\}\bigr)=t,\quad T_\omega(\{\ord_p\})=1 .
\]
Note that \(T_\omega(A_2(t))=t/2\) while
\(T_\omega(\{t\ord_{p_1}\})=T_\omega(\{t\ord_{p_2}\})=t\); the height of a
set is strictly smaller than the infimum of the heights of its points, so
the minimax in \eqref{eq:bj-alexander-taylor-height} is genuine.

\emph{Capacities.}  Each \(A_m(t)\) is compact, so
\Cref{prop:ch48-at-compact}\,\cref{itm:ch48-atc:curve} applies and, since
\(n=1\),
\[
 \Capa_\omega\bigl(A_m(t)\bigr)
 =\Capa^*_\omega\bigl(A_m(t)\bigr)
 =\min\Bigl\{1,\frac mt\Bigr\} .
\]
The maximizing admissible potentials are explicit.  For \(t\ge m\) take
\(u(s\,\ord_{p_i}):=\max\{-s/t,-1\}\) and \(u=0\) off the \(m\) rays; then
\(\sum_ia_i=m/t\le1\), \(-1\le u\le0\), and
\[
 \MA_\omega(u)=\Bigl(1-\frac mt\Bigr)\delta_{v_{\triv}}
   +\frac1t\sum_{i=1}^m\delta_{t\,\ord_{p_i}},
\]
of mass \(1\), giving \(\MA_\omega(u)(A_m(t))=m/t\).  For \(t\le m\) take
instead \(u(s\,\ord_{p_i}):=\max\{-s/m,-t/m\}\), so that
\(\sum_ia_i=1\), \(-1\le-t/m\le u\le0\), and
\(\MA_\omega(u)=\frac1m\sum_i\delta_{t\,\ord_{p_i}}\), of mass \(1\)
carried entirely by \(A_m(t)\); this gives \(\Capa_\omega(A_m(t))=1\).

\emph{Subadditivity, strict and not.}  With \(m=2\),
\[
 \Capa_\omega\bigl(A_2(t)\bigr)=\min\Bigl\{1,\frac2t\Bigr\},
 \qquad
 \Capa_\omega\bigl(\{t\ord_{p_1}\}\bigr)
 +\Capa_\omega\bigl(\{t\ord_{p_2}\}\bigr)=2\min\Bigl\{1,\frac1t\Bigr\}.
\]
For \(t\ge1\) the two sides agree and subadditivity is an equality; for
\(t<1\) the left side is \(1\) and the right side is \(2\), so
subadditivity is strict, with a loss of \(1\).  This is the exact analogue
of \Cref{ex:ch25-subadditivity-strict}, where the same phenomenon occurs on
a segment of length four with a loss of \(1\).

\emph{Tails of a ray, and the sharpness of \eqref{eq:bj-cap-sublevel-bound}.}
Let \(t>0\) and let
\(R_p^{\ge t}:=\{s\,\ord_p\mid s\ge t\}\cup\{v_{p,\triv}\}\), a compact
set.  A normalized \(\phi\) is nonincreasing along \(R_p\), so
\(\sup_{R_p^{\ge t}}\phi=\phi(t\,\ord_p)\) and
\(T_\omega(R_p^{\ge t})=T_\omega(\{t\ord_p\})=t\); hence
\[
 \Capa_\omega\bigl(R_p^{\ge t}\bigr)=\min\Bigl\{1,\frac1t\Bigr\} .
\]
Now let \(\phi_0:=\log|\mathfrak m_p|\), so \(\phi_0(s\,\ord_p)=-s\) and
\(\phi_0=0\) off \(R_p\).  Then
\(\{\phi_0\le-t\}=R_p^{\ge t}\) and
\[
 \Capa_\omega\bigl(\{\phi_0\le-t\}\bigr)=\min\Bigl\{1,\frac1t\Bigr\}
 =\min\Bigl\{1,\frac nt\Bigr\},
\]
so the bound \eqref{eq:bj-cap-sublevel-bound} is attained exactly, for
every \(t\ge1\).  The same is true over a discretely valued base, where
\eqref{eq:ch25-p1-capacity} gives the identical formula; the two
computations agree because in both cases the object is a ray carrying a
linear potential and both chapters normalize \(\ddc\) by the sum of
outgoing slopes.

\emph{The energy criterion, tested twice.}  The potential \(\phi_0\) has
\(E_\omega(\phi_0)=-\infty\) by \Cref{ex:ch46-e1-curve}, and indeed
\[
 \int_1^\infty t^n\,\Capa_\omega(\{\phi_0\le-t\})\,\rd t
 =\int_1^\infty t\cdot\frac1t\,\rd t=+\infty ,
\]
so \eqref{eq:bj-cap-energy-criterion} gives no information, as it must.
Now take instead
\[
 \phi_1(s\,\ord_p):=-4\bigl((1+s)^{1/4}-1\bigr),\qquad
 \phi_1=0\ \text{off}\ R_p .
\]
Its slope is \(h'(s)=-(1+s)^{-3/4}\), so \(h\) is convex, nonincreasing,
and \(a_p=1\le1\): thus \(\phi_1\in\PSH(\omega)\), with
\(\sup\phi_1=0\) and \(\phi_1(v_{p,\triv})=-\infty\).  By
\Cref{prop:ch46-dirichlet-vertical}\,\cref{itm:ch46-dir:energy},
\[
 E_\omega(\phi_1)=-\tfrac12\int_0^\infty|h'(s)|^2\,\rd s
 =-\tfrac12\int_0^\infty(1+s)^{-3/2}\,\rd s=-1,
 \qquad J(\phi_1)=1,
\]
so \(\phi_1\in\cE^1(\omega)\).  Its Monge--Amp\`ere measure is
\(\MA_\omega(\phi_1)=h''(s)\,\rd s=\tfrac34(1+s)^{-7/4}\rd s\) on the ray,
of total mass \(\tfrac34\cdot\tfrac43=1\), with no atom at
\(v_{\triv}\) and none at \(v_{p,\triv}\); and
\[
 \int\phi_1\,\MA_\omega(\phi_1)
 =-3\int_0^\infty\bigl[(1+s)^{-3/2}-(1+s)^{-7/4}\bigr]\rd s
 =-3\bigl(2-\tfrac43\bigr)=-2,
\]
so that \(E^\vee_\omega(\MA_\omega(\phi_1))
=E_\omega(\phi_1)-\int\phi_1\MA_\omega(\phi_1)=-1+2=1=J(\phi_1)\), in
agreement with \Cref{thm:ch48-ma-homeomorphism}\,\cref{itm:ch48-ma:energy}
for \(n=1\).  Finally
\(\{\phi_1\le-s\}=R_p^{\ge t_s}\) with \(t_s=(1+s/4)^4-1\), so
\[
 \Capa_\omega\bigl(\{\phi_1\le-s\}\bigr)
 =\min\Bigl\{1,\frac1{(1+s/4)^4-1}\Bigr\}\le\frac{256}{s^4}
 \quad(s\ge1),
\]
and
\(\int_1^\infty s\,\Capa_\omega(\{\phi_1\le-s\})\rd s\le
\int_1^\infty256\,s^{-3}\rd s=128\).  Thus
\eqref{eq:bj-cap-energy-criterion} gives \(J(\phi_1)\le129\), whereas the
true value is \(J(\phi_1)=1\).  The criterion is therefore a very lossy
numerical bound and a sharp qualitative one: it converges exactly in the
finite-energy case and diverges exactly in the infinite-energy one, for
this one-parameter family.
\end{example}

\subsection{Pluripolar sets, countability, and the failure of capacitability}
\label{sec:ch48-ex-pluripolar}

\begin{example}[Proper-support points: countable versus uncountable]
\label{ex:ch48-proper-support-capacity}
Keep \(X=\mathbf P^1\), \(\omega=c_1(\cO(1))\), \(V=1\).  For
\(S\subset X(k)\) put
\[
 B_S:=\{v_{p,\triv}\mid p\in S\}\subset X^{\An}\setminus X^{\lin},
\]
the set of proper-support points attached to \(S\).  The behaviour of
\(B_S\) depends on the cardinality of \(S\) in the sharpest possible way.

\emph{(1) A single point is pluripolar, with capacity \(0\).}  The function
\(\phi_0=\log|\mathfrak m_p|\) of \Cref{ex:ch48-p1-capacity} lies in
\(\PSH(\omega)\), has \(\sup\phi_0=0\), and satisfies
\(\phi_0(v_{p,\triv})=\lim_{t\to\infty}(-t)=-\infty\).  Hence
\(\{v_{p,\triv}\}\) is pluripolar, and by
\Cref{prop:ch48-pluripolar-capacity},
\[
 \Capa^*_\omega\bigl(\{v_{p,\triv}\}\bigr)
 \le\inf_{t>0}\Capa_\omega\bigl(\{\phi_0<-t\}\bigr)
 =\inf_{t>0}\min\Bigl\{1,\frac1t\Bigr\}=0 .
\]
Equivalently \(T_\omega(v_{p,\triv})=+\infty\), which is
\Cref{prop:bj-linear-growth-potential}: proper-support points do not have
linear growth.

\emph{(2) A countable set is pluripolar, with capacity \(0\).}  Let
\(S=\{p_1,p_2,\dots\}\) be countable and put
\(\phi:=\sum_{j\ge1}2^{-j}\log|\mathfrak m_{p_j}|\), which lies in
\(\PSH(\omega)\) by
\Cref{prop:ch48-pluripolar-negligible}\,\cref{itm:ch48-pn:ideal} and equals
\(-\infty\) at every \(v_{p_j,\triv}\); so \(B_S\) is pluripolar and
\(\Capa^*_\omega(B_S)=0\).  The vanishing can also be seen directly, and
the computation is instructive.  Fix \(M\ge1\) and let
\[
 U_M:=\bigcup_{j\ge1}\bigl\{s\,\ord_{p_j}\mid s>M4^j\bigr\}
        \cup\{v_{p_j,\triv}\},
\]
an open neighbourhood of \(B_S\).  Let \(u\in\mathcal U_\omega\), with
initial slopes \(a_j\) on \(R_{p_j}\) and \(h_j:=u|_{R_{p_j}}\).  Since
\(h_j\) is convex with \(h_j(0)=0\) and \(h_j\ge-1\), the mean value
estimate \(h_j(N)-h_j(0)\le N\,h_j'(N)\) gives
\(-1\le h_j(N)\le N h_j'(N)\), that is,
\(-h_j'(N)\le1/N\); and by \eqref{eq:bj-curve-laplacian} the mass that
\(\MA_\omega(u)\) puts on the closed tail beyond \(N\) is exactly
\(-h_j'(N)\).  With \(N=M4^j\),
\[
 \MA_\omega(u)(U_M)\le\sum_{j\ge1}\frac1{M4^j}=\frac1{3M},
 \qquad\text{so}\qquad
 \Capa_\omega(U_M)\le\frac1{3M}\longrightarrow0 .
\]
The convergent series \(\sum_j4^{-j}\) is the whole point: a countable
family of tails can be made to have small total capacity by pushing the
\(j\)-th tail out fast enough.

\emph{(3) An uncountable set is not pluripolar.}  Let \(S\) be uncountable
and suppose \(\phi\in\PSH(\omega)\) with \(\sup\phi=0\) satisfies
\(\phi=-\infty\) on \(B_S\).  As in \Cref{ex:ch48-p1-capacity}, write
\(h_p:=\phi|_{R_p}\) and \(a_p:=-h_p'(0^+)\), so that
\(\sum_{p\in F}a_p\le1\) for every finite \(F\subset X(k)\).  If
\(a_p=0\) then \(h_p'\ge h_p'(0^+)=0\) by convexity, so \(h_p\) is
nondecreasing; being also nonincreasing with \(h_p(0)=0\), it vanishes
identically and \(\phi(v_{p,\triv})=0\ne-\infty\).  Hence \(a_p>0\) for
every \(p\in S\).  But an uncountable family of strictly positive reals
cannot have all its finite partial sums bounded by \(1\): the sets
\(\{p\in S\mid a_p>1/m\}\) increase to \(S\), so one of them is infinite,
and a finite subset of it of cardinality \(>m\) already violates the
bound.  This contradiction shows that \(B_S\) is not pluripolar.

\emph{(4) Its inner capacity is \(0\) and its outer capacity is \(1\).}
Every \(u\in\mathcal U_\omega\) is bounded, hence lies in
\(\cE^1(\omega)\) by \Cref{thm:bj-e1-basic}, so \(\MA_\omega(u)\) is
carried by \(X^{\lin}\) by \Cref{prop:ch46-no-mass-outside-lin}; since
\(B_S\cap X^{\lin}=\varnothing\), every compact \(K\subset B_S\) has
\(\MA_\omega(u)(K)=0\) and therefore
\[
 \Capa_\omega(K)=0\qquad\text{for every compact }K\subset B_S .
\]
On the other hand, let \(U\) be \emph{any} open set containing \(B_S\).
For \(p\in S\) the points \(s\,\ord_p\) converge to \(v_{p,\triv}\) as
\(s\to\infty\), so there is \(N_p\) with
\(\{s\,\ord_p\mid s>N_p\}\subset U\).  The sets
\(S_m:=\{p\in S\mid N_p<m\}\) increase to \(S\), which is uncountable, so
some \(S_{m_0}\) is infinite.  Given \(m\ge m_0\), choose distinct
\(q_1,\dots,q_m\in S_{m_0}\) and put
\[
 u(s\,\ord_{q_i}):=\max\Bigl\{-\frac sm,-\frac{m_0}m\Bigr\},
 \qquad u:=0\ \text{off}\ R_{q_1}\cup\dots\cup R_{q_m}.
\]
Then \(\sum_ia_i=m\cdot\tfrac1m=1\) and
\(-1\le-\tfrac{m_0}m\le u\le0\), so \(u\in\mathcal U_\omega\), and
\[
 \MA_\omega(u)=\frac1m\sum_{i=1}^m\delta_{m_0\ord_{q_i}} ,
\]
a probability measure carried by the \(m\) points \(m_0\ord_{q_i}\), each of
which lies in \(U\) because \(m_0>N_{q_i}\).  Hence
\(\MA_\omega(u)(U)=1\) and \(\Capa_\omega(U)=1\).  As \(U\) was arbitrary,
\[
 \Capa^*_\omega(B_S)=1 .
\]

\emph{(5) What this shows.}  The set \(B_S\) with \(S\) uncountable is
\emph{not capacitable}: the supremum of \(\Capa_\omega(K)\) over compact
\(K\subset B_S\) is \(0\), while \(\Capa^*_\omega(B_S)=1\).  Consequently
\(B_S\) is not \(K\)-analytic, and the Choquet capacitability statement
\cite[Theorem~13.29]{BJ22} does not apply to it; there is no contradiction,
only a reminder that capacitability is a hypothesis on the set.  It also
shows that \Cref{thm:bj-small-set-equivalence} is not vacuous in the
direction one might expect to be easy: \(B_S\) has all of its compact
subsets of capacity zero, all of its points pluripolar, and every countable
subset pluripolar, and yet it is not pluripolar and its outer capacity is
as large as possible.

\emph{(6) The exact analogue in Thuillier's theory.}  This is the trivially
valued form of \cite[Remarque~3.6.12]{Thu05}, recorded in
\Cref{rem:ch25-thuillier-vs-bfj}: over a residue field of cardinality
greater than \(\aleph_0\), the set of type~\(\mathrm I\) points reducing to
the distinct residues of the Gauss point is closed, lies in the
\(-\infty\) locus of a subharmonic function on each connected component of
the punctured space, and has infinite outer relative capacity.  In both
settings the mechanism is the same: the definition of ``polar'' or
``pluripolar'' asks for \emph{one} function, and a single psh function has
only countably many rays along which it can decrease, because the sum of
its initial slopes is bounded by the degree.  The countability clause in
\cite[Th\'eor\`eme~3.6.11]{Thu05} is thus not a technical convenience, and
neither is the appearance of \(\sigma\)-ideals in
\Cref{prop:ch48-pluripolar-negligible}\,\cref{itm:ch48-pn:ideal} and in
\Cref{prop:polar-sigma-ideal}.
\end{example}

\begin{example}[A pluripolar set of positive outer capacity does not exist,
                and why that is a theorem]
\label{ex:ch48-pluripolar-cap-zero}
It is worth contrasting \Cref{ex:ch48-proper-support-capacity}\,(4) with
\Cref{prop:ch48-pluripolar-capacity}.  The latter says that a pluripolar
set always has vanishing \emph{outer} capacity, and its proof is three
lines: the pole function's open sublevel sets form a neighbourhood basis
of small capacity.  The former exhibits a set whose \emph{inner} capacity
vanishes and whose outer capacity is \(1\); by
\Cref{prop:ch48-pluripolar-capacity} that set cannot be pluripolar, and
\Cref{ex:ch48-proper-support-capacity}\,(3) confirms it directly.  The two
computations are consistent, and together they show that the passage from
inner to outer capacity is exactly where the content lies.

For a numerical illustration, take \(S\) countable and compare the two
neighbourhood systems.  Along \(U_M\) of
\Cref{ex:ch48-proper-support-capacity}\,(2) the capacity is at most
\(1/(3M)\); if instead one takes the neighbourhoods
\(U'_M:=\bigcup_j\{s\ord_{p_j}\mid s>M\}\cup\{v_{p_j,\triv}\}\), with the
\emph{same} cut-off \(M\) on every ray, then, putting
\(m:=\lceil M\rceil+1\) and
\[
 u(s\,\ord_{p_i}):=\max\Bigl\{-\frac sm,-\frac{M+1}m\Bigr\}
 \quad(1\le i\le m),\qquad u:=0\ \text{elsewhere},
\]
one has \(\sum_ia_i=1\) and \(-1\le-(M+1)/m\le u\le0\), so
\(u\in\mathcal U_\omega\), while
\(\MA_\omega(u)=\frac1m\sum_i\delta_{(M+1)\ord_{p_i}}\) is a probability
measure carried by \(U'_M\).  Hence \(\Capa_\omega(U'_M)=1\) for every
\(M\), no matter how large.  The capacity of a neighbourhood
of a countable pluripolar set therefore depends entirely on how fast the
cut-offs are allowed to grow, and it is the freedom to let them grow ray by
ray that a countable index set provides and an uncountable one does not.
\end{example}

\section{Scope and hypotheses}
\label{sec:ch48-scope}

\subsection{Where each hypothesis is used}
\label{sec:ch48-hypothesis-use}

The standing hypotheses of \Cref{sec:ch48-hypotheses} are: \(k\)
algebraically closed and trivially valued; \(X\) integral projective of
dimension \(n\) and unibranch; \(\omega\) ample; the probability
normalization.  Each is used, and it is worth recording where.

\emph{Ampleness of \(\omega\).}  Used for the existence of bounded psh
functions --- the constants are \(\omega\)-psh
(\Cref{lem:ch34-psh-nonempty}), without which
\(\mathcal U_\omega=\varnothing\) and \eqref{eq:bj-bt-capacity} would be
identically \(0\).  Used again for continuity of envelopes of continuous
obstacles (\Cref{cor:ch45-continuity-envelopes}), hence for
\Cref{thm:ch45-continuous-regularization} and thus for
\Cref{thm:bj-quasicontinuity}.  For a merely nef class the regularizing net
of \Cref{thm:ch45-continuous-regularization} still exists but need not
consist of continuous functions, and the proof of quasicontinuity given
here breaks at the point where \(\{u_g-\phi>\delta\}\) is asserted to be
open.  Nothing in \Cref{sec:ch48-potentials,sec:ch48-measures,sec:ch48-ma}
needs more than nefness together with \(V>0\), but we have not stated those
sections in that generality.

\emph{Unibranchness.}  Used only through the envelope property, and there
in four ways, listed in \Cref{imp:ch48-envelope}: the obstacle envelope is
an operator, \(\PSH^{\sup}(\omega)\) is compact, the continuous
regularization exists, and orthogonality holds.  Compactness enters
\Cref{thm:ch48-completeness} and
\Cref{prop:bj-strong-compactness-criterion}; the operator property enters
\Cref{prop:ch48-at-compact} and the discharge
\Cref{prop:ch48-discharge}; regularization enters
\Cref{thm:bj-quasicontinuity}.  By
\Cref{thm:bj-final-unibranch-characterization} the hypothesis is not
removable: on a variety that is not unibranch,
\(\PSH^{\sup}(\omega)\) fails to be compact for every ample \(\omega\)
(\Cref{rem:ch45-compactness-equivalent}), and with it every statement of
this chapter that rests on compactness.

\emph{Integrality.}  Used for the uniqueness clause of
\Cref{thm:ch47-uniqueness}, hence for the injectivity in
\Cref{thm:ch48-ma-homeomorphism}\,\cref{itm:ch48-ma:homeo}.  For an
equidimensional reducible \(X\) the fibres of \(\MA_\omega\) are the
locally constant functions rather than the constants, and the quotient in
\eqref{eq:bj-ma-bilipschitz} must be taken accordingly.

\emph{Algebraic closedness of \(k\).}  Used only to know that \(X(k)\) is
infinite, in \Cref{ex:ch46-energy-not-weakly-continuous} and hence in
\Cref{prop:ch48-no-strong-compactness} and
\Cref{ex:ch48-weak-not-strong,ex:ch48-p1-capacity,ex:ch48-proper-support-capacity},
and through the standing conventions of \Cref{chap:trivial-analytification}
for the description of \(X^{\An}\).

\emph{The probability normalization.}  Used everywhere, and its only
substantive consequences are \(E_\omega(\phi+c)=E_\omega(\phi)+c\), which
makes \(E^\vee_\omega\) finite exactly on probability measures
(\Cref{lem:ch47-mass-forced}) and makes the Fenchel defect
translation invariant
(\Cref{prop:ch48-defect-elementary}\,\cref{itm:ch48-def:transl}); and
\(\Capa_\omega(X^{\An})=1\).  Every comparison with
\Cref{chap:bfj-capacity,chap:bfj-finite-energy-measures} carries a factor
\(V\), displayed at each occurrence.

\subsection{What is not claimed}
\label{sec:ch48-not-claimed}

Eight statements that a reader might expect are \emph{not} asserted here,
and the reasons differ.

\emph{Strong compactness.}  Energy sublevels are weakly compact and not
strongly compact (\Cref{prop:ch48-no-strong-compactness}).  Nothing in this
chapter produces a strongly convergent subnet from a bounded-energy net;
the substitute is completeness plus the Cauchy property
(\Cref{prop:ch48-criteria}\,\cref{itm:ch48-crit:cauchy}).

\emph{Metrizability by \(\mathbb I\).}  \(\mathbb I\) is a quasi-metric and
not a metric (\Cref{ex:ch48-not-a-metric}).  Metrizability of the strong
topology is \Cref{thm:stability-norm-completion}, proved in
\Cref{chap:bj-stability} by different means, and no isometry between
\(d_1\) and \(\mathbb I\) is asserted anywhere.

\emph{Weak continuity of \(\MA_\omega\).}  Not asserted.  What is proved is
strong continuity in both directions
(\Cref{thm:ch48-ma-homeomorphism}), and continuity along monotone nets
(\Cref{prop:bj-strong-monotone}).
\Cref{ex:ch48-weak-not-strong} shows that weak convergence of potentials
and of their Monge--Amp\`ere measures can hold simultaneously with a
uniform positive strong distance.

\emph{Increasing-net continuity of the energy.}  Imported
(\Cref{imp:ch48-imported}\,\cref{itm:ch48-import:increasing}), not proved.
\Cref{warn:ch48-monotone-asymmetry} explains why the decreasing proof does
not reverse.

\emph{Strong subadditivity of the capacity, and capacitability.}  Not
proved; \cite[Theorem~13.29]{BJ22} is quoted where it is used, and
\Cref{ex:ch48-proper-support-capacity} shows that capacitability fails for
general subsets.

\emph{``Negligible \(\Rightarrow\) pluripolar'' with a rate.}  Not
available; see \Cref{sec:ch48-what-is-open}.

\emph{Any comparison with Thuillier's relative capacity.}  Not asserted in
either direction.  \Cref{sec:ch48-part3} explains why his hypotheses fail
over a trivially valued field, and \Cref{rem:ch25-thuillier-vs-bfj} records
the four differences between the two set functions.

\emph{Sequential forms of the compactness statements.}  Not available.
\(X^{\An}\) is compact but not metrizable, \(\PSH^{\sup}(\omega)\) is not
known to be metrizable, and every convergence statement above is about
nets; the only sequential statements are
\Cref{prop:ch48-capacity-subsequence} and Step~1 of the proof of
\Cref{thm:bj-quasicontinuity}, both of which extract a sequence by hand
from a countable amount of data.  See \Cref{rem:ch48-nets}.

\section{Chapter dictionary}
\label{sec:ch48-dictionary}

\begingroup
\begin{longtable}{p{.24\textwidth}p{.38\textwidth}p{.28\textwidth}}
\toprule
object & content & where\\
\midrule
\endfirsthead
\toprule
object & content & where\\
\midrule
\endhead
weak topology &
pointwise convergence on \(X^{\divis}\); Hausdorff &
\eqref{eq:bj-e1-weak},
\Cref{lem:ch48-net-criterion}\\
strong topology &
coarsest refinement making \(E_\omega\) continuous; equivalently weak
convergence plus \(E_\omega(\phi_i)\to E_\omega(\phi)\) &
\Cref{def:bj-strong-e1},
\Cref{lem:ch48-net-criterion}\\
potential gauge &
\(\mathbb I=I+|\sup\phi-\sup\psi|\) induces the strong topology;
quasi-metric, not a metric &
\eqref{eq:bj-strong-e1-I},
\Cref{thm:bj-strong-e1-quasimetric},
\Cref{ex:ch48-not-a-metric}\\
curve case &
for \(n=1\), \(I=\lVert\phi-\psi\rVert^2\) for a fixed inner product, and
\Cref{thm:bj-strong-e1-quasimetric} is proved outright &
\Cref{prop:ch48-curve-strong}\\
monotone criteria &
decreasing \(\Rightarrow\) strong; increasing \(\Rightarrow\) strong
(imported); uniform \(\Rightarrow\) strong &
\Cref{prop:bj-strong-monotone},
\Cref{prop:ch48-uniform-strong}\\
completeness &
every \(\mathbb I\)-Cauchy net converges; equals
\Cref{thm:bj-cy-envelope-equivalence}\,(4) &
\Cref{thm:ch48-completeness}\\
Fenchel defect &
\(J_\mu(\phi)=E^\vee(\mu)-E(\phi)+\int\phi\rd\mu\ge0\), \(=J_\psi(\phi)\)
when \(\mu=\MA(\psi)\), \(=0\) iff \(\MA(\phi)=\mu\) &
\eqref{eq:bj-fenchel-defect},
\Cref{prop:ch48-defect-elementary}\\
measure gauge &
\(I^\vee(\mu,\nu)=\inf_\phi(J_\mu(\phi)+J_\nu(\phi))\); symmetric,
separating &
\eqref{eq:bj-measure-Idual},
\Cref{prop:ch48-Idual-elementary}\\
strong measures &
weak \(+\)\,\(E^\vee\) convergence \(\Leftrightarrow\) \(I^\vee\to0\)
\(\Leftrightarrow\) uniform pairing on \(J\)-sublevels; complete &
\Cref{thm:bj-strong-measures}\\
bi-Lipschitz &
\(\frac1{(n+1)C_n}I\le I^\vee(\MA\phi,\MA\psi)\le\frac n{n+1}I\); equals
\(\frac14I\) when \(n=1\) &
\Cref{prop:bj-ma-strong-embedding},
\Cref{prop:ch48-curve-Idual}\\
main theorem &
\(\MA_\omega\) is a strongly continuous, strongly open surjection with
fibres \(\phi+\mathbf R\), inducing a bi-Lipschitz homeomorphism
\(\cE^1/\mathbf R\simeq\cM^1\) &
\Cref{thm:ch48-ma-homeomorphism}\\
energies compared &
\(\frac1nJ(\phi)\le E^\vee(\MA\phi)\le nJ(\phi)\), equality for \(n=1\) &
\Cref{thm:ch48-ma-homeomorphism}\,\cref{itm:ch48-ma:energy}\\
capacity &
\(\Capa_\omega(A)=\sup\{\MA_\omega(\psi)(A)\mid-1\le\psi\le0\}\), a
probability; \(\Capa_\omega(\{v_{\triv}\})=1\) &
\eqref{eq:bj-bt-capacity},
\Cref{prop:ch48-capacity-elementary}\\
CLN estimate &
\(\int(-\phi)\MA_\omega(\psi_1,\dots,\psi_n)\le n\) for
\(\sup\phi=0\), \(-1\le\psi_i\le0\) &
\Cref{lem:ch48-cln}\\
sublevels &
\(\Capa(\{\phi\le-t\})\le n/t\), attained;
\(J(\phi)\le1+\int_1^\infty t^n\Capa(\{\phi\le-t\})\rd t\) &
\Cref{prop:bj-capacity-sublevel},
\Cref{ex:ch48-p1-capacity}\\
fundamental estimate &
\(\int(\psi-\varphi)\MA(u_1,\dots,u_n)\le4M
(\int(\psi-\varphi)\MA(\varphi/2))^{2^{-n}}\) &
\Cref{lem:ch48-fundamental}\\
quasicontinuity &
psh functions are continuous off open sets of arbitrarily small capacity;
one set for a countable family &
\eqref{eq:bj-quasicontinuity},
\Cref{thm:bj-quasicontinuity}\\
capacity vs.\ measure &
\(\mu(B)\le t^n\Capa(B)+(n+1)E^\vee(\mu)/t\); finite-energy measures ignore
capacity-null sets &
\Cref{prop:ch48-measure-vs-capacity}\\
three smallness notions &
pluripolar \(\Leftrightarrow\) negligible \(\Leftrightarrow\)
\(\Capa^*=0\) &
\eqref{eq:bj-small-set-equivalence},
\Cref{thm:bj-small-set-equivalence}\\
Alexander--Taylor &
\(\min\{1,T^{-n}\}\le\Capa^*\le n/T\); both bounds proved for compact
sets, equality when \(n=1\) &
\eqref{eq:bj-alexander-taylor},
\Cref{prop:ch48-at-upper},
\Cref{prop:ch48-at-compact}\\
compactness &
weak sublevels compact, strong sublevels not &
\Cref{prop:bj-strong-compactness-criterion},
\Cref{prop:ch48-no-strong-compactness}\\
counterexample &
uncountably many proper-support points: inner capacity \(0\), outer
capacity \(1\), not pluripolar &
\Cref{ex:ch48-proper-support-capacity}\\
\bottomrule
\end{longtable}
\endgroup

Four separations should be carried forward.

\emph{Weak versus strong.}  Every compactness statement in
\Crefrange{chap:bj-energy}{chap:bj-strong-topologies} is weak; every
statement identifying a limit is strong.  \Cref{ex:ch48-weak-not-strong}
shows that the two can disagree on both sides of \(\MA_\omega\)
simultaneously, and that the discrepancy is measured by a single number.

\emph{Inner versus outer capacity.}  \(\Capa_\omega\) is defined on Borel
sets by a supremum of masses and \(\Capa^*_\omega\) on all sets by an
infimum over open neighbourhoods.  They agree on open sets, and
\(\Capa_\omega\le\Capa^*_\omega\); \Cref{ex:ch48-proper-support-capacity}
shows the gap can be the whole interval \([0,1]\).  Every statement about
smallness is about \(\Capa^*_\omega\).

\emph{Proved versus imported.}  Three statements are imported and are
listed once, in \Cref{imp:ch48-imported}: the uniform pairing estimate,
continuity of the energy along increasing nets, and the two arrows of
\Cref{thm:bj-small-set-equivalence} that go from smallness back to
pluripolarity.  Everything else is proved, and
\Cref{prop:ch48-curve-strong},
\Cref{prop:ch48-curve-Idual} and \Cref{prop:ch48-at-compact} remove two of
the three in dimension one and for compact sets respectively.

\emph{Capacity versus energy.}  These measure different things.
\Cref{ex:ch48-weak-not-strong} exhibits measures sitting on sets of full
capacity whose dual energies do not converge, and
\Cref{ex:ch48-p1-capacity} exhibits a potential of infinite energy whose
sublevel capacities decay at the maximal permitted rate.  The bridge
between them is \Cref{prop:bj-capacity-sublevel}, and it is one-directional
in each of its two halves.

\section{Source and provenance}
\label{sec:ch48-source}

\subsection{Source note}
\label{sec:ch48-source-note}

Strong measure geometry is \cite[\S10]{BJ22}, strong potential geometry is
\cite[\S12]{BJ22}, and capacity is \cite[\S\S7.7 and~13]{BJ22}.  The proof
that negligible sets are pluripolar is \cite[Theorem~13.17]{BJ22} and
depends on the envelope property.
The capacity theory is therefore not being used to reprove that property;
its logical direction is the reverse.  \Cref{chap:bfj-finite-energy-measures}
uses analogous words in a
different, nontrivially valued synthetic setting, with its own hypotheses
and normalization.

Statement by statement, the attributions are as follows.
\Cref{def:bj-strong-e1} and \Cref{thm:bj-strong-e1-quasimetric} are
\cite[\S12.1 and Theorem~12.4]{BJ22}; the continuity of the energy along
increasing nets used in \Cref{prop:bj-strong-monotone} and
\Cref{thm:ch48-completeness} is \cite[Theorem~12.5]{BJ22}; the completeness
statement is the implication (1)\(\Rightarrow\)(4) of
\cite[Theorem~12.8]{BJ22}.  On the measure side,
\Cref{def:bj-measure-quasimetric} and \Cref{thm:bj-strong-measures} are
\cite[Theorems~10.2--10.3, 10.12 and~10.14]{BJ22}, the uniform pairing
estimate being \cite[Theorem~10.3]{BJ22}; strong density of the image of
\(\MA_\omega\) is \cite[Corollary~10.13]{BJ22}; the independence of the
polarization is \cite[Corollary~9.21 and Proposition~10.16]{BJ22}.  The
capacity is \cite[\S7.7]{BJ22}, its Choquet properties
\cite[Propositions~13.3--13.4, Corollary~13.13 and Theorem~13.29]{BJ22},
the energy criterion \cite[Lemma~7.47]{BJ22}, quasicontinuity
\cite[Theorem~13.9 and Corollary~13.10]{BJ22}, the Alexander--Taylor
inequality \cite[Theorem~13.14]{BJ22}, the implication ``pluripolar
\(\Rightarrow\) negligible'' \cite[Proposition~5.3]{BJ22}, and the
equivalence \eqref{eq:bj-small-set-equivalence}
\cite[Theorems~13.14 and~13.17]{BJ22}.  Orthogonality for lower
semicontinuous obstacles, used in \Cref{prop:ch48-at-compact}, is
\cite[Theorem~13.1]{BJ22}.  The metric \(d_1\) of
\Cref{rem:ch48-d1-comparison} is \cite[Theorem~B]{BJKS1}.

The complex prototypes are named where they occur and are collected in
\Cref{sec:ch48-history}.  The chapter's only non-BJ inputs are the formal
capacity calculus of \Cref{sec:capacity-preview} and the net and measure
technology of \Cref{app:nets}.

\subsection{What is supplied here and is not in the source}
\label{sec:ch48-additions}

The following material is not in \cite{BJ22} in this form and is supplied
here.

\Cref{lem:ch48-net-criterion}: the identification of the strong topology by
its subbasis and its net criterion, and the observation that upper
semicontinuity of \(E_\omega\) makes the strengthening one-sided.  The
statement is folklore; it is written out because
\Cref{rem:ch48-usc-is-the-point} is used to justify the phrase ``the
coarsest refinement'' in \Cref{def:bj-strong-e1}, and because
the corresponding statement over a discretely valued base
(\Cref{prop:ch26-strong-basic}) is proved there and not here.

\Cref{sec:ch48-curve-case}: the complete proof of
\Cref{thm:bj-strong-e1-quasimetric} for \(n=1\), in both directions, from
the pre-Hilbert structure of \Cref{prop:ch47-curve-I-norm}.  The two
mechanisms it isolates --- a uniform bound on the pairing over a
\(J\)-sublevel, and strong density of finitely supported divisorial
measures --- are exactly the two that the imported estimate supplies in
general, and \Cref{rem:ch48-curve-isolates} says so.

\Cref{prop:ch48-defect-elementary} and
\Cref{prop:ch48-Idual-elementary}: the elementary structure of the Fenchel
defect and of \(I^\vee\), in particular the identity
\(J_{\MA_\omega(\psi)}=J_\psi\), the separation property deduced from
\Cref{thm:bj-maximizers-converge}, and the disclaimer
\Cref{rem:ch48-Idual-not-convex}.  These are implicit in
\cite[\S10]{BJ22}; they are written out because they are the exact
analogues of
\Cref{prop:ch28-fenchel-defect,prop:ch28-dual-order} and because the
comparison of the two regimes is one of the stated purposes of the
chapter.

\Cref{prop:ch48-curve-Idual}: the exact value
\(I^\vee(\MA\phi,\MA\psi)=\tfrac14I(\phi,\psi)\) when \(n=1\), with the
midpoint as minimizer.  This shows that the lower bi-Lipschitz constant
\(\frac1{(n+1)C_n}\) of \Cref{prop:bj-ma-strong-embedding} is attained, and
it supplies the numbers used in \Cref{ex:ch48-weak-not-strong}.

\Cref{thm:ch48-ma-homeomorphism}: the packaging of
\Cref{prop:bj-ma-strong-embedding} as a homeomorphism statement, together
with the splitting \(\cE^1\simeq\cE^1_0\times\mathbf R\), the openness of
\(\MA_\omega\), and the two-sided comparison
\(\frac1nJ(\phi)\le E^\vee_\omega(\MA_\omega\phi)\le nJ(\phi)\).  The
comparison is elementary but does not appear in this form in the source,
and it is what makes \Cref{prop:ch48-measure-vs-capacity} quantitative.

\Cref{lem:ch48-cln} and \Cref{lem:ch48-fundamental}: the
Chern--Levine--Nirenberg estimate with the constant \(n\), and the
fundamental capacity estimate with the constant \(4M\), both in the
probability normalization.  The first is where the identity
\(\MA_\omega(0)=\delta_{v_{\triv}}\) removes the additive constant that
\Cref{lem:ch25-normalized-bound} has to carry; the second is a literal
transfer of \Cref{lem:fundamental-capacity-estimate}, and
\Cref{rem:ch48-fundamental-transfer} records what the transfer costs and
why it does not serve increasing nets.

The complete proofs of \Cref{prop:bj-capacity-sublevel} and
\Cref{thm:bj-quasicontinuity}.  In the source these are stated with proofs
of a different shape: quasicontinuity there is proved by regularizing with
Fubini--Study functions in varying classes, whereas
\Cref{thm:ch45-continuous-regularization} makes a fixed-class continuous
regularization available, and the proof shortens to one application of
\Cref{lem:ch48-fundamental} and one Chebyshev inequality.  The
lower bound \(\Capa_\omega(\{\phi\le-t\})\ge t^{-n}\MA_\omega(\phi)
(\{\phi\le-t\})\) of Step~2, which is what makes
\eqref{eq:bj-cap-energy-criterion} an equality-driven criterion rather than
a citation, is written out here.

\Cref{prop:ch48-measure-vs-capacity}: that a finite-energy measure is
dominated by a power of the capacity, with the explicit exponent
\(1/(n+1)\).  This settles
\Cref{rem:ch47-polar-scope}\,\cref{itm:ch47-polar:pluripolar}, deferred from
\Cref{chap:bj-finite-energy}, and it does so using only the easy half of
\Cref{thm:bj-small-set-equivalence}.  Over a discretely valued base the
corresponding domination is a \emph{hypothesis}
(\Cref{def:ch29-capacity-domination}); here it is a theorem, because the
Calabi--Yau theorem makes every finite-energy measure a Monge--Amp\`ere
measure.

\Cref{prop:ch48-pluripolar-capacity},
\Cref{prop:ch48-pluripolar-negligible},
\Cref{lem:ch48-normalize-pole}, \Cref{prop:ch48-at-upper},
\Cref{prop:ch48-at-compact}: the elementary half of
\eqref{eq:bj-small-set-equivalence} and of
\eqref{eq:bj-alexander-taylor}, proved in full, together with the compact
case of the hard half.  The construction \(\sum_j2^{-j}\phi_j\) that turns
an infinite Alexander--Taylor height into a pole function uses that
\(v_{\triv}\) is a common maximum point, and has no counterpart over a
discretely valued base.  The proof of
\Cref{prop:ch48-at-compact} through the equilibrium potential of a compact
set and orthogonality is written out because it identifies exactly which
step fails for a non-compact set, which is the content of
\Cref{sec:ch48-what-is-open}.

\Cref{prop:ch48-discharge}: the explicit discharge of the two steps that
\Cref{thm:bj-obstacle-orthogonality-energy} deferred to this chapter, with
\Cref{rem:ch48-logical-direction} auditing the dependency.

\Cref{prop:ch48-criteria} and \Cref{prop:ch48-no-strong-compactness}: the
catalogue of convergence criteria and the failure of strong compactness,
the latter with explicit numbers.

\Cref{ex:ch48-not-a-metric},
\Cref{ex:ch48-weak-not-strong},
\Cref{ex:ch48-p1-capacity},
\Cref{ex:ch48-proper-support-capacity},
\Cref{ex:ch48-pluripolar-cap-zero}: all new.  They exist because the two
claims this chapter makes --- that the strong topologies are the right
ones, and that the three notions of smallness agree --- are exactly the
claims a plausible reading can appear to establish without doing so.  A
reader who has watched a sequence of the simplest PL psh functions on
\(\mathbf P^1\) converge weakly on both sides of \(\MA_\omega\) while
staying at distance \(1\) and \(\tfrac14\) respectively, computed a
capacity on a skeleton and found the Alexander--Taylor inequality to be an
equality, and seen a set whose compact subsets all have capacity \(0\) and
whose outer capacity is \(1\), will not mistake the weak topology for the
strong one or inner capacity for outer capacity.

\subsection{Historical note}
\label{sec:ch48-history}

Capacity entered pluripotential theory through Bedford and Taylor.  Their
\cite{BT76} solved the Dirichlet problem for the complex Monge--Amp\`ere
operator on bounded psh functions; \cite{BT82} introduced the capacity that
bears their name, proved quasicontinuity, and proved that negligible sets
are pluripolar; \cite{BT87} added the fine topology.  The estimate that
makes all of it work --- what is called here the Chern--Levine--Nirenberg
inequality, \Cref{lem:ch48-cln} --- goes back to \cite{CLN69}, where it was
introduced for a different purpose, the construction of intrinsic norms.
Josefson's theorem \cite{Jos78}, that a locally pluripolar subset of
\(\mathbf C^n\) is globally pluripolar, is what allows the complex theory to
convert a capacity statement into a pole function; there is no
non-Archimedean analogue, and its absence is the reason that one arrow of
\Cref{thm:bj-small-set-equivalence} is imported here.  Klimek's book
\cite{Kl91} is the standard systematic account of the local theory, and it
is the one that Thuillier follows on curves.

The compact Kähler theory is due to Guedj and Zeriahi.  \cite{GZ05}
constructs the intrinsic capacities on a compact Kähler manifold and proves
the comparison between the capacity of a set and its global extremal
function --- the complex prototype of
\eqref{eq:bj-alexander-taylor}; \cite{GZ07} introduces the weighted energy
classes, of which \(\cE^1\) is the one used here, and proves the criterion
that a potential has finite energy exactly when the capacities of its
sublevel sets decay fast enough, the prototype of
\eqref{eq:bj-cap-energy-criterion}.  The variational solution of the
Monge--Amp\`ere equation in \(\cE^1\) is
Berman--Boucksom--Guedj--Zeriahi \cite{BBGZ13}, building on Cegrell's
energy classes \cite{Ceg98} and on Ko\l{}odziej's capacity estimates
\cite{Kol98}; the textbook account is \cite{GZ17}.  The metric geometry of
\(\cE^1\), and in particular the fact that the natural gauge is a genuine
metric only after the right normalization, is Darvas's
\cite{Dar15,Dar17}; \Cref{ex:ch48-not-a-metric} is the non-Archimedean
shadow of that discussion.

The non-Archimedean story has three episodes.  Thuillier's thesis
\cite{Thu05} built, for Berkovich curves over a nontrivially valued field,
a complete potential theory with an energy-theoretic relative capacity, and
proved that its polar sets are exactly the sets of vanishing outer
capacity, \emph{without} using quasicontinuity; \Cref{sec:ch25-thuillier}
gives that theory and \Cref{sec:ch48-part3} explains why it does not
transport to the present setting.  Boucksom, Favre and Jonsson
\cite{BFJ15,BFJ16} built the model-theoretic theory over a discretely
valued base, where quasicontinuity became the central technical result;
\Cref{chap:bfj-capacity} is an account of it, and
\Cref{rem:ch25-curve-converse} records that the converse to ``pluripolar
\(\Rightarrow\) negligible'' is not available there.  Boucksom and Jonsson
\cite{BJ22} built the trivially valued theory, where the envelope property
supplies both the continuous regularization and the missing converse; their
\S\S10, 12 and 13 are the source of this chapter.

Two threads run out of it.  On the metric side,
\cite{BJKS1} identifies \(\cE^1\) and \(\cM^1\) with completions of spaces
of graded norms and filtrations for a genuine metric \(d_1\), and
\Cref{chap:bj-stability} develops that identification; the strong topology
constructed here is what makes the identification possible, and
\Cref{rem:ch48-d1-comparison} records exactly what is and is not claimed
about the comparison.  On the measure side, \cite{BJ25FiniteEnergy}
develops a synthetic theory of finite-energy measures in which the strong
topology is an axiom rather than a theorem; that theory is the subject of
\Cref{chap:bfj-finite-energy-measures} and should not be attributed
retroactively to \cite{BJ22}.  Finally, the regularity of the Green
functions whose existence is guaranteed by
\Cref{cor:ch47-divisorial-green} --- and which are used in
\Cref{prop:ch48-curve-strong} --- is genuinely subtle: by
\cite[Theorem~A]{BJGreen24} they need not be \(\mathbf Q\)-PL, so the
Cauchy--Schwarz argument of \Cref{sec:ch48-curve-case} cannot be replaced
by a finite computation on one test configuration.
